\documentclass{report}
\usepackage[utf8]{inputenc}
\usepackage[T1]{fontenc}
\usepackage[french]{babel}
\usepackage{amsmath, amssymb, amsthm}
\usepackage{graphicx}
\usepackage{xcolor}
\usepackage{tikz}
\usepackage{hyperref}
\usepackage{geometry}
\geometry{a4paper, margin=2cm}

\title{Applications lin\'eaires, noyau ($\ker$), image ($\text{Im}$) et d\'emonstration du th\'eor\`eme du rang}
\author{Charles EDOU NZE}
\date{\today}

\begin{document}
\maketitle
\tableofcontents
\newpage

\chapter*{Jalon 8 : Applications linéaires, noyau ($\ker$), image ($\text{Im}$) et démonstration du théorème du rang}

\section*{1. Présentation du concept clé}

\begin{figure}[h]
\centering
\begin{tikzpicture}
  \draw[thick, fill=blue!10] (0,0) ellipse (2cm and 3cm);
  \node at (0, 3.5) {Espace vectoriel $E$};
  \draw[thick, fill=blue!30] (0,-1) ellipse (1cm and 1.5cm);
  \node at (0, -0.5) {$\ker(f)$};
  \filldraw[black] (0,-1.5) circle (1.5pt) node[anchor=north] {$0_E$};
  \draw[thick, fill=red!10] (6,0) ellipse (2cm and 3cm);
  \node at (6, 3.5) {Espace vectoriel $F$};
  \draw[thick, fill=red!30] (6,1) ellipse (1cm and 1.5cm);
  \node at (6, 1.5) {$\text{Im}(f)$};
  \filldraw[black] (6,0.5) circle (1.5pt) node[anchor=north] {$0_F$};
  \draw[->, thick] (1,-1) to[bend right=20] node[midway, below] {$f$} (5,0.5);
  \draw[->, thick] (1,1.5) to[bend left=20] node[midway, above] {$f$} (5,1.5);
\end{tikzpicture}
\end{figure}


 Imaginez une machine à étirer ou à faire tourner de la pâte à modeler.
  - L'\textbf{Application Linéaire}, c'est une machine respectueuse : si vous doublez la quantité de pâte au départ, vous aurez le double à l'arrivée. Si vous mettez deux morceaux de couleurs différentes, le résultat sera le même que si vous les aviez traités séparément puis mélangés.
  - Le \textbf{Noyau ($\ker$)}, c'est la "poubelle" de la machine : ce sont tous les morceaux de pâte qui finissent écrasés en un tout petit point (zéro) après le passage dans la machine.
  - L'\textbf{Image ($\text{Im}$)}, c'est l'ensemble de toutes les formes que la machine est capable de fabriquer.
  - Le \textbf{Théorème du Rang}, c'est la loi de conservation : la complexité totale de votre pâte au départ est égale à la somme de ce qui a été écrasé (Noyau) et de ce qui a été transformé (Image).
 La plupart des phénomènes physiques simples sont linéaires. En IA, transformer une image ou un texte revient souvent à appliquer une suite de ces machines. Savoir ce qui est "perdu" (le noyau) est vital.
 Imaginez projeter l'ombre d'un cube 3D sur une feuille 2D. L'ombre est l'\textbf{Image} (2D). La direction de la lumière qui "écrase" la profondeur est le \textbf{Noyau} (1D). $3 = 1 + 2$. C'est le théorème du rang !

\section*{2. Formalisation}

\subsection*{A. Définitions Formelles}
Soient $E$ et $F$ deux espaces vectoriels sur un corps commutatif $\mathbb{K}$.
1.  \textbf{Application Linéaire ($f : E \to F$) :} Une application $f: E \to F$ est dite linéaire si et seulement si elle vérifie les deux propriétés suivantes :
    -   $\forall x, y \in E, f(x + y) = f(x) + f(y)$ (Additivité).
    -   $\forall \lambda \in \mathbb{K}, \forall x \in E, f(\lambda \cdot x) = \lambda \cdot f(x)$ (Homogénéité).
    L'ensemble des applications linéaires de $E$ dans $F$ est noté $\mathcal{L}(E, F)$.

2.  \textbf{Noyau ($\ker f$) :} Le noyau de $f$ est l'ensemble des vecteurs de $E$ dont l'image par $f$ est le vecteur nul de $F$.
    $\ker f = \{ x \in E \mid f(x) = 0_F \}$.
    Le noyau $\ker f$ est un sous-espace vectoriel de $E$.

3.  \textbf{Image ($\text{Im } f$) :} L'image de $f$ est l'ensemble de toutes les images des vecteurs de $E$ par $f$.
    $\text{Im } f = f(E) = \{ f(x) \mid x \in E \}$.
    L'image $\text{Im } f$ est un sous-espace vectoriel de $F$.

4.  \textbf{Rang ($\text{rg } f$) :} On appelle rang de $f$ la dimension de son image.
    $\text{rg } f = \dim(\text{Im } f)$.

\subsection*{B. Théorèmes, Propositions \& Lemmes}
> \textbf{Caractérisation de l'injectivité :}
> Une application linéaire $f: E \to F$ est injective si et seulement si son noyau est réduit au vecteur nul de $E$.
> $f \text{ est injective } \iff \ker f = \{ 0_E \}$.

> \textbf{Théorème du Rang (Théorème Fondamental) :}
> Soit $E$ un espace vectoriel de dimension finie sur un corps $\mathbb{K}$. Pour toute application linéaire $f : E \to F$ (où $F$ est un $\mathbb{K}$-espace vectoriel), on a la relation fondamentale :
> $$\dim E = \dim(\ker f) + \text{rg}(f)$$

\section*{3. Démonstrations}

\subsection*{Démonstration du Théorème Pivot : Théorème du Rang}
Soient $E$ et $F$ deux $\mathbb{K}$-espaces vectoriels, avec $E$ de dimension finie. Soit $f \in \mathcal{L}(E, F)$ une application linéaire de $E$ vers $F$. Supposons que $\dim E = n$, où $n \in \mathbb{N}$ est un entier naturel.

1.  \textbf{Initialisation / Cadre :}
    -   Soit $p$ la dimension du noyau de $f$, notée $p = \dim(\ker f)$. Par définition, $p \in \mathbb{N}$ et $0 \le p \le n$.
    -   Si $p > 0$, soit $\mathcal{B}_{\ker f} = (e_1, ..., e_p)$ une base du sous-espace vectoriel $\ker f$. Si $p=0$, alors $\ker f = \{0_E\}$ et la base est vide.
    -   D'après le théorème de la base incomplète, puisque $\mathcal{B}_{\ker f}$ est une famille libre dans $E$, nous pouvons la compléter en une base de l'espace vectoriel $E$. Soit cette base $\mathcal{B}_E = (e_1, ..., e_p, e_{p+1}, ..., e_n)$, où les vecteurs $e_{p+1}, ..., e_n$ sont choisis de manière appropriée.
    -   Notre objectif est de démontrer que la famille de vecteurs $\mathcal{B}_{\text{Im } f} = (f(e_{p+1}), ..., f(e_n))$ constitue une base de l'image de $f$, $\text{Im } f$. Si cette démonstration est concluante, alors la dimension de l'image sera $\dim(\text{Im } f) = n - p$. En substituant cette relation dans l'équation du théorème du rang, nous obtiendrons $\dim E = \dim(\ker f) + \dim(\text{Im } f)$, c'est-à-dire $n = p + (n-p)$, ce qui validera le théorème.

2.  \textbf{Étape 1 : Montrons que la famille $(f(e_{p+1}), ..., f(e_n))$ est génératrice de $\text{Im } f$}
    -   Soit $y$ un vecteur arbitraire appartenant à l'image de $f$, c'est-à-dire $y \in \text{Im } f \subseteq F$.
    -   Par définition de l'image, il existe un vecteur $x \in E$ tel que $y = f(x)$.
    -   Puisque $\mathcal{B}_E = (e_1, ..., e_n)$ est une base de l'espace vectoriel $E$, tout vecteur $x \in E$ peut être exprimé comme une combinaison linéaire unique des vecteurs de cette base. Ainsi, il existe des scalaires uniques $\lambda_1, ..., \lambda_n \in \mathbb{K}$ tels que :
        $x = \sum_{i=1}^n \lambda_i e_i$.
    -   Appliquons l'application linéaire $f$ au vecteur $x$ :
        $y = f\left(\sum_{i=1}^n \lambda_i e_i\right)$.
    -   Par la propriété de linéarité de $f$ (additivité et homogénéité), nous pouvons écrire :
        $y = \sum_{i=1}^n \lambda_i f(e_i)$.
    -   Nous pouvons séparer cette somme en deux parties :
        $y = \sum_{i=1}^p \lambda_i f(e_i) + \sum_{i=p+1}^n \lambda_i f(e_i)$.
    -   Par définition du noyau, pour tout $i \in \{1, ..., p\}$, le vecteur $e_i$ appartient à $\ker f$. Par conséquent, l'image de $e_i$ par $f$ est le vecteur nul de $F$, c'est-à-dire $f(e_i) = 0_F$.
    -   Substituons cette information dans la somme :
        $y = \sum_{i=1}^p \lambda_i \cdot 0_F + \sum_{i=p+1}^n \lambda_i f(e_i)$.
    -   Puisque $\lambda_i \cdot 0_F = 0_F$ pour tout $\lambda_i \in \mathbb{K}$, la première somme est nulle :
        $y = 0_F + \sum_{i=p+1}^n \lambda_i f(e_i)$.
    -   Ainsi, nous obtenons :
        $y = \sum_{i=p+1}^n \lambda_i f(e_i)$.
    -   Puisque tout vecteur $y \in \text{Im } f$ peut être exprimé comme une combinaison linéaire des vecteurs $(f(e_{p+1}), ..., f(e_n))$, la famille $\mathcal{B}_{\text{Im } f} = (f(e_{p+1}), ..., f(e_n))$ est une famille génératrice de $\text{Im } f$.

3.  \textbf{Étape 2 : Montrons que la famille $(f(e_{p+1}), ..., f(e_n))$ est libre}
    -   Considérons une combinaison linéaire nulle des vecteurs de la famille $\mathcal{B}_{\text{Im } f}$. Soient $\mu_{p+1}, ..., \mu_n \in \mathbb{K}$ des scalaires tels que :
        $\sum_{i=p+1}^n \mu_i f(e_i) = 0_F$.
    -   En utilisant la linéarité de l'application $f$, nous pouvons réécrire cette somme comme :
        $f\left(\sum_{i=p+1}^n \mu_i e_i\right) = 0_F$.
    -   Par définition du noyau, cela signifie que le vecteur $v = \sum_{i=p+1}^n \mu_i e_i$ appartient à $\ker f$.
    -   Puisque $\mathcal{B}_{\ker f} = (e_1, ..., e_p)$ est une base de $\ker f$, tout vecteur de $\ker f$ peut être exprimé comme une combinaison linéaire des éléments de $\mathcal{B}_{\ker f}$. Par conséquent, il existe des scalaires $\alpha_1, ..., \alpha_p \in \mathbb{K}$ tels que :
        $v = \sum_{j=1}^p \alpha_j e_j$.
    -   Ainsi, nous avons l'égalité :
        $\sum_{i=p+1}^n \mu_i e_i = \sum_{j=1}^p \alpha_j e_j$.
    -   Regroupons tous les termes du même côté de l'équation pour obtenir le vecteur nul $0_E$ :
        $\sum_{i=p+1}^n \mu_i e_i - \sum_{j=1}^p \alpha_j e_j = 0_E$.
    -   Réarrangeons les termes pour correspondre à l'ordre de la base $\mathcal{B}_E$ :
        $\sum_{j=1}^p (-\alpha_j) e_j + \sum_{i=p+1}^n \mu_i e_i = 0_E$.
    -   La famille $\mathcal{B}_E = (e_1, ..., e_p, e_{p+1}, ..., e_n)$ est une base de $E$, ce qui implique qu'elle est une famille libre. Par conséquent, la seule combinaison linéaire de ses vecteurs qui est égale au vecteur nul $0_E$ est celle où tous les coefficients sont nuls.
    -   De l'équation $\sum_{j=1}^p (-\alpha_j) e_j + \sum_{i=p+1}^n \mu_i e_i = 0_E$, il découle que tous les scalaires associés doivent être nuls :
        -   Pour $j \in \{1, ..., p\}$, $-\alpha_j = 0$, ce qui implique $\alpha_j = 0$.
        -   Pour $i \in \{p+1, ..., n\}$, $\mu_i = 0$.
    -   En particulier, tous les scalaires $\mu_i$ pour $i \in \{p+1, ..., n\}$ sont nuls.
    -   Puisque tous les scalaires $\mu_i$ sont nuls, la famille $\mathcal{B}_{\text{Im } f} = (f(e_{p+1}), ..., f(e_n))$ est une famille libre.

4.  \textbf{Conclusion :}
    -   Ayant démontré que la famille $(f(e_{p+1}), ..., f(e_n))$ est à la fois génératrice de $\text{Im } f$ et libre, nous pouvons affirmer qu'elle constitue une base de $\text{Im } f$. Le nombre de vecteurs dans cette base est $n - p$.
    -   Par définition du rang, $\text{rg}(f) = \dim(\text{Im } f)$. Donc, nous avons $\text{rg}(f) = n - p$.
    -   En substituant les définitions de $n$ et $p$, nous obtenons : $\text{rg}(f) = \dim E - \dim(\ker f)$.
    -   En réarrangeant les termes, nous obtenons l'énoncé du théorème du rang : $\dim E = \dim(\ker f) + \text{rg}(f)$.

\section*{4. Exercices d'Application}

\subsection*{Exercice 1 : Application Directe (Noyau et Image)}
 Soit $f : \mathbb{R}^3 \to \mathbb{R}^2$ une application définie pour tout $(x,y,z) \in \mathbb{R}^3$ par $f(x,y,z) = (x+y, y+z)$. Déterminer le noyau de $f$, $\ker f$, et le rang de $f$, $\text{rg } f$.

1.  \textbf{Détermination du Noyau ($\ker f$) :}
    -   Pour déterminer le noyau de $f$, nous cherchons l'ensemble des vecteurs $(x,y,z) \in \mathbb{R}^3$ tels que $f(x,y,z) = 0_{\mathbb{R}^2}$, où $0_{\mathbb{R}^2} = (0,0)$ est le vecteur nul de l'espace d'arrivée $\mathbb{R}^2$.
    -   Cela conduit au système d'équations linéaires suivant :
        1.  $x+y = 0$
        2.  $y+z = 0$
    -   De la première équation (1), nous déduisons une expression pour $x$ en fonction de $y$ :
        $x = -y$.
    -   De la seconde équation (2), nous déduisons une expression pour $z$ en fonction de $y$ :
        $z = -y$.
    -   Ainsi, tout vecteur $(x,y,z)$ appartenant à $\ker f$ doit être de la forme $(-y, y, -y)$ pour un certain scalaire $y \in \mathbb{R}$.
    -   Nous pouvons factoriser le scalaire $y$ de ce vecteur :
        $(-y, y, -y) = y \cdot (-1, 1, -1)$.
    -   Par conséquent, le noyau de $f$ est l'ensemble des multiples scalaires du vecteur $(-1, 1, -1)$ :
        $\ker f = \{ y \cdot (-1, 1, -1) \mid y \in \mathbb{R} \}$.
    -   Ceci est l'espace vectoriel engendré par le vecteur $(-1, 1, -1)$, noté $\text{Vect}((-1, 1, -1))$.
    -   Le vecteur $(-1, 1, -1)$ est non nul. Par conséquent, il forme une base de $\ker f$.
    -   La dimension du noyau est le nombre de vecteurs dans cette base :
        $\dim(\ker f) = 1$.

2.  \textbf{Détermination du Rang ($\text{rg } f$) en utilisant le Théorème du Rang :}
    -   L'espace de départ $E = \mathbb{R}^3$ est de dimension finie, et sa dimension est $\dim \mathbb{R}^3 = 3$.
    -   Nous pouvons appliquer le théorème du rang pour l'application linéaire $f : \mathbb{R}^3 \to \mathbb{R}^2$ :
        $\dim E = \dim(\ker f) + \text{rg}(f)$.
    -   En substituant les valeurs connues de $\dim E$ et $\dim(\ker f)$ :
        $3 = 1 + \text{rg}(f)$.
    -   Par conséquent, nous pouvons calculer le rang de $f$ :
        $\text{rg}(f) = 3 - 1 = 2$.

\textbf{Conclusion :}
Le noyau $\ker f$ est une droite vectorielle dans $\mathbb{R}^3$, engendrée par le vecteur $(-1, 1, -1)$. Sa dimension est $\dim(\ker f) = 1$.
Le rang de $f$ est $\text{rg}(f) = 2$. Puisque l'espace d'arrivée $F = \mathbb{R}^2$ a une dimension de 2, et que $\text{Im } f$ est un sous-espace vectoriel de $\mathbb{R}^2$ de dimension 2, il s'ensuit que $\text{Im } f = \mathbb{R}^2$. L'application $f$ est donc surjective.

\subsection*{Exercice 2 : Niveau Avancé (Projecteurs)}
 Soit $E$ un $\mathbb{K}$-espace vectoriel. Soit $p \in \mathcal{L}(E, E)$ un endomorphisme de $E$ (appelé projecteur) tel que $p \circ p = p$, c'est-à-dire $p(p(x)) = p(x)$ pour tout $x \in E$. Démontrer que $E$ est la somme directe du noyau de $p$ et de l'image de $p$, notée $E = \ker p \oplus \text{Im } p$.

Pour démontrer que $E = \ker p \oplus \text{Im } p$, nous devons établir deux conditions :
1.  \textbf{Somme :} Pour tout vecteur $x \in E$, il existe au moins un couple de vecteurs $(y, z)$ tel que $y \in \ker p$, $z \in \text{Im } p$ et $x = y + z$. (Ceci montre que $E = \ker p + \text{Im } p$).
2.  \textbf{Intersection réduite au vecteur nul :} L'intersection du noyau et de l'image est réduite au vecteur nul de $E$, c'est-à-dire $\ker p \cap \text{Im } p = \{0_E\}$. (Ceci montre que la somme est directe).

1.  \textbf{Démonstration de la Somme ($E = \ker p + \text{Im } p$) :}
    -   Soit $x$ un vecteur arbitraire de $E$. Nous cherchons à décomposer $x$ en une somme d'un élément du noyau et d'un élément de l'image.
    -   Considérons la décomposition suivante :
        $x = (x - p(x)) + p(x)$.
    -   Posons $z = p(x)$.
    -   Par définition de l'image, puisque $x \in E$, le vecteur $z = p(x)$ appartient à l'image de $p$, donc $z \in \text{Im } p$.
    -   Posons $y = x - p(x)$.
    -   Pour montrer que $y \in \ker p$, nous devons vérifier que $p(y) = 0_E$. Calculons $p(y)$ :
        $p(y) = p(x - p(x))$.
    -   Par la propriété de linéarité de $p$, nous pouvons distribuer $p$ sur la soustraction :
        $p(x - p(x)) = p(x) - p(p(x))$.
    -   En utilisant la propriété du projecteur $p \circ p = p$, nous savons que $p(p(x)) = p(x)$.
    -   Substituons cette propriété dans l'expression de $p(y)$ :
        $p(y) = p(x) - p(x)$.
    -   Ce qui simplifie à :
        $p(y) = 0_E$.
    -   Par conséquent, $y \in \ker p$.
    -   Nous avons ainsi montré que pour tout $x \in E$, $x$ peut s'écrire comme la somme d'un vecteur $y \in \ker p$ et d'un vecteur $z \in \text{Im } p$. Ceci établit que $E = \ker p + \text{Im } p$.

2.  \textbf{Démonstration de l'Intersection réduite au vecteur nul ($\ker p \cap \text{Im } p = \{0_E\}$) :}
    -   Pour démontrer que la somme est directe, nous devons montrer que l'intersection $\ker p \cap \text{Im } p$ est réduite au seul vecteur nul $0_E$.
    -   Soit $u$ un vecteur arbitraire appartenant à cette intersection, c'est-à-dire $u \in \ker p \cap \text{Im } p$.
    -   Puisque $u \in \ker p$, par définition du noyau, nous avons :
        $p(u) = 0_E$.
    -   Puisque $u \in \text{Im } p$, par définition de l'image, il existe un vecteur $v \in E$ tel que :
        $u = p(v)$.
    -   Substituons cette expression de $u$ dans l'équation $p(u) = 0_E$ :
        $p(p(v)) = 0_E$.
    -   En utilisant la propriété du projecteur $p \circ p = p$, nous avons $p(p(v)) = p(v)$.
    -   Donc, l'équation devient :
        $p(v) = 0_E$.
    -   Puisque nous avons établi que $u = p(v)$, nous en déduisons que :
        $u = 0_E$.
    -   L'intersection $\ker p \cap \text{Im } p$ est donc réduite au seul vecteur nul $\{0_E\}$.

\textbf{Conclusion :}
Ayant démontré que $E = \ker p + \text{Im } p$ (la somme) et que $\ker p \cap \text{Im } p = \{0_E\}$ (l'intersection est réduite au vecteur nul), nous pouvons conclure que l'espace vectoriel $E$ est la somme directe du noyau de $p$ et de l'image de $p$, c'est-à-dire $E = \ker p \oplus \text{Im } p$.

\section*{5. Application en Intelligence Artificielle}
-   \textbf{Le Pont Théorique :} Les couches denses (Dense Layers / Linear Layers) des réseaux de neurones sont des applications affines (Linéaire + Translation). Les concepts d'applications linéaires, de noyau et d'image sont donc directement applicables à la compréhension de ces couches.
-   \textbf{Exemple Concret :} Dans la \textbf{réduction de dimension} et la \textbf{Compression de Réseaux}, on cherche à savoir si une couche linéaire a un gros \textbf{Noyau}. Si $\dim(\ker f)$ est grand, cela signifie que beaucoup de combinaisons d'entrées sont "oubliées" ou projetées sur le même point (le vecteur nul) par le réseau, indiquant une redondance ou une faible expressivité dans certaines directions. Le \textbf{Rang} de la matrice de poids d'une couche détermine la "capacité expressive" de cette couche, c'est-à-dire la dimension de l'espace des sorties qu'elle peut effectivement générer. Un réseau "Low-Rank" (de faible rang) est plus rapide à entraîner et moins sujet au sur-apprentissage (Overfitting) car il impose une contrainte sur la complexité des transformations qu'il peut effectuer, favorisant ainsi la généralisation.

\section*{6. Liens Sémantiques}
-   \textbf{Concepts Précédents requis :} Jalon 7 (Espaces vectoriels abstraits)
-   \textbf{Concepts Futurs dépendants :} Jalon-9, Jalon 30 (Trigonalisation d'endomorphismes et décomposition de Dunford.), Jalon 46 (Matrice jacobienne)

\newpage
\chapter*{Exercices d\'Application}
\addcontentsline{toc}{chapter}{Exercices d'Application}
\chapter*{Exercice 1 : Applications linéaires élémentaires (Difficulté : \*)}

\section*{Énoncé du problème}

On considère l'application $f: \mathbb{R}^2 \to \mathbb{R}^3$ définie pour tout vecteur $(x, y) \in \mathbb{R}^2$ par :
$$f(x, y) = (x+y, x-y, 2x)$$
où $\mathbb{R}$ est le corps des nombres réels. L'ensemble $\mathbb{R}^2$ est l'espace vectoriel des couples de nombres réels, et $\mathbb{R}^3$ est l'espace vectoriel des triplets de nombres réels, tous deux munis des opérations d'addition vectorielle et de multiplication par un scalaire usuelles sur le corps $\mathbb{R}$.

1.  Démontrer que $f$ est une application linéaire.
2.  Déterminer le noyau de $f$, noté $\text{ker}(f)$. En donner une base et préciser sa dimension.
3.  Déterminer l'image de $f$, notée $\text{Im}(f)$. En donner une base et préciser sa dimension (qui est le rang de $f$).
4.  Vérifier le théorème du rang pour l'application $f$.

\section*{Correction détaillée}

\subsection*{1. Démontrer que $f$ est une application linéaire.}

Pour qu'une application $f: E \to F$ soit linéaire, où $E$ et $F$ sont des espaces vectoriels sur un corps $K$, elle doit satisfaire deux propriétés fondamentales :
a) Additivité : $f(\mathbf{u} + \mathbf{v}) = f(\mathbf{u}) + f(\mathbf{v})$ pour tous vecteurs $\mathbf{u}, \mathbf{v} \in E$.
b) Homogénéité : $f(\lambda \mathbf{u}) = \lambda f(\mathbf{u})$ pour tout scalaire $\lambda \in K$ et tout vecteur $\mathbf{u} \in E$.

Dans le cadre de cet exercice, l'espace de départ est $E = \mathbb{R}^2$, l'espace d'arrivée est $F = \mathbb{R}^3$, et le corps des scalaires est $K = \mathbb{R}$.

Soient $\mathbf{u} = (x_1, y_1)$ et $\mathbf{v} = (x_2, y_2)$ deux vecteurs arbitraires de l'espace vectoriel $\mathbb{R}^2$.
Soit $\lambda \in \mathbb{R}$ un scalaire arbitraire.

\subsubsection*{Propriété a) : Additivité}

Calculons d'abord la somme des vecteurs $\mathbf{u}$ et $\mathbf{v}$ dans $\mathbb{R}^2$:
$$\mathbf{u} + \mathbf{v} = (x_1, y_1) + (x_2, y_2) = (x_1+x_2, y_1+y_2)$$
Appliquons ensuite la fonction $f$ à ce vecteur somme :
$$f(\mathbf{u} + \mathbf{v}) = f(x_1+x_2, y_1+y_2)$$
Par la définition de $f$, où la première composante est la somme des coordonnées, la deuxième est leur différence, et la troisième est le double de la première coordonnée :
$$f(\mathbf{u} + \mathbf{v}) = ((x_1+x_2)+(y_1+y_2), (x_1+x_2)-(y_1+y_2), 2(x_1+x_2))$$
Réorganisons les termes au sein de chaque composante pour regrouper les termes relatifs à $\mathbf{u}$ et $\mathbf{v}$ :
$$f(\mathbf{u} + \mathbf{v}) = ((x_1+y_1)+(x_2+y_2), (x_1-y_1)+(x_2-y_2), 2x_1+2x_2)$$

Calculons maintenant la somme des images $f(\mathbf{u})$ et $f(\mathbf{v})$ :
$$f(\mathbf{u}) = (x_1+y_1, x_1-y_1, 2x_1)$$
$$f(\mathbf{v}) = (x_2+y_2, x_2-y_2, 2x_2)$$
Effectuons l'addition de ces deux vecteurs dans $\mathbb{R}^3$ :
$$f(\mathbf{u}) + f(\mathbf{v}) = ((x_1+y_1)+(x_2+y_2), (x_1-y_1)+(x_2-y_2), (2x_1)+(2x_2))$$
En comparant les expressions obtenues pour $f(\mathbf{u} + \mathbf{v})$ et $f(\mathbf{u}) + f(\mathbf{v})$, nous observons qu'elles sont identiques.
Ainsi, la propriété d'additivité $f(\mathbf{u} + \mathbf{v}) = f(\mathbf{u}) + f(\mathbf{v})$ est vérifiée.

\subsubsection*{Propriété b) : Homogénéité}

Calculons d'abord le produit du scalaire $\lambda$ par le vecteur $\mathbf{u}$ dans $\mathbb{R}^2$:
$$\lambda \mathbf{u} = \lambda (x_1, y_1) = (\lambda x_1, \lambda y_1)$$
Appliquons ensuite la fonction $f$ à ce vecteur :
$$f(\lambda \mathbf{u}) = f(\lambda x_1, \lambda y_1)$$
Par la définition de $f$ :
$$f(\lambda \mathbf{u}) = (\lambda x_1 + \lambda y_1, \lambda x_1 - \lambda y_1, 2\lambda x_1)$$
Factorisons le scalaire $\lambda$ dans chaque composante du vecteur résultant :
$$f(\lambda \mathbf{u}) = (\lambda(x_1 + y_1), \lambda(x_1 - y_1), \lambda(2x_1))$$
Nous pouvons extraire le scalaire $\lambda$ du vecteur :
$$f(\lambda \mathbf{u}) = \lambda(x_1 + y_1, x_1 - y_1, 2x_1)$$
Par la définition de $f(\mathbf{u})$, nous reconnaissons l'expression entre parenthèses :
$$f(\lambda \mathbf{u}) = \lambda f(\mathbf{u})$$
Ainsi, la propriété d'homogénéité $f(\lambda \mathbf{u}) = \lambda f(\mathbf{u})$ est vérifiée.

Puisque les deux propriétés d'additivité et d'homogénéité sont satisfaites, l'application $f: \mathbb{R}^2 \to \mathbb{R}^3$ est bien une application linéaire.

\subsection*{2. Déterminer le noyau de $f$, $\text{ker}(f)$. En donner une base et sa dimension.}

Le noyau de $f$, noté $\text{ker}(f)$, est défini comme l'ensemble des vecteurs $\mathbf{u} \in \mathbb{R}^2$ tels que leur image par $f$ est le vecteur nul de l'espace d'arrivée $\mathbb{R}^3$.
Formellement : $\text{ker}(f) = \{ \mathbf{u} \in \mathbb{R}^2 \mid f(\mathbf{u}) = \mathbf{0}_{\mathbb{R}^3} \}$.
Soit un vecteur $\mathbf{u} = (x, y) \in \mathbb{R}^2$. Le vecteur nul de $\mathbb{R}^3$ est $\mathbf{0}_{\mathbb{R}^3} = (0, 0, 0)$.
L'équation $f(x, y) = (0, 0, 0)$ se traduit par le système d'équations linéaires suivant :
1.  $x+y = 0$
2.  $x-y = 0$
3.  $2x = 0$

Nous allons résoudre ce système étape par étape :
À partir de l'équation (3) :
$$2x = 0$$
En divisant par 2 (qui est non nul dans $\mathbb{R}$), nous obtenons :
$$x = 0$$
Substituons la valeur de $x=0$ dans l'équation (1) :
$$0+y = 0$$
$$y = 0$$
Vérifions ces valeurs $x=0$ et $y=0$ dans l'équation (2) :
$$0-0 = 0$$
L'équation (2) est satisfaite.

Par conséquent, le seul vecteur $(x, y)$ de $\mathbb{R}^2$ qui satisfait la condition $f(x, y) = \mathbf{0}_{\mathbb{R}^3}$ est le vecteur nul $(0, 0)$.
Donc, le noyau de $f$ est $\text{ker}(f) = \{ (0, 0) \}$.

L'espace vectoriel $\{(0,0)\}$ est l'espace vectoriel réduit au vecteur nul. Une base de cet espace est l'ensemble vide $\emptyset$.
La dimension de $\text{ker}(f)$, notée $\text{dim}(\text{ker}(f))$, est le nombre de vecteurs dans sa base. Puisque la base est l'ensemble vide, la dimension est 0.
$$\text{dim}(\text{ker}(f)) = 0$$

\subsection*{3. Déterminer l'image de $f$, $\text{Im}(f)$. En donner une base et sa dimension.}

L'image de $f$, notée $\text{Im}(f)$, est l'ensemble des vecteurs de l'espace d'arrivée $\mathbb{R}^3$ qui sont des images d'au moins un vecteur de l'espace de départ $\mathbb{R}^2$ par l'application $f$.
Formellement : $\text{Im}(f) = \{ f(x,y) \mid (x,y) \in \mathbb{R}^2 \}$.
Considérons un vecteur générique $f(x,y)$ dans $\text{Im}(f)$ :
$$f(x,y) = (x+y, x-y, 2x)$$
Nous pouvons décomposer ce vecteur en une somme de vecteurs, en regroupant les termes dépendant de $x$ et ceux dépendant de $y$ :
$$f(x,y) = (x, x, 2x) + (y, -y, 0)$$
En factorisant les scalaires $x$ et $y$ de chaque vecteur :
$$f(x,y) = x(1,1,2) + y(1,-1,0)$$
Cette expression montre que tout vecteur dans l'image de $f$ peut être écrit comme une combinaison linéaire des vecteurs $v_1 = (1,1,2)$ et $v_2 = (1,-1,0)$.
Par conséquent, l'image de $f$ est l'espace vectoriel engendré par ces deux vecteurs :
$$\text{Im}(f) = \text{Vect}((1,1,2), (1,-1,0))$$

Pour déterminer une base de $\text{Im}(f)$, nous devons vérifier si les vecteurs $v_1$ et $v_2$ sont linéairement indépendants.
Soient $\lambda_1, \lambda_2 \in \mathbb{R}$ deux scalaires tels que leur combinaison linéaire est égale au vecteur nul de $\mathbb{R}^3$:
$$\lambda_1 v_1 + \lambda_2 v_2 = \mathbf{0}_{\mathbb{R}^3}$$
$$\lambda_1(1,1,2) + \lambda_2(1,-1,0) = (0,0,0)$$
Ceci conduit au système d'équations linéaires suivant, en égalant les composantes :
1.  $\lambda_1 + \lambda_2 = 0$
2.  $\lambda_1 - \lambda_2 = 0$
3.  $2\lambda_1 = 0$

Nous allons résoudre ce système étape par étape :
À partir de l'équation (3) :
$$2\lambda_1 = 0$$
En divisant par 2, nous obtenons :
$$\lambda_1 = 0$$
Substituons la valeur de $\lambda_1=0$ dans l'équation (1) :
$$0 + \lambda_2 = 0$$
$$\lambda_2 = 0$$
Vérifions ces valeurs $\lambda_1=0$ et $\lambda_2=0$ dans l'équation (2) :
$$0 - 0 = 0$$
L'équation (2) est satisfaite.

Puisque la seule solution au système est $\lambda_1 = 0$ et $\lambda_2 = 0$, les vecteurs $v_1 = (1,1,2)$ et $v_2 = (1,-1,0)$ sont linéairement indépendants.
Comme ils engendrent $\text{Im}(f)$ et sont linéairement indépendants, ils forment une base de $\text{Im}(f)$.
Une base de $\text{Im}(f)$ est donc $\{(1,1,2), (1,-1,0)\}$.
La dimension de $\text{Im}(f)$, notée $\text{dim}(\text{Im}(f))$, est le nombre de vecteurs dans cette base, soit 2.
Le rang de $f$, noté $\text{rg}(f)$, est par définition la dimension de son image.
$$\text{rg}(f) = \text{dim}(\text{Im}(f)) = 2$$

\subsection*{4. Vérifier le théorème du rang.}

Le théorème du rang est un résultat fondamental en algèbre linéaire qui relie la dimension de l'espace de départ, la dimension du noyau, et la dimension de l'image d'une application linéaire. Pour une application linéaire $f: E \to F$, où $E$ est un espace vectoriel de dimension finie, le théorème stipule que :
$$\text{dim}(E) = \text{dim}(\text{ker}(f)) + \text{dim}(\text{Im}(f))$$

Appliquons ce théorème à l'application linéaire $f: \mathbb{R}^2 \to \mathbb{R}^3$ que nous avons étudiée :
L'espace de départ $E$ est $\mathbb{R}^2$. Sa dimension est $\text{dim}(\mathbb{R}^2) = 2$.
D'après nos calculs de la partie 2, la dimension du noyau de $f$ est $\text{dim}(\text{ker}(f)) = 0$.
D'après nos calculs de la partie 3, la dimension de l'image de $f$ est $\text{dim}(\text{Im}(f)) = 2$.

Vérifions l'égalité du théorème du rang en substituant ces valeurs :
$$\text{dim}(\text{ker}(f)) + \text{dim}(\text{Im}(f)) = 0 + 2$$
$$\text{dim}(\text{ker}(f)) + \text{dim}(\text{Im}(f)) = 2$$
Nous comparons cette somme à la dimension de l'espace de départ :
$$\text{dim}(E) = 2$$
L'égalité est vérifiée :
$$2 = 2$$
Le théorème du rang est bien vérifié pour l'application linéaire $f$.

\newpage
\chapter*{Exercice 2 : Application linéaire de dérivation (Difficulté : *)}

\section*{Énoncé du problème}

Soit $\mathbb{K}$ le corps des nombres réels, $\mathbb{K} = \mathbb{R}$.
On considère l'espace vectoriel $E = \mathbb{R}_2[X]$ des polynômes à coefficients réels de degré au plus 2. Autrement dit, $E = \{ a_0 + a_1 X + a_2 X^2 \mid a_0, a_1, a_2 \in \mathbb{R} \}$.

On définit l'application $D: E \to E$ par $D(P) = P'$, où $P'$ désigne le polynôme dérivé de $P$.

1.  Justifier que $D$ est une application linéaire.
2.  Déterminer le noyau de $D$, noté $\text{Ker}(D)$. En donner une base et sa dimension.
3.  Déterminer l'image de $D$, notée $\text{Im}(D)$. En donner une base et sa dimension.
4.  Vérifier le théorème du rang pour l'application $D$.

\section*{Correction détaillée}

\subsection*{Question 1 : Justifier que $D$ est une application linéaire.}

Pour démontrer que l'application $D: E \to E$ est une application linéaire, nous devons vérifier qu'elle satisfait deux propriétés fondamentales : l'additivité et l'homogénéité.

1.  \textbf{Additivité} : Pour tous polynômes $P, Q \in E$, nous devons montrer que $D(P+Q) = D(P) + D(Q)$.
    Soient $P(X)$ et $Q(X)$ deux polynômes arbitraires de l'espace vectoriel $E = \mathbb{R}_2[X]$.
    Par définition de $E$, ils s'écrivent sous la forme :
    $P(X) = a_0 + a_1 X + a_2 X^2$
    $Q(X) = b_0 + b_1 X + b_2 X^2$
    où $a_0, a_1, a_2, b_0, b_1, b_2$ sont des coefficients réels appartenant à $\mathbb{R}$.

    Calculons la somme $P(X) + Q(X)$ :
    $P(X) + Q(X) = (a_0 + a_1 X + a_2 X^2) + (b_0 + b_1 X + b_2 X^2)$
    $P(X) + Q(X) = (a_0+b_0) + (a_1+b_1)X + (a_2+b_2)X^2$
    Ce polynôme est bien un élément de $E$ car ses coefficients $(a_0+b_0), (a_1+b_1), (a_2+b_2)$ sont des réels.

    Appliquons l'opérateur $D$ à cette somme :
    $D(P+Q) = (P+Q)'$
    $D(P+Q) = ((a_0+b_0) + (a_1+b_1)X + (a_2+b_2)X^2)'$
    En utilisant les règles de dérivation des polynômes, la dérivée est :
    $D(P+Q) = (a_1+b_1) + 2(a_2+b_2)X$

    Calculons maintenant $D(P)$ et $D(Q)$ séparément :
    $D(P) = P' = (a_0 + a_1 X + a_2 X^2)' = a_1 + 2a_2 X$
    $D(Q) = Q' = (b_0 + b_1 X + b_2 X^2)' = b_1 + 2b_2 X$

    Calculons la somme $D(P) + D(Q)$ :
    $D(P) + D(Q) = (a_1 + 2a_2 X) + (b_1 + 2b_2 X)$
    $D(P) + D(Q) = (a_1+b_1) + (2a_2+2b_2)X$
    $D(P) + D(Q) = (a_1+b_1) + 2(a_2+b_2)X$

    En comparant les expressions obtenues pour $D(P+Q)$ et $D(P)+D(Q)$, nous constatons qu'elles sont identiques :
    $D(P+Q) = (a_1+b_1) + 2(a_2+b_2)X$
    $D(P) + D(Q) = (a_1+b_1) + 2(a_2+b_2)X$
    Donc, la propriété d'additivité $D(P+Q) = D(P) + D(Q)$ est vérifiée.

2.  \textbf{Homogénéité} : Pour tout polynôme $P \in E$ et tout scalaire $\lambda \in \mathbb{R}$, nous devons montrer que $D(\lambda P) = \lambda D(P)$.
    Soit $P(X) = a_0 + a_1 X + a_2 X^2 \in E$, avec $a_0, a_1, a_2 \in \mathbb{R}$.
    Soit $\lambda \in \mathbb{R}$ un scalaire arbitraire.

    Calculons le produit scalaire $\lambda P(X)$ :
    $\lambda P(X) = \lambda (a_0 + a_1 X + a_2 X^2)$
    $\lambda P(X) = (\lambda a_0) + (\lambda a_1)X + (\lambda a_2)X^2$
    Ce polynôme est bien un élément de $E$ car ses coefficients $(\lambda a_0), (\lambda a_1), (\lambda a_2)$ sont des réels.

    Appliquons l'opérateur $D$ à ce produit scalaire :
    $D(\lambda P) = (\lambda P)'$
    $D(\lambda P) = ((\lambda a_0) + (\lambda a_1)X + (\lambda a_2)X^2)'$
    En utilisant les règles de dérivation des polynômes, la dérivée est :
    $D(\lambda P) = (\lambda a_1) + 2(\lambda a_2)X$

    Calculons maintenant $\lambda D(P)$ :
    $D(P) = P' = (a_0 + a_1 X + a_2 X^2)' = a_1 + 2a_2 X$
    $\lambda D(P) = \lambda (a_1 + 2a_2 X)$
    $\lambda D(P) = (\lambda a_1) + (\lambda \cdot 2a_2)X$
    $\lambda D(P) = (\lambda a_1) + 2(\lambda a_2)X$

    En comparant les expressions obtenues pour $D(\lambda P)$ et $\lambda D(P)$, nous constatons qu'elles sont identiques :
    $D(\lambda P) = (\lambda a_1) + 2(\lambda a_2)X$
    $\lambda D(P) = (\lambda a_1) + 2(\lambda a_2)X$
    Donc, la propriété d'homogénéité $D(\lambda P) = \lambda D(P)$ est vérifiée.

Puisque l'application $D$ satisfait à la fois la propriété d'additivité et la propriété d'homogénéité, nous concluons que $D$ est une application linéaire de $E$ dans $E$.

\subsection*{Question 2 : Déterminer le noyau de $D$, $\text{Ker}(D)$. En donner une base et sa dimension.}

Par définition, le noyau de l'application linéaire $D: E \to E$, noté $\text{Ker}(D)$, est l'ensemble de tous les polynômes $P \in E$ dont l'image par $D$ est le polynôme nul de $E$. Le polynôme nul de $E$, noté $0_E$, est le polynôme $0 + 0X + 0X^2$.

Formellement :
$\text{Ker}(D) = \{ P \in E \mid D(P) = 0_E \}$

Soit un polynôme $P(X) \in E$. Il s'écrit sous la forme générale :
$P(X) = a_0 + a_1 X + a_2 X^2$, où $a_0, a_1, a_2 \in \mathbb{R}$.

L'application $D$ est la dérivation, donc $D(P) = P'$.
Calculons la dérivée de $P(X)$ :
$P'(X) = (a_0 + a_1 X + a_2 X^2)'$
$P'(X) = 0 + a_1 \cdot 1 + a_2 \cdot (2X)$
$P'(X) = a_1 + 2a_2 X$

Pour que $P(X)$ appartienne au noyau de $D$, il faut que $P'(X) = 0_E$.
$P'(X) = 0_E \iff a_1 + 2a_2 X = 0 + 0X$

Par identification des coefficients des polynômes, pour que deux polynômes soient égaux, leurs coefficients respectifs doivent être égaux.
Le coefficient du terme constant ($X^0$) de $P'(X)$ est $a_1$. Le coefficient du terme constant de $0_E$ est $0$.
Donc, $a_1 = 0$.

Le coefficient du terme en $X$ ($X^1$) de $P'(X)$ est $2a_2$. Le coefficient du terme en $X$ de $0_E$ est $0$.
Donc, $2a_2 = 0$, ce qui implique $a_2 = 0$.

Ainsi, pour que $P(X)$ soit dans $\text{Ker}(D)$, les coefficients $a_1$ et $a_2$ doivent être nuls. Le coefficient $a_0$ n'est soumis à aucune contrainte.
Le polynôme $P(X)$ prend donc la forme :
$P(X) = a_0 + 0 \cdot X + 0 \cdot X^2$
$P(X) = a_0$

Le noyau de $D$ est l'ensemble de tous les polynômes constants :
$\text{Ker}(D) = \{ a_0 \mid a_0 \in \mathbb{R} \}$
Cet ensemble est l'espace vectoriel des polynômes de degré au plus 0, noté $\mathbb{R}_0[X]$.

\textbf{Détermination d'une base de $\text{Ker}(D)$ :}
Tout polynôme $P(X) \in \text{Ker}(D)$ peut s'écrire $P(X) = a_0 \cdot 1$, où $1$ est le polynôme constant égal à $1$.
Cela signifie que le polynôme $1$ engendre l'espace vectoriel $\text{Ker}(D)$.
Pour vérifier que $\{1\}$ est une base, nous devons également montrer qu'il est linéairement indépendant.
Considérons une combinaison linéaire du polynôme $1$ égale au polynôme nul :
$\alpha \cdot 1 = 0_E$, où $\alpha \in \mathbb{R}$.
Cette équation se réduit à $\alpha = 0$.
Puisque la seule solution est $\alpha = 0$, le polynôme $1$ est linéairement indépendant.
Par conséquent, l'ensemble $B_{\text{Ker}(D)} = \{1\}$ est une base de $\text{Ker}(D)$.

\textbf{Détermination de la dimension de $\text{Ker}(D)$ :}
La dimension d'un espace vectoriel est le nombre de vecteurs dans n'importe laquelle de ses bases.
Puisque la base $B_{\text{Ker}(D)}$ contient un seul vecteur, la dimension du noyau est :
$\dim(\text{Ker}(D)) = 1$.

\subsection*{Question 3 : Déterminer l'image de $D$, $\text{Im}(D)$. En donner une base et sa dimension.}

Par définition, l'image de l'application linéaire $D: E \to E$, notée $\text{Im}(D)$, est l'ensemble de tous les polynômes $Q \in E$ pour lesquels il existe au moins un polynôme $P \in E$ tel que $D(P) = Q$.

Formellement :
$\text{Im}(D) = \{ Q \in E \mid \exists P \in E \text{ tel que } D(P) = Q \}$

Soit un polynôme $P(X) \in E$. Il s'écrit sous la forme :
$P(X) = a_0 + a_1 X + a_2 X^2$, où $a_0, a_1, a_2 \in \mathbb{R}$.

L'image de $P(X)$ par $D$ est $D(P) = P'(X)$.
$D(P) = P'(X) = a_1 + 2a_2 X$.

Tout polynôme $Q(X)$ dans l'image de $D$ est donc de la forme $a_1 + 2a_2 X$.
Posons $b_0 = a_1$ et $b_1 = 2a_2$. Puisque $a_1$ et $a_2$ peuvent prendre n'importe quelle valeur réelle, $b_0$ peut être n'importe quel réel, et $b_1$ peut être n'importe quel réel.
Ainsi, tout polynôme dans l'image est de la forme $b_0 + b_1 X$, où $b_0, b_1 \in \mathbb{R}$.
Ceci correspond à la définition de l'espace vectoriel des polynômes de degré au plus 1, noté $\mathbb{R}_1[X]$.
Donc, nous avons montré que $\text{Im}(D) \subseteq \mathbb{R}_1[X]$.

Pour montrer l'égalité $\text{Im}(D) = \mathbb{R}_1[X]$, nous devons également prouver que $\mathbb{R}_1[X] \subseteq \text{Im}(D)$. C'est-à-dire, pour tout polynôme $Q(X) \in \mathbb{R}_1[X]$, il existe un polynôme $P(X) \in E$ tel que $D(P) = Q$.

Soit $Q(X)$ un polynôme arbitraire dans $\mathbb{R}_1[X]$. Il s'écrit sous la forme :
$Q(X) = b_0 + b_1 X$, où $b_0, b_1 \in \mathbb{R}$.

Nous cherchons un polynôme $P(X) = c_0 + c_1 X + c_2 X^2 \in E$ (avec $c_0, c_1, c_2 \in \mathbb{R}$) tel que $P'(X) = Q(X)$.
La dérivée de $P(X)$ est $P'(X) = c_1 + 2c_2 X$.
Nous voulons que $c_1 + 2c_2 X = b_0 + b_1 X$.

Par identification des coefficients :
Le coefficient du terme constant : $c_1 = b_0$.
Le coefficient du terme en $X$ : $2c_2 = b_1$, ce qui implique $c_2 = \frac{b_1}{2}$.

Le coefficient $c_0$ n'est pas déterminé par cette équation. Nous pouvons choisir n'importe quelle valeur réelle pour $c_0$. Par exemple, choisissons $c_0 = 0$.
Alors, le polynôme $P(X)$ est :
$P(X) = 0 + b_0 X + \frac{b_1}{2} X^2 = b_0 X + \frac{b_1}{2} X^2$.

Puisque $b_0 \in \mathbb{R}$ et $b_1 \in \mathbb{R}$, il s'ensuit que $b_0 \in \mathbb{R}$ et $\frac{b_1}{2} \in \mathbb{R}$.
Le polynôme $P(X) = b_0 X + \frac{b_1}{2} X^2$ est bien un polynôme de degré au plus 2 à coefficients réels, donc $P(X) \in E$.
Et sa dérivée est $D(P) = (b_0 X + \frac{b_1}{2} X^2)' = b_0 + 2 \cdot \frac{b_1}{2} X = b_0 + b_1 X = Q(X)$.
Ceci prouve que tout polynôme de $\mathbb{R}_1[X]$ est l'image d'un polynôme de $E$ par $D$.
Donc, $\mathbb{R}_1[X] \subseteq \text{Im}(D)$.

Ayant montré $\text{Im}(D) \subseteq \mathbb{R}_1[X]$ et $\mathbb{R}_1[X] \subseteq \text{Im}(D)$, nous concluons que :
$\text{Im}(D) = \mathbb{R}_1[X]$.

\textbf{Détermination d'une base de $\text{Im}(D)$ :}
L'espace vectoriel $\mathbb{R}_1[X]$ est l'ensemble des polynômes de degré au plus 1.
Les polynômes $1$ et $X$ sont des éléments de $\mathbb{R}_1[X]$.
Pour montrer que $B_{\text{Im}(D)} = \{1, X\}$ est une base de $\text{Im}(D)$, nous devons prouver qu'ils engendrent $\text{Im}(D)$ et qu'ils sont linéairement indépendants.

1.  \textbf{Engendrement} : Tout polynôme $Q(X) = b_0 + b_1 X \in \mathbb{R}_1[X]$ peut être écrit comme une combinaison linéaire de $1$ et $X$ :
    $Q(X) = b_0 \cdot 1 + b_1 \cdot X$.
    Donc, $\{1, X\}$ engendre $\text{Im}(D)$.

2.  \textbf{Indépendance linéaire} : Considérons une combinaison linéaire de $1$ et $X$ égale au polynôme nul $0_E$:
    $\alpha_0 \cdot 1 + \alpha_1 \cdot X = 0_E$, où $\alpha_0, \alpha_1 \in \mathbb{R}$.
    Cette équation s'écrit $\alpha_0 + \alpha_1 X = 0 + 0X$.
    Par identification des coefficients, nous obtenons :
    $\alpha_0 = 0$
    $\alpha_1 = 0$
    Puisque la seule solution est $\alpha_0 = 0$ et $\alpha_1 = 0$, les polynômes $1$ et $X$ sont linéairement indépendants.

Par conséquent, l'ensemble $B_{\text{Im}(D)} = \{1, X\}$ est une base de $\text{Im}(D)$.

\textbf{Détermination de la dimension de $\text{Im}(D)$ :}
La dimension de l'image est le nombre de vecteurs dans sa base.
Puisque la base $B_{\text{Im}(D)}$ contient deux vecteurs, la dimension de l'image est :
$\dim(\text{Im}(D)) = 2$.

\subsection*{Question 4 : Vérifier le théorème du rang pour l'application $D$.}

Le théorème du rang est un résultat fondamental en algèbre linéaire qui relie la dimension de l'espace de départ, la dimension du noyau et la dimension de l'image d'une application linéaire. Pour une application linéaire $L: V \to W$, le théorème du rang stipule que :
$\dim(V) = \dim(\text{Ker}(L)) + \dim(\text{Im}(L))$.

Dans notre cas, l'application linéaire est $D: E \to E$. L'espace de départ est $V = E = \mathbb{R}_2[X]$.

1.  \textbf{Détermination de la dimension de l'espace de départ $E$ :}
    L'espace vectoriel $E = \mathbb{R}_2[X]$ est l'ensemble des polynômes de degré au plus 2.
    Une base canonique de $\mathbb{R}_2[X]$ est l'ensemble des polynômes $\{1, X, X^2\}$.
    Cet ensemble est composé de trois vecteurs linéairement indépendants qui engendrent $E$.
    Par conséquent, la dimension de l'espace de départ $E$ est :
    $\dim(E) = 3$.

2.  \textbf{Récupération de la dimension du noyau $\text{Ker}(D)$ :}
    D'après les résultats de la Question 2, nous avons déterminé que $\dim(\text{Ker}(D)) = 1$.

3.  \textbf{Récupération de la dimension de l'image $\text{Im}(D)$ :}
    D'après les résultats de la Question 3, nous avons déterminé que $\dim(\text{Im}(D)) = 2$.

4.  \textbf{Vérification du théorème du rang :}
    Nous devons vérifier si l'égalité $\dim(E) = \dim(\text{Ker}(D)) + \dim(\text{Im}(D))$ est satisfaite avec les valeurs obtenues.
    Substituons les dimensions calculées dans l'équation du théorème du rang :
    $\dim(\text{Ker}(D)) + \dim(\text{Im}(D)) = 1 + 2$
    $\dim(\text{Ker}(D)) + \dim(\text{Im}(D)) = 3$

    Nous comparons ce résultat avec la dimension de l'espace de départ :
    $\dim(E) = 3$.

    Nous constatons que :
    $\dim(E) = 3$
    $\dim(\text{Ker}(D)) + \dim(\text{Im}(D)) = 3$

    L'égalité $\dim(E) = \dim(\text{Ker}(D)) + \dim(\text{Im}(D))$ est donc satisfaite.
    Par conséquent, le théorème du rang est vérifié pour l'application linéaire $D$.

\newpage

\chapter*{Exercice 3 : Noyau, Image, Théorème du Rang avec Paramètre (Difficulté : **)}

Soit $\mathbb{K} = \mathbb{R}$ le corps des nombres réels.
Soit $E = \mathbb{R}^3$ l'espace vectoriel des triplets de nombres réels, muni de sa base canonique $\mathcal{B} = (e_1, e_2, e_3)$, où $e_1 = (1, 0, 0)$, $e_2 = (0, 1, 0)$ et $e_3 = (0, 0, 1)$.

On considère l'application linéaire $f_a: E \to E$ définie, pour tout paramètre $a \in \mathbb{R}$, par sa matrice $M_a \in \mathcal{M}_3(\mathbb{R})$ dans la base canonique $\mathcal{B}$ :
$$M_a = \begin{pmatrix} 1 & 1 & 1 \\ 1 & a & 1 \\ 1 & 1 & a \end{pmatrix}$$

1.  Déterminer les valeurs du paramètre $a \in \mathbb{R}$ pour lesquelles l'application $f_a$ est un automorphisme de $E$.
2.  Pour les valeurs de $a$ pour lesquelles $f_a$ n'est pas un automorphisme, déterminer une base du noyau $\ker(f_a)$ et une base de l'image $\text{Im}(f_a)$.
3.  Dans tous les cas (que $f_a$ soit un automorphisme ou non), vérifier le théorème du rang pour l'application $f_a$.

\section*{Correction détaillée}

\subsection*{Question 1 : Détermination des valeurs de $a$ pour lesquelles $f_a$ est un automorphisme}

Soit $f_a: E \to E$ une application linéaire. Par définition, $f_a$ est un automorphisme de l'espace vectoriel $E$ si et seulement si $f_a$ est une application linéaire bijective de $E$ vers $E$.
Une application linéaire $f_a$ est un automorphisme si et seulement si sa matrice représentative $M_a$ dans une base donnée (ici, la base canonique $\mathcal{B}$) est une matrice inversible.
Pour une matrice carrée $M_a \in \mathcal{M}_3(\mathbb{R})$, la condition d'inversibilité est équivalente à la condition que son déterminant soit non nul, c'est-à-dire $\det(M_a) \neq 0$.

Calculons le déterminant de la matrice $M_a$:
$$M_a = \begin{pmatrix} 1 & 1 & 1 \\ 1 & a & 1 \\ 1 & 1 & a \end{pmatrix}$$

Pour simplifier le calcul du déterminant, nous allons effectuer des opérations élémentaires sur les lignes de la matrice. Ces opérations ne modifient pas la valeur du déterminant. L'objectif est de transformer la matrice en une matrice triangulaire supérieure.
Appliquons les opérations suivantes :
1.  $L_2 \leftarrow L_2 - L_1$ (soustraction de la première ligne à la deuxième ligne)
2.  $L_3 \leftarrow L_3 - L_1$ (soustraction de la première ligne à la troisième ligne)

Le déterminant de $M_a$ est alors :
$$ \det(M_a) = \det \begin{pmatrix} 1 & 1 & 1 \\ 1-1 & a-1 & 1-1 \\ 1-1 & 1-1 & a-1 \end{pmatrix} $$
En effectuant les soustractions, nous obtenons la matrice suivante :
$$ \det(M_a) = \det \begin{pmatrix} 1 & 1 & 1 \\ 0 & a-1 & 0 \\ 0 & 0 & a-1 \end{pmatrix} $$

La matrice obtenue est une matrice triangulaire supérieure. Le déterminant d'une matrice triangulaire (qu'elle soit supérieure ou inférieure) est égal au produit de ses éléments diagonaux.
Par conséquent :
$$ \det(M_a) = 1 \cdot (a-1) \cdot (a-1) $$
$$ \det(M_a) = (a-1)^2 $$

L'application $f_a$ est un automorphisme de $E$ si et seulement si $\det(M_a) \neq 0$.
Nous devons donc résoudre l'inéquation :
$$ (a-1)^2 \neq 0 $$
Cette inéquation est vérifiée si et seulement si la base de la puissance n'est pas nulle :
$$ a-1 \neq 0 $$
En ajoutant $1$ aux deux membres de l'inéquation, nous obtenons :
$$ a \neq 1 $$

Par conséquent, l'application linéaire $f_a$ est un automorphisme de $E$ si et seulement si le paramètre $a \in \mathbb{R}$ est différent de $1$.

\subsection*{Question 2 : Bases du noyau et de l'image pour $f_a$ lorsque $f_a$ n'est pas un automorphisme}

D'après les résultats de la Question 1, l'application $f_a$ n'est pas un automorphisme de $E$ si et seulement si $a=1$. Nous allons donc étudier ce cas spécifique.

Pour $a=1$, la matrice de l'application $f_1$ est obtenue en substituant la valeur $a=1$ dans l'expression de $M_a$:
$$M_1 = \begin{pmatrix} 1 & 1 & 1 \\ 1 & 1 & 1 \\ 1 & 1 & 1 \end{pmatrix}$$

\subsubsection*{Détermination d'une base du noyau $\ker(f_1)$}

Le noyau de l'application linéaire $f_1$, noté $\ker(f_1)$, est l'ensemble de tous les vecteurs $X \in E = \mathbb{R}^3$ tels que $f_1(X) = 0_E$, où $0_E$ est le vecteur nul de $E$.
En termes matriciels, si $X = \begin{pmatrix} x \\ y \\ z \end{pmatrix}$ est un vecteur de $\mathbb{R}^3$, alors $X \in \ker(f_1)$ si et seulement si $M_1 X = 0_E$.
Ceci conduit au système d'équations linéaires homogènes :
$$ \begin{pmatrix} 1 & 1 & 1 \\ 1 & 1 & 1 \\ 1 & 1 & 1 \end{pmatrix} \begin{pmatrix} x \\ y \\ z \end{pmatrix} = \begin{pmatrix} 0 \\ 0 \\ 0 \end{pmatrix} $$
En effectuant le produit matriciel, nous obtenons le système suivant :
$$ \begin{cases} 1x + 1y + 1z = 0 \\ 1x + 1y + 1z = 0 \\ 1x + 1y + 1z = 0 \end{cases} $$
Les trois équations de ce système sont identiques. Par conséquent, le système se réduit à une seule équation linéaire :
$$ x + y + z = 0 $$
Cette équation définit un plan passant par l'origine dans l'espace vectoriel $\mathbb{R}^3$. Pour trouver une base du noyau, nous exprimons une variable en fonction des autres. Choisissons d'exprimer $z$ en fonction de $x$ et $y$:
$$ z = -x - y $$
Un vecteur $X \in \ker(f_1)$ peut donc s'écrire sous la forme :
$$ X = \begin{pmatrix} x \\ y \\ -x-y \end{pmatrix} $$
Nous pouvons décomposer ce vecteur en une combinaison linéaire de vecteurs où $x$ et $y$ sont les coefficients (scalaires) :
$$ X = \begin{pmatrix} x \\ 0 \\ -x \end{pmatrix} + \begin{pmatrix} 0 \\ y \\ -y \end{pmatrix} $$
$$ X = x \begin{pmatrix} 1 \\ 0 \\ -1 \end{pmatrix} + y \begin{pmatrix} 0 \\ 1 \\ -1 \end{pmatrix} $$
Soient les vecteurs $v_1 = \begin{pmatrix} 1 \\ 0 \\ -1 \end{pmatrix}$ et $v_2 = \begin{pmatrix} 0 \\ 1 \\ -1 \end{pmatrix}$.
Ces vecteurs $v_1$ et $v_2$ engendrent le noyau $\ker(f_1)$, c'est-à-dire $\ker(f_1) = \text{Vect}(v_1, v_2)$.
Pour qu'ils forment une base du noyau, il faut également qu'ils soient linéairement indépendants.
Considérons une combinaison linéaire nulle de $v_1$ et $v_2$ avec des scalaires $\alpha, \beta \in \mathbb{R}$:
$$ \alpha v_1 + \beta v_2 = 0_E $$
$$ \alpha \begin{pmatrix} 1 \\ 0 \\ -1 \end{pmatrix} + \beta \begin{pmatrix} 0 \\ 1 \\ -1 \end{pmatrix} = \begin{pmatrix} 0 \\ 0 \\ 0 \end{pmatrix} $$
En effectuant l'addition vectorielle, nous obtenons :
$$ \begin{pmatrix} \alpha \cdot 1 + \beta \cdot 0 \\ \alpha \cdot 0 + \beta \cdot 1 \\ \alpha \cdot (-1) + \beta \cdot (-1) \end{pmatrix} = \begin{pmatrix} 0 \\ 0 \\ 0 \end{pmatrix} $$
Ce qui conduit au système d'équations :
$$ \begin{cases} \alpha = 0 \\ \beta = 0 \\ -\alpha - \beta = 0 \end{cases} $$
Les deux premières équations impliquent directement $\alpha = 0$ et $\beta = 0$. La troisième équation est alors satisfaite : $-(0) - (0) = 0$.
Puisque la seule combinaison linéaire de $v_1$ et $v_2$ qui donne le vecteur nul est celle où tous les scalaires sont nuls, les vecteurs $v_1$ et $v_2$ sont linéairement indépendants.

Par conséquent, une base du noyau de $f_1$ est $\mathcal{B}_{\ker(f_1)} = \left\{ \begin{pmatrix} 1 \\ 0 \\ -1 \end{pmatrix}, \begin{pmatrix} 0 \\ 1 \\ -1 \end{pmatrix} \right\}$.
La dimension du noyau de $f_1$ est $\dim(\ker(f_1)) = 2$.

\subsubsection*{Détermination d'une base de l'image $\text{Im}(f_1)$}

L'image de l'application linéaire $f_1$, notée $\text{Im}(f_1)$, est l'espace vectoriel engendré par les vecteurs colonnes de la matrice $M_1$.
Les vecteurs colonnes de $M_1$ sont :
$$ C_1 = \begin{pmatrix} 1 \\ 1 \\ 1 \end{pmatrix}, \quad C_2 = \begin{pmatrix} 1 \\ 1 \\ 1 \end{pmatrix}, \quad C_3 = \begin{pmatrix} 1 \\ 1 \\ 1 \end{pmatrix} $$
L'image est l'ensemble des combinaisons linéaires de ces vecteurs colonnes :
$$ \text{Im}(f_1) = \text{Vect}(C_1, C_2, C_3) $$
Nous observons que les trois vecteurs colonnes sont identiques : $C_1 = C_2 = C_3$.
Par conséquent, l'espace engendré par ces trois vecteurs est le même que l'espace engendré par un seul d'entre eux.
$$ \text{Im}(f_1) = \text{Vect}\left(\begin{pmatrix} 1 \\ 1 \\ 1 \end{pmatrix}\right) $$
Soit le vecteur $w_1 = \begin{pmatrix} 1 \\ 1 \\ 1 \end{pmatrix}$. Ce vecteur est non nul. Un ensemble constitué d'un unique vecteur non nul est toujours linéairement indépendant.
Ainsi, le vecteur $w_1$ forme à lui seul une base de l'image de $f_1$.

Par conséquent, une base de l'image de $f_1$ est $\mathcal{B}_{\text{Im}(f_1)} = \left\{ \begin{pmatrix} 1 \\ 1 \\ 1 \end{pmatrix} \right\}$.
La dimension de l'image de $f_1$ est $\dim(\text{Im}(f_1)) = 1$.

\subsection*{Question 3 : Vérification du théorème du rang dans tous les cas}

Le théorème du rang est un résultat fondamental en algèbre linéaire qui relie la dimension de l'espace de départ, la dimension du noyau et la dimension de l'image d'une application linéaire.
Pour toute application linéaire $f: E \to F$, où $E$ et $F$ sont des espaces vectoriels de dimension finie, le théorème du rang stipule que :
$$ \dim(E) = \dim(\ker(f)) + \dim(\text{Im}(f)) $$
où $\dim(\text{Im}(f))$ est également appelé le rang de $f$, noté $\text{rang}(f)$.

Dans notre exercice, l'espace de départ est $E = \mathbb{R}^3$. Sa dimension est $\dim(E) = 3$.
Nous allons vérifier le théorème du rang pour les deux cas distincts du paramètre $a \in \mathbb{R}$.

\subsubsection*{Cas 1 : $a=1$ (où $f_1$ n'est pas un automorphisme)}

D'après les calculs effectués dans la Question 2 pour le cas $a=1$:
*   La dimension du noyau de $f_1$ est $\dim(\ker(f_1)) = 2$.
*   La dimension de l'image de $f_1$ est $\dim(\text{Im}(f_1)) = 1$.

Appliquons le théorème du rang avec ces valeurs :
$$ \dim(E) = \dim(\ker(f_1)) + \dim(\text{Im}(f_1)) $$
$$ 3 = 2 + 1 $$
$$ 3 = 3 $$
L'égalité est vérifiée. Le théorème du rang est donc confirmé pour l'application $f_1$ lorsque $a=1$.

\subsubsection*{Cas 2 : $a \neq 1$ (où $f_a$ est un automorphisme)}

D'après les résultats de la Question 1, si $a \neq 1$, l'application $f_a: E \to E$ est un automorphisme de $E$.
Par définition, un automorphisme est une application linéaire bijective.
1.  \textbf{Injectivité de $f_a$}: Puisque $f_a$ est injective, son noyau est réduit au seul vecteur nul de l'espace de départ $E$.
    $$ \ker(f_a) = \{0_E\} $$
    Par conséquent, la dimension du noyau de $f_a$ est :
    $$ \dim(\ker(f_a)) = 0 $$
2.  \textbf{Surjectivité de $f_a$}: Puisque $f_a$ est surjective et que l'espace d'arrivée est $E$, son image est égale à l'espace d'arrivée $E$.
    $$ \text{Im}(f_a) = E $$
    Par conséquent, la dimension de l'image de $f_a$ est :
    $$ \dim(\text{Im}(f_a)) = \dim(E) = 3 $$

Appliquons le théorème du rang avec ces valeurs pour $a \neq 1$:
$$ \dim(E) = \dim(\ker(f_a)) + \dim(\text{Im}(f_a)) $$
$$ 3 = 0 + 3 $$
$$ 3 = 3 $$
L'égalité est vérifiée. Le théorème du rang est donc confirmé pour l'application $f_a$ lorsque $a \neq 1$.

En conclusion, le théorème du rang est vérifié pour toutes les valeurs du paramètre $a \in \mathbb{R}$.

\newpage

\chapter*{Exercice 4 : Noyau, Image et Théorème du Rang (Difficulté : **)}

\section*{Énoncé du problème}

Soit $\mathbb{K}$ le corps des nombres réels, noté $\mathbb{K} = \mathbb{R}$.
On considère les $\mathbb{K}$-espaces vectoriels $E = \mathbb{R}^3$ et $F = \mathbb{R}^2$.
Soit $f$ l'application de $E$ vers $F$ définie pour tout vecteur $\mathbf{v} = (x,y,z) \in E$ par :
$$ f(x,y,z) = (x-y, y-z) $$
On admettra que $f$ est une application linéaire de $E$ dans $F$.

1.  Déterminer une base et la dimension du noyau de $f$, noté $\operatorname{Ker}(f)$.
2.  Déterminer une base et la dimension de l'image de $f$, notée $\operatorname{Im}(f)$.
3.  Vérifier le théorème du rang pour l'application linéaire $f$.

---

\section*{Correction détaillée}

\subsection*{Question 1 : Détermination du noyau de $f$}

Le noyau de l'application linéaire $f$, noté $\operatorname{Ker}(f)$, est défini comme l'ensemble des vecteurs $\mathbf{v} \in E$ tels que $f(\mathbf{v}) = \mathbf{0}_F$, où $\mathbf{0}_F$ représente le vecteur nul de l'espace vectoriel $F$.
Dans notre cas, l'espace de départ est $E = \mathbb{R}^3$ et l'espace d'arrivée est $F = \mathbb{R}^2$. Le vecteur nul de $F$ est $\mathbf{0}_F = (0,0)$.

Soit un vecteur $\mathbf{v} = (x,y,z) \in \mathbb{R}^3$.
L'appartenance de $\mathbf{v}$ au noyau $\operatorname{Ker}(f)$ est caractérisée par l'équation vectorielle $f(x,y,z) = (0,0)$.
En substituant la définition de l'application $f$, nous obtenons :
$$ (x-y, y-z) = (0,0) $$
Cette égalité de vecteurs dans $\mathbb{R}^2$ est équivalente au système d'équations linéaires suivant, en égalant les composantes correspondantes :
$$ \begin{cases} x - y = 0 & \quad (L_1) \\ y - z = 0 & \quad (L_2) \end{cases} $$

Résolvons ce système d'équations pour exprimer $x, y, z$ en fonction d'un paramètre.
À partir de l'équation $(L_1)$, $x - y = 0$:
En ajoutant le terme $y$ aux deux membres de l'équation, nous obtenons :
$$ x - y + y = 0 + y $$
$$ x = y \quad (L_3) $$

À partir de l'équation $(L_2)$, $y - z = 0$:
En ajoutant le terme $z$ aux deux membres de l'équation, nous obtenons :
$$ y - z + z = 0 + z $$
$$ y = z \quad (L_4) $$

En combinant les résultats des équations $(L_3)$ et $(L_4)$, nous déduisons que $x=y$ et $y=z$. Par la propriété de transitivité de l'égalité, cela implique que $x=y=z$.
Ainsi, tout vecteur $\mathbf{v} = (x,y,z)$ appartenant à $\operatorname{Ker}(f)$ doit satisfaire la condition $x = y = z$.
Les vecteurs du noyau sont donc de la forme $(x,x,x)$ pour tout scalaire $x \in \mathbb{R}$.

Nous pouvons factoriser le scalaire $x$ de ce vecteur pour exprimer sa structure :
$$ (x,x,x) = x \cdot (1,1,1) $$
Le noyau $\operatorname{Ker}(f)$ est donc l'ensemble de tous les multiples scalaires du vecteur $(1,1,1)$.
Par définition, cela signifie que $\operatorname{Ker}(f)$ est l'espace vectoriel engendré par le vecteur $(1,1,1)$.
$$ \operatorname{Ker}(f) = \operatorname{Vect}((1,1,1)) $$

Pour déterminer une base de $\operatorname{Ker}(f)$, nous devons trouver un ensemble de vecteurs qui est à la fois générateur de $\operatorname{Ker}(f)$ et linéairement indépendant.
L'ensemble $B_{\operatorname{Ker}(f)} = \{ (1,1,1) \}$ est un ensemble générateur de $\operatorname{Ker}(f)$ par construction.
Vérifions la liberté linéaire de cet ensemble. Soit $\alpha \in \mathbb{R}$ un scalaire tel que $\alpha \cdot (1,1,1) = \mathbf{0}_E$.
$$ (\alpha \cdot 1, \alpha \cdot 1, \alpha \cdot 1) = (0,0,0) $$
$$ (\alpha, \alpha, \alpha) = (0,0,0) $$
En égalant les composantes de ces vecteurs, nous obtenons $\alpha = 0$.
Puisque le seul scalaire $\alpha$ qui satisfait cette condition est $\alpha = 0$, le vecteur $(1,1,1)$ est linéairement indépendant.
De plus, le vecteur $(1,1,1)$ est non nul.
Par conséquent, l'ensemble $B_{\operatorname{Ker}(f)} = \{ (1,1,1) \}$ est une base de $\operatorname{Ker}(f)$.

La dimension du noyau de $f$, qui est le nombre de vecteurs dans une base de $\operatorname{Ker}(f)$, est $\dim(\operatorname{Ker}(f)) = 1$.

\subsection*{Question 2 : Détermination de l'image de $f$}

L'image de l'application linéaire $f$, notée $\operatorname{Im}(f)$, est définie comme l'ensemble des vecteurs $\mathbf{w} \in F$ pour lesquels il existe au moins un vecteur $\mathbf{v} \in E$ tel que $f(\mathbf{v}) = \mathbf{w}$.
Pour une application linéaire $f: E \to F$, l'image $\operatorname{Im}(f)$ est un sous-espace vectoriel de $F$.
De plus, l'image d'une application linéaire est engendrée par les images des vecteurs d'une base quelconque de l'espace de départ $E$.

Considérons la base canonique de l'espace vectoriel de départ $E = \mathbb{R}^3$, notée $B_E = \{ \mathbf{e}_1, \mathbf{e}_2, \mathbf{e}_3 \}$, où les vecteurs sont définis comme suit :
*   $\mathbf{e}_1 = (1,0,0) \in \mathbb{R}^3$
*   $\mathbf{e}_2 = (0,1,0) \in \mathbb{R}^3$
*   $\mathbf{e}_3 = (0,0,1) \in \mathbb{R}^3$

Calculons les images de ces vecteurs par l'application $f$:
*   Pour le vecteur $\mathbf{e}_1 = (1,0,0)$:
    $$ f(\mathbf{e}_1) = f(1,0,0) = (1-0, 0-0) = (1,0) $$
*   Pour le vecteur $\mathbf{e}_2 = (0,1,0)$:
    $$ f(\mathbf{e}_2) = f(0,1,0) = (0-1, 1-0) = (-1,1) $$
*   Pour le vecteur $\mathbf{e}_3 = (0,0,1)$:
    $$ f(\mathbf{e}_3) = f(0,0,1) = (0-0, 0-1) = (0,-1) $$

L'image de $f$ est l'espace vectoriel engendré par ces trois vecteurs images :
$$ \operatorname{Im}(f) = \operatorname{Vect}((1,0), (-1,1), (0,-1)) $$
Ces vecteurs appartiennent à l'espace vectoriel $F = \mathbb{R}^2$. La dimension de $F$ est $\dim(F) = 2$. Par conséquent, la dimension de $\operatorname{Im}(f)$ ne peut pas excéder 2, c'est-à-dire $\dim(\operatorname{Im}(f)) \le 2$.
Nous devons extraire une base de l'ensemble de vecteurs générateurs $S = \{ (1,0), (-1,1), (0,-1) \}$.

Considérons les deux premiers vecteurs de l'ensemble $S$: $\mathbf{u}_1 = (1,0)$ et $\mathbf{u}_2 = (-1,1)$.
Vérifions si ces deux vecteurs sont linéairement indépendants dans $\mathbb{R}^2$.
Soient $\alpha, \beta \in \mathbb{R}$ des scalaires tels que $\alpha \mathbf{u}_1 + \beta \mathbf{u}_2 = \mathbf{0}_F$.
$$ \alpha(1,0) + \beta(-1,1) = (0,0) $$
Effectuons la multiplication scalaire pour chaque terme :
$$ (\alpha \cdot 1, \alpha \cdot 0) + (\beta \cdot (-1), \beta \cdot 1) = (0,0) $$
$$ (\alpha, 0) + (-\beta, \beta) = (0,0) $$
Effectuons l'addition vectorielle des deux vecteurs résultants :
$$ (\alpha - \beta, 0 + \beta) = (0,0) $$
$$ (\alpha - \beta, \beta) = (0,0) $$
Cette égalité de vecteurs est équivalente au système d'équations linéaires suivant, en égalant les composantes :
$$ \begin{cases} \alpha - \beta = 0 & \quad (L_5) \\ \beta = 0 & \quad (L_6) \end{cases} $$
De l'équation $(L_6)$, nous avons directement $\beta = 0$.
Substituons cette valeur de $\beta$ dans l'équation $(L_5)$:
$$ \alpha - 0 = 0 $$
$$ \alpha = 0 $$
Puisque les seuls scalaires $\alpha$ et $\beta$ qui satisfont l'équation $\alpha \mathbf{u}_1 + \beta \mathbf{u}_2 = \mathbf{0}_F$ sont $\alpha = 0$ et $\beta = 0$, les vecteurs $\mathbf{u}_1 = (1,0)$ et $\mathbf{u}_2 = (-1,1)$ sont linéairement indépendants.

Nous avons trouvé deux vecteurs linéairement indépendants, $(1,0)$ et $(-1,1)$, qui appartiennent à l'espace vectoriel $F = \mathbb{R}^2$.
Puisque la dimension de $\mathbb{R}^2$ est $\dim(\mathbb{R}^2) = 2$, tout ensemble de 2 vecteurs linéairement indépendants dans $\mathbb{R}^2$ forme une base de $\mathbb{R}^2$.
Par conséquent, l'ensemble $B_{\operatorname{Im}(f)} = \{ (1,0), (-1,1) \}$ est une base de $\operatorname{Im}(f)$.
De plus, cela implique que $\operatorname{Im}(f) = \mathbb{R}^2$.

Pour une complétude rigoureuse, nous pouvons vérifier que le troisième vecteur image, $(0,-1)$, est une combinaison linéaire des vecteurs de la base $B_{\operatorname{Im}(f)}$.
Cherchons des scalaires $c_1, c_2 \in \mathbb{R}$ tels que $c_1(1,0) + c_2(-1,1) = (0,-1)$.
$$ (c_1 \cdot 1 + c_2 \cdot (-1), c_1 \cdot 0 + c_2 \cdot 1) = (0,-1) $$
$$ (c_1 - c_2, c_2) = (0,-1) $$
Ce qui conduit au système d'équations :
$$ \begin{cases} c_1 - c_2 = 0 \\ c_2 = -1 \end{cases} $$
De la seconde équation, nous avons $c_2 = -1$.
Substituons la valeur de $c_2 = -1$ dans la première équation :
$$ c_1 - (-1) = 0 $$
$$ c_1 + 1 = 0 $$
$$ c_1 = -1 $$
Ainsi, nous avons $(-1)(1,0) + (-1)(-1,1) = (-1,0) + (1,-1) = (0,-1)$.
Le vecteur $(0,-1)$ est bien une combinaison linéaire de $(1,0)$ et $(-1,1)$, ce qui confirme qu'il est redondant pour la génération de l'espace et que $B_{\operatorname{Im}(f)}$ est bien une base.

La dimension de l'image de $f$, qui est le nombre de vecteurs dans la base $B_{\operatorname{Im}(f)}$, est $\dim(\operatorname{Im}(f)) = 2$.
Le rang de l'application linéaire $f$, noté $\operatorname{rg}(f)$, est défini comme la dimension de son image.
Donc, $\operatorname{rg}(f) = \dim(\operatorname{Im}(f)) = 2$.

\subsection*{Question 3 : Vérification du théorème du rang}

Le théorème du rang est un résultat fondamental de l'algèbre linéaire qui établit une relation entre la dimension de l'espace de départ, la dimension du noyau et la dimension de l'image d'une application linéaire.
Pour toute application linéaire $f: E \to F$, le théorème du rang stipule que :
$$ \dim(E) = \dim(\operatorname{Ker}(f)) + \dim(\operatorname{Im}(f)) $$

Appliquons ce théorème aux résultats obtenus pour l'application linéaire $f$ :
*   La dimension de l'espace vectoriel de départ $E = \mathbb{R}^3$ est $\dim(E) = 3$.
*   D'après les calculs détaillés de la Question 1, la dimension du noyau de $f$ est $\dim(\operatorname{Ker}(f)) = 1$.
*   D'après les calculs détaillés de la Question 2, la dimension de l'image de $f$ est $\dim(\operatorname{Im}(f)) = 2$.

Substituons ces valeurs numériques dans la formule du théorème du rang :
$$ 3 = 1 + 2 $$
Effectuons l'addition dans le membre de droite de l'égalité :
$$ 3 = 3 $$
L'égalité est vérifiée. Par conséquent, le théorème du rang est satisfait pour l'application linéaire $f$ considérée.

\newpage
\chapter*{Exercice 5 : Noyau, Image et Rang de la composition d'applications linéaires (Difficulté : ***)}

\textbf{Énoncé du Problème}

Soient $K$ un corps commutatif. Soient $E, F, G$ trois $K$-espaces vectoriels de dimensions finies.
Soient $f: E \to F$ et $g: F \to G$ deux applications linéaires.

1.  Démontrer que $\ker(f) \subseteq \ker(g \circ f)$ et que $\text{Im}(g \circ f) \subseteq \text{Im}(g)$.
2.  On considère l'application linéaire restreinte $g_f: \text{Im}(f) \to G$ définie par $g_f(y) = g(y)$ pour tout $y \in \text{Im}(f)$.
    a) Démontrer que $\ker(g_f) = \text{Im}(f) \cap \ker(g)$.
    b) Démontrer que $\text{Im}(g_f) = \text{Im}(g \circ f)$.
3.  En utilisant le théorème du rang pour l'application $g_f$ et l'application $f$, établir la relation suivante :
    $$ \dim(\ker(g \circ f)) = \dim(\ker(f)) + \dim(\text{Im}(f) \cap \ker(g)) $$
4.  En déduire la relation liant les rangs des applications linéaires :
    $$ \text{rg}(g \circ f) = \text{rg}(f) - \dim(\text{Im}(f) \cap \ker(g)) $$

---

\section*{Correction détaillée}

1.  \textbf{Démonstration de $\ker(f) \subseteq \ker(g \circ f)$ et $\text{Im}(g \circ f) \subseteq \text{Im}(g)$}

    *   \textbf{Pour $\ker(f) \subseteq \ker(g \circ f)$ :}
        Soit $x \in \ker(f)$. Par définition du noyau, cela signifie que $f(x) = 0_F$, où $0_F$ est le vecteur nul de $F$.
        Appliquons l'application $g$ à $f(x)$ :
        $(g \circ f)(x) = g(f(x)) = g(0_F)$.
        Puisque $g$ est une application linéaire, elle envoie le vecteur nul sur le vecteur nul : $g(0_F) = 0_G$.
        Ainsi, $(g \circ f)(x) = 0_G$.
        Par définition du noyau, $x \in \ker(g \circ f)$.
        Donc, tout élément de $\ker(f)$ est aussi un élément de $\ker(g \circ f)$, ce qui prouve $\ker(f) \subseteq \ker(g \circ f)$.

    *   \textbf{Pour $\text{Im}(g \circ f) \subseteq \text{Im}(g)$ :}
        Soit $z \in \text{Im}(g \circ f)$. Par définition de l'image, cela signifie qu'il existe un vecteur $x \in E$ tel que $z = (g \circ f)(x)$.
        On peut réécrire $z = g(f(x))$.
        Posons $y = f(x)$. Alors, par définition de l'image de $f$, $y \in \text{Im}(f)$.
        Nous avons $z = g(y)$.
        Puisque $y \in F$ et $z = g(y)$, par définition de l'image de $g$, $z \in \text{Im}(g)$.
        Donc, tout élément de $\text{Im}(g \circ f)$ est aussi un élément de $\text{Im}(g)$, ce qui prouve $\text{Im}(g \circ f) \subseteq \text{Im}(g)$.

2.  \textbf{Propriétés de l'application linéaire restreinte $g_f$}

    L'application $g_f: \text{Im}(f) \to G$ est définie par $g_f(y) = g(y)$ pour tout $y \in \text{Im}(f)$.
    Il est important de noter que $\text{Im}(f)$ est un sous-espace vectoriel de $F$, et $g_f$ est la restriction de $g$ à ce sous-espace.

    a) \textbf{Démonstration de $\ker(g_f) = \text{Im}(f) \cap \ker(g)$}

        *   \textbf{Inclusion $\subseteq$ :}
            Soit $y \in \ker(g_f)$. Par définition du noyau de $g_f$, cela signifie que $y \in \text{Im}(f)$ et $g_f(y) = 0_G$.
            Par définition de $g_f$, on a $g_f(y) = g(y)$. Donc, $g(y) = 0_G$.
            Puisque $g(y) = 0_G$, par définition du noyau de $g$, on a $y \in \ker(g)$.
            Ainsi, $y$ appartient à la fois à $\text{Im}(f)$ et à $\ker(g)$. Donc $y \in \text{Im}(f) \cap \ker(g)$.
            Ceci prouve $\ker(g_f) \subseteq \text{Im}(f) \cap \ker(g)$.

        *   \textbf{Inclusion $\supseteq$ :}
            Soit $y \in \text{Im}(f) \cap \ker(g)$. Par définition de l'intersection, cela signifie que $y \in \text{Im}(f)$ et $y \in \ker(g)$.
            Puisque $y \in \text{Im}(f)$, l'application $g_f(y)$ est bien définie.
            Puisque $y \in \ker(g)$, par définition du noyau de $g$, on a $g(y) = 0_G$.
            Par définition de $g_f$, on a $g_f(y) = g(y)$. Donc, $g_f(y) = 0_G$.
            Ainsi, $y \in \ker(g_f)$.
            Ceci prouve $\text{Im}(f) \cap \ker(g) \subseteq \ker(g_f)$.

        Les deux inclusions étant établies, on conclut que $\ker(g_f) = \text{Im}(f) \cap \ker(g)$.

    b) \textbf{Démonstration de $\text{Im}(g_f) = \text{Im}(g \circ f)$}

        *   \textbf{Inclusion $\subseteq$ :}
            Soit $z \in \text{Im}(g_f)$. Par définition de l'image de $g_f$, cela signifie qu'il existe un vecteur $y \in \text{Im}(f)$ tel que $z = g_f(y)$.
            Par définition de $g_f$, on a $z = g(y)$.
            Puisque $y \in \text{Im}(f)$, par définition de l'image de $f$, il existe un vecteur $x \in E$ tel que $y = f(x)$.
            En substituant $y$, nous obtenons $z = g(f(x))$.
            Par définition de la composition d'applications, $z = (g \circ f)(x)$.
            Puisque $x \in E$ et $z = (g \circ f)(x)$, par définition de l'image de $g \circ f$, on a $z \in \text{Im}(g \circ f)$.
            Ceci prouve $\text{Im}(g_f) \subseteq \text{Im}(g \circ f)$.

        *   \textbf{Inclusion $\supseteq$ :}
            Soit $z \in \text{Im}(g \circ f)$. Par définition de l'image de $g \circ f$, cela signifie qu'il existe un vecteur $x \in E$ tel que $z = (g \circ f)(x)$.
            On peut réécrire $z = g(f(x))$.
            Posons $y = f(x)$. Alors, par définition de l'image de $f$, $y \in \text{Im}(f)$.
            Puisque $y \in \text{Im}(f)$, l'application $g_f(y)$ est bien définie, et par définition $g_f(y) = g(y)$.
            Ainsi, $z = g(f(x)) = g(y) = g_f(y)$.
            Puisque $y \in \text{Im}(f)$ et $z = g_f(y)$, par définition de l'image de $g_f$, on a $z \in \text{Im}(g_f)$.
            Ceci prouve $\text{Im}(g \circ f) \subseteq \text{Im}(g_f)$.

        Les deux inclusions étant établies, on conclut que $\text{Im}(g_f) = \text{Im}(g \circ f)$.

3.  \textbf{Établissement de la relation $\dim(\ker(g \circ f)) = \dim(\ker(f)) + \dim(\text{Im}(f) \cap \ker(g))$}

    Nous allons appliquer le théorème du rang aux applications linéaires $f$ et $g_f$.

    *   \textbf{Théorème du rang pour $g_f: \text{Im}(f) \to G$ :}
        Puisque $\text{Im}(f)$ est un sous-espace vectoriel de $F$ et est de dimension finie (car $F$ est de dimension finie), le théorème du rang s'applique à $g_f$.
        $$ \dim(\text{Im}(f)) = \dim(\ker(g_f)) + \dim(\text{Im}(g_f)) $$
        En utilisant les résultats de la question 2a) et 2b) :
        $$ \dim(\text{Im}(f)) = \dim(\text{Im}(f) \cap \ker(g)) + \dim(\text{Im}(g \circ f)) \quad (*)$$

    *   \textbf{Théorème du rang pour $f: E \to F$ :}
        $$ \dim(E) = \dim(\ker(f)) + \dim(\text{Im}(f)) $$
        On peut exprimer $\dim(\text{Im}(f))$ comme :
        $$ \dim(\text{Im}(f)) = \dim(E) - \dim(\ker(f)) \quad (**) $$

    *   \textbf{Théorème du rang pour $g \circ f: E \to G$ :}
        $$ \dim(E) = \dim(\ker(g \circ f)) + \dim(\text{Im}(g \circ f)) $$
        On peut exprimer $\dim(\text{Im}(g \circ f))$ comme :
        $$ \dim(\text{Im}(g \circ f)) = \dim(E) - \dim(\ker(g \circ f)) \quad (***) $$

    Substituons $(**)$ et $(***)$ dans l'équation $(*)$ :
    $$ (\dim(E) - \dim(\ker(f))) = \dim(\text{Im}(f) \cap \ker(g)) + (\dim(E) - \dim(\ker(g \circ f))) $$
    Simplifions en soustrayant $\dim(E)$ des deux côtés :
    $$ -\dim(\ker(f)) = \dim(\text{Im}(f) \cap \ker(g)) - \dim(\ker(g \circ f)) $$
    Réarrangeons les termes pour isoler $\dim(\ker(g \circ f))$ :
    $$ \dim(\ker(g \circ f)) = \dim(\ker(f)) + \dim(\text{Im}(f) \cap \ker(g)) $$
    Ceci établit la relation demandée.

4.  \textbf{Déduction de la relation $\text{rg}(g \circ f) = \text{rg}(f) - \dim(\text{Im}(f) \cap \ker(g))$}

    Par définition, le rang d'une application linéaire est la dimension de son image.
    Donc $\text{rg}(f) = \dim(\text{Im}(f))$ et $\text{rg}(g \circ f) = \dim(\text{Im}(g \circ f))$.

    Reprenons l'équation $(*)$ établie à la question 3 :
    $$ \dim(\text{Im}(f)) = \dim(\text{Im}(f) \cap \ker(g)) + \dim(\text{Im}(g \circ f)) $$
    En substituant les termes par leurs rangs correspondants :
    $$ \text{rg}(f) = \dim(\text{Im}(f) \cap \ker(g)) + \text{rg}(g \circ f) $$
    Pour obtenir la relation souhaitée, il suffit de réarranger les termes :
    $$ \text{rg}(g \circ f) = \text{rg}(f) - \dim(\text{Im}(f) \cap \ker(g)) $$
    Ceci conclut la déduction.

\newpage
\chapter*{Exercice 6 : L'opérateur de différence finie sur les polynômes (Difficulté : ***)}

\section*{Énoncé du problème}

Soit $n \in \mathbb{N}$ un entier naturel. On considère l'espace vectoriel réel $E = \mathbb{R}_n[X]$, l'ensemble des polynômes à coefficients réels de degré inférieur ou égal à $n$. La base canonique de $E$ est $B = (1, X, X^2, \dots, X^n)$.

On définit l'application $L: E \to E$ par :
$$L(P)(X) = P(X+1) - P(X)$$
pour tout polynôme $P \in E$.

1.  Démontrer que $L$ est une application linéaire.
2.  Déterminer le noyau $\mathrm{Ker}(L)$ de $L$. Donner une base de $\mathrm{Ker}(L)$ et sa dimension.
3.  Déterminer l'image $\mathrm{Im}(L)$ de $L$. Donner une base de $\mathrm{Im}(L)$ et sa dimension.
4.  Calculer le rang de $L$ et vérifier le théorème du rang.
5.  Déterminer la matrice $M_B(L)$ de $L$ dans la base canonique $B$ de $E$.

---

\section*{Correction détaillée}

\subsection*{Question 1 : Démontrer que $L$ est une application linéaire.}

Pour que $L$ soit linéaire, elle doit satisfaire deux propriétés :
a) $L(P+Q) = L(P) + L(Q)$ pour tous $P, Q \in E$.
b) $L(\lambda P) = \lambda L(P)$ pour tout $P \in E$ et tout scalaire $\lambda \in \mathbb{R}$.

Soient $P, Q \in E$ et $\lambda \in \mathbb{R}$.

a) Calculons $L(P+Q)(X)$:
$$L(P+Q)(X) = (P+Q)(X+1) - (P+Q)(X)$$
Par définition de l'addition de polynômes, $(P+Q)(X+1) = P(X+1) + Q(X+1)$ et $(P+Q)(X) = P(X) + Q(X)$.
Donc,
$$L(P+Q)(X) = (P(X+1) + Q(X+1)) - (P(X) + Q(X))$$
Regroupons les termes :
$$L(P+Q)(X) = (P(X+1) - P(X)) + (Q(X+1) - Q(X))$$
Par définition de $L$, ceci est égal à $L(P)(X) + L(Q)(X)$.
Ainsi, $L(P+Q) = L(P) + L(Q)$.

b) Calculons $L(\lambda P)(X)$:
$$L(\lambda P)(X) = (\lambda P)(X+1) - (\lambda P)(X)$$
Par définition de la multiplication d'un polynôme par un scalaire, $(\lambda P)(X+1) = \lambda P(X+1)$ et $(\lambda P)(X) = \lambda P(X)$.
Donc,
$$L(\lambda P)(X) = \lambda P(X+1) - \lambda P(X)$$
Mettons $\lambda$ en facteur :
$$L(\lambda P)(X) = \lambda (P(X+1) - P(X))$$
Par définition de $L$, ceci est égal à $\lambda L(P)(X)$.
Ainsi, $L(\lambda P) = \lambda L(P)$.

Les deux propriétés étant vérifiées, $L$ est bien une application linéaire.
De plus, si $P \in \mathbb{R}_n[X]$, alors $P(X+1)$ est également un polynôme de degré au plus $n$. La différence $P(X+1)-P(X)$ est donc un polynôme de degré au plus $n$. L'application $L$ est donc bien définie de $E$ vers $E$.

\subsection*{Question 2 : Déterminer le noyau $\mathrm{Ker}(L)$.}

Le noyau $\mathrm{Ker}(L)$ est l'ensemble des polynômes $P \in E$ tels que $L(P) = 0$.
$$P \in \mathrm{Ker}(L) \iff L(P)(X) = 0 \text{ pour tout } X \in \mathbb{R}$$
$$P(X+1) - P(X) = 0 \iff P(X+1) = P(X)$$
Un polynôme qui vérifie $P(X+1) = P(X)$ pour tout $X \in \mathbb{R}$ est un polynôme constant.
Pour le prouver, soit $P(X) = a_k X^k + a_{k-1} X^{k-1} + \dots + a_1 X + a_0$ un polynôme de degré $k \le n$, avec $a_k \neq 0$ si $k \ge 1$.
Alors $P(X+1) = a_k (X+1)^k + a_{k-1} (X+1)^{k-1} + \dots + a_1 (X+1) + a_0$.
Développons $(X+1)^k$ : $(X+1)^k = X^k + k X^{k-1} + \binom{k}{2} X^{k-2} + \dots + 1$.
$$P(X+1) = a_k (X^k + k X^{k-1} + \dots) + a_{k-1} (X^{k-1} + \dots) + \dots + a_0$$
$$P(X+1) - P(X) = a_k k X^{k-1} + (\text{termes de degré inférieur})$$
Si $k \ge 1$, alors le polynôme $P(X+1) - P(X)$ a pour degré $k-1$ et son coefficient dominant est $a_k k$.
Pour que $P(X+1) - P(X) = 0$, il faut que tous ses coefficients soient nuls. En particulier, $a_k k$ doit être nul.
Puisque $a_k \neq 0$ (par hypothèse $P$ est de degré $k$), il faut que $k=0$.
Si $k=0$, le polynôme $P(X)$ est de degré 0, c'est-à-dire une constante $P(X) = a_0$.
Dans ce cas, $L(a_0) = a_0 - a_0 = 0$.
Donc, les seuls polynômes dans $E$ qui appartiennent à $\mathrm{Ker}(L)$ sont les polynômes constants.
L'ensemble des polynômes constants dans $E$ est $\mathbb{R}_0[X]$.
Une base pour $\mathbb{R}_0[X]$ est le polynôme constant $1$.
Par conséquent, $\mathrm{Ker}(L) = \mathrm{Vect}(1)$.
La dimension de $\mathrm{Ker}(L)$ est $\mathrm{dim}(\mathrm{Ker}(L)) = 1$.

\subsection*{Question 3 : Déterminer l'image $\mathrm{Im}(L)$.}

L'image $\mathrm{Im}(L)$ est l'ensemble des polynômes $Q \in E$ tels que $Q = L(P)$ pour un certain $P \in E$.
D'après la question 2, si $P(X)$ est un polynôme de degré $k \ge 1$, alors $L(P)(X)$ est un polynôme de degré $k-1$.
Si $P(X)$ est un polynôme constant (de degré 0), $L(P)(X) = 0$.
Puisque $P \in \mathbb{R}_n[X]$, le degré maximal de $P$ est $n$.
Si $P$ est de degré $n$, $L(P)$ est de degré $n-1$.
Si $P$ est de degré $k < n$, $L(P)$ est de degré $k-1$ (si $k \ge 1$) ou 0 (si $k=0$).
Donc, tous les polynômes dans $\mathrm{Im}(L)$ ont un degré inférieur ou égal à $n-1$.
Cela signifie que $\mathrm{Im}(L) \subseteq \mathbb{R}_{n-1}[X]$.

Nous connaissons la dimension de l'espace de départ $E = \mathbb{R}_n[X]$, qui est $\mathrm{dim}(E) = n+1$.
Nous avons trouvé $\mathrm{dim}(\mathrm{Ker}(L)) = 1$.
D'après le théorème du rang (voir question 4), $\mathrm{dim}(\mathrm{Im}(L)) = \mathrm{dim}(E) - \mathrm{dim}(\mathrm{Ker}(L)) = (n+1) - 1 = n$.

Nous avons donc $\mathrm{Im}(L) \subseteq \mathbb{R}_{n-1}[X]$ et $\mathrm{dim}(\mathrm{Im}(L)) = n$.
Puisque $\mathrm{dim}(\mathbb{R}_{n-1}[X]) = n$, et que $\mathrm{Im}(L)$ est un sous-espace de $\mathbb{R}_{n-1}[X]$ de même dimension, nous pouvons conclure que $\mathrm{Im}(L) = \mathbb{R}_{n-1}[X]$.

Une base pour $\mathrm{Im}(L) = \mathbb{R}_{n-1}[X]$ est la base canonique $(1, X, X^2, \dots, X^{n-1})$.
La dimension de $\mathrm{Im}(L)$ est $n$.

\subsection*{Question 4 : Calculer le rang de $L$ et vérifier le théorème du rang.}

Le rang d'une application linéaire $L$, noté $\mathrm{rang}(L)$, est la dimension de son image $\mathrm{Im}(L)$.
D'après la question 3, $\mathrm{dim}(\mathrm{Im}(L)) = n$.
Donc, $\mathrm{rang}(L) = n$.

Le théorème du rang stipule que pour une application linéaire $L: E \to F$, où $E$ est un espace vectoriel de dimension finie, on a :
$$\mathrm{dim}(E) = \mathrm{dim}(\mathrm{Ker}(L)) + \mathrm{dim}(\mathrm{Im}(L))$$
Dans notre cas :
*   $\mathrm{dim}(E) = \mathrm{dim}(\mathbb{R}_n[X]) = n+1$.
*   $\mathrm{dim}(\mathrm{Ker}(L)) = 1$ (d'après la question 2).
*   $\mathrm{dim}(\mathrm{Im}(L)) = n$ (d'après la question 3).

Vérifions le théorème du rang :
$$(n+1) = 1 + n$$
L'égalité est bien vérifiée.

\subsection*{Question 5 : Déterminer la matrice $M_B(L)$ de $L$ dans la base canonique $B$ de $E$.}

La base canonique de $E = \mathbb{R}_n[X]$ est $B = (P_0, P_1, \dots, P_n)$ où $P_k(X) = X^k$.
Pour construire la matrice $M_B(L)$, nous devons calculer $L(P_k)$ pour chaque $k \in \{0, \dots, n\}$ et exprimer le résultat comme une combinaison linéaire des éléments de la base $B$. Le vecteur colonne $j$ de la matrice sera constitué des coefficients de $L(P_j)$.

1.  Pour $P_0(X) = 1$:
    $L(1)(X) = (1) - (1) = 0$.
    Donc, $L(1) = 0 \cdot 1 + 0 \cdot X + \dots + 0 \cdot X^n$. La première colonne de $M_B(L)$ est $(0, 0, \dots, 0)^T$.

2.  Pour $P_k(X) = X^k$, avec $k \ge 1$:
    $L(X^k)(X) = (X+1)^k - X^k$.
    Utilisons la formule du binôme de Newton : $(X+1)^k = \sum_{j=0}^k \binom{k}{j} X^j = \binom{k}{0} X^0 + \binom{k}{1} X^1 + \dots + \binom{k}{k-1} X^{k-1} + \binom{k}{k} X^k$.
    $L(X^k)(X) = \left( \sum_{j=0}^k \binom{k}{j} X^j \right) - X^k$.
    Comme $\binom{k}{k}=1$, le terme $X^k$ s'annule :
    $L(X^k)(X) = \sum_{j=0}^{k-1} \binom{k}{j} X^j = \binom{k}{0} + \binom{k}{1} X + \binom{k}{2} X^2 + \dots + \binom{k}{k-1} X^{k-1}$.
    Ce polynôme est de degré $k-1$.

Les colonnes de la matrice $M_B(L)$ sont les vecteurs de coordonnées de $L(X^j)$ dans la base $(1, X, \dots, X^n)$.
*   Colonne 0 (pour $X^0=1$): $(0, 0, \dots, 0)^T$.
*   Colonne 1 (pour $X^1=X$): $L(X) = \binom{1}{0} = 1$. Coordonnées: $(1, 0, \dots, 0)^T$.
*   Colonne 2 (pour $X^2$): $L(X^2) = \binom{2}{0} + \binom{2}{1} X = 1 + 2X$. Coordonnées: $(1, 2, 0, \dots, 0)^T$.
*   Colonne 3 (pour $X^3$): $L(X^3) = \binom{3}{0} + \binom{3}{1} X + \binom{3}{2} X^2 = 1 + 3X + 3X^2$. Coordonnées: $(1, 3, 3, 0, \dots, 0)^T$.
*   Généralement, pour la colonne $j$ (correspondant à $X^j$): les coefficients sont $\binom{j}{0}, \binom{j}{1}, \dots, \binom{j}{j-1}$, suivis de zéros.

La matrice $M_B(L)$ est de taille $(n+1) \times (n+1)$. Ses coefficients $(M_B(L))_{i,j}$ (où $i$ est l'indice de ligne et $j$ l'indice de colonne, tous deux allant de 0 à $n$) sont donnés par :
$$(M_B(L))_{i,j} = \begin{cases} \binom{j}{i} & \text{si } i < j \\ 0 & \text{si } i \ge j \end{cases}$$
(Notez que $\binom{j}{i}=0$ si $i>j$ ou $i<0$, donc la condition $i<j$ est suffisante).

Exemple pour $n=3$, $E = \mathbb{R}_3[X]$, base $(1, X, X^2, X^3)$:
$L(1) = 0$
$L(X) = 1$
$L(X^2) = 1 + 2X$
$L(X^3) = 1 + 3X + 3X^2$

La matrice $M_B(L)$ est :
$$
M_B(L) = \begin{pmatrix}
(M_B(L))_{0,0} & (M_B(L))_{0,1} & (M_B(L))_{0,2} & (M_B(L))_{0,3} \\
(M_B(L))_{1,0} & (M_B(L))_{1,1} & (M_B(L))_{1,2} & (M_B(L))_{1,3} \\
(M_B(L))_{2,0} & (M_B(L))_{2,1} & (M_B(L))_{2,2} & (M_B(L))_{2,3} \\
(M_B(L))_{3,0} & (M_B(L))_{3,1} & (M_B(L))_{3,2} & (M_B(L))_{3,3}
\end{pmatrix}
=
\begin{pmatrix}
0 & \binom{1}{0} & \binom{2}{0} & \binom{3}{0} \\
0 & 0 & \binom{2}{1} & \binom{3}{1} \\
0 & 0 & 0 & \binom{3}{2} \\
0 & 0 & 0 & 0
\end{pmatrix}
=
\begin{pmatrix}
0 & 1 & 1 & 1 \\
0 & 0 & 2 & 3 \\
0 & 0 & 0 & 3 \\
0 & 0 & 0 & 0
\end{pmatrix}
$$
Cette matrice est strictement triangulaire supérieure. Son rang est $n$, comme attendu.

\newpage
\chapter*{Exercice 7 : Endomorphismes idempotents et stabilité (Difficulté : \textbf{})}

Soit $\mathbb{K}$ un corps commutatif.
Soit $E$ un espace vectoriel de dimension finie $n \geq 1$ sur $\mathbb{K}$.
Soit $u \in \mathcal{L}(E)$ un endomorphisme de $E$ tel que $u \circ u = u$, c'est-à-dire $u^2 = u$. Un tel endomorphisme est appelé un projecteur ou un endomorphisme idempotent.

1.  Démontrer que $E = \text{Ker}(u) \oplus \text{Im}(u)$.
    *(Indication : Considérer l'expression $x = (x - u(x)) + u(x)$ pour tout $x \in E$.)*
2.  Démontrer que $\text{Im}(u) = \{ y \in E \mid u(y) = y \}$.
    *(Autrement dit, $\text{Im}(u) = \text{Ker}(u - \text{Id}_E)$.)*
3.  Soit $v \in \mathcal{L}(E)$ un autre endomorphisme de $E$. On suppose que $u$ et $v$ commutent, c'est-à-dire $u \circ v = v \circ u$.
    Démontrer que $\text{Ker}(u)$ est stable par $v$, et que $\text{Im}(u)$ est stable par $v$.
    *(Un sous-espace $F$ de $E$ est stable par $v$ si $v(F) \subset F$.)*
4.  En déduire que si $u$ et $v$ commutent, alors la restriction de $v$ à $\text{Im}(u)$, notée $v|_{\text{Im}(u)}$, est un endomorphisme de $\text{Im}(u)$, et la restriction de $v$ à $\text{Ker}(u)$, notée $v|_{\text{Ker}(u)}$, est un endomorphisme de $\text{Ker}(u)$.

\section*{Correction détaillée}

1.  \textbf{Démontrer que $E = \text{Ker}(u) \oplus \text{Im}(u)$.}

    Pour établir qu'un espace vectoriel $E$ est la somme directe de deux sous-espaces vectoriels $F_1$ et $F_2$, il faut montrer que $F_1 \cap F_2 = \{0_E\}$ (l'intersection est réduite au vecteur nul) et que $E = F_1 + F_2$ (la somme est génératrice).

    a) \textbf{Montrons que $\text{Ker}(u) \cap \text{Im}(u) = \{0_E\}$.}
    Soit $y$ un vecteur appartenant à l'intersection $\text{Ker}(u) \cap \text{Im}(u)$.
    Puisque $y \in \text{Ker}(u)$, la définition du noyau implique que $u(y) = 0_E$.
    Puisque $y \in \text{Im}(u)$, la définition de l'image implique qu'il existe un vecteur $x \in E$ tel que $y = u(x)$.
    Appliquons l'endomorphisme $u$ à l'égalité $y = u(x)$ :
    $u(y) = u(u(x)) = u^2(x)$.
    Par hypothèse, nous savons que $u^2 = u$. Donc $u^2(x) = u(x)$.
    Ainsi, nous avons $u(y) = u(x)$.
    En substituant $u(x)$ par $y$, nous obtenons $u(y) = y$.
    Nous avons donc deux expressions pour $u(y)$ : $u(y) = 0_E$ (car $y \in \text{Ker}(u)$) et $u(y) = y$.
    En égalisant ces deux expressions, nous obtenons $y = 0_E$.
    Par conséquent, le seul vecteur commun à $\text{Ker}(u)$ et $\text{Im}(u)$ est le vecteur nul, ce qui prouve que $\text{Ker}(u) \cap \text{Im}(u) = \{0_E\}$.

    b) \textbf{Montrons que $E = \text{Ker}(u) + \text{Im}(u)$.}
    Soit $x$ un vecteur quelconque de $E$. Nous voulons montrer qu'il peut s'écrire comme la somme d'un vecteur de $\text{Ker}(u)$ et d'un vecteur de $\text{Im}(u)$.
    Considérons la décomposition suggérée : $x = (x - u(x)) + u(x)$.
    Analysons le premier terme, $x - u(x)$ :
    Appliquons $u$ à ce terme : $u(x - u(x)) = u(x) - u(u(x))$.
    Ceci peut s s'écrire $u(x) - u^2(x)$.
    Puisque $u^2 = u$ par hypothèse, nous avons $u^2(x) = u(x)$.
    Donc, $u(x - u(x)) = u(x) - u(x) = 0_E$.
    Par définition du noyau, cela signifie que le vecteur $x - u(x)$ appartient à $\text{Ker}(u)$.
    Analysons le second terme, $u(x)$ :
    Par définition même de l'image d'un endomorphisme, tout vecteur de la forme $u(z)$ pour un certain $z \in E$ appartient à $\text{Im}(u)$. Ici, $u(x)$ appartient clairement à $\text{Im}(u)$.
    Ainsi, tout vecteur $x \in E$ peut s'écrire comme la somme d'un vecteur de $\text{Ker}(u)$ (à savoir $x - u(x)$) et d'un vecteur de $\text{Im}(u)$ (à savoir $u(x)$).
    Ceci prouve que $E = \text{Ker}(u) + \text{Im}(u)$.

    En combinant les résultats de a) et b), nous concluons que $E = \text{Ker}(u) \oplus \text{Im}(u)$.

2.  \textbf{Démontrer que $\text{Im}(u) = \{ y \in E \mid u(y) = y \}$.}

    Notons $F = \{ y \in E \mid u(y) = y \}$. Nous devons montrer l'égalité des ensembles $\text{Im}(u)$ et $F$. Pour cela, nous démontrons la double inclusion.

    a) \textbf{Montrons que $\text{Im}(u) \subset F$.}
    Soit $y \in \text{Im}(u)$. Par définition, il existe un vecteur $x \in E$ tel que $y = u(x)$.
    Appliquons $u$ au vecteur $y$ : $u(y) = u(u(x)) = u^2(x)$.
    Puisque $u^2 = u$ par hypothèse, nous avons $u^2(x) = u(x)$.
    En substituant, $u(y) = u(x)$.
    Comme $y = u(x)$, on obtient $u(y) = y$.
    Par conséquent, $y$ satisfait la condition $u(y) = y$, ce qui signifie $y \in F$.
    Donc, tout élément de $\text{Im}(u)$ est un élément de $F$, d'où $\text{Im}(u) \subset F$.

    b) \textbf{Montrons que $F \subset \text{Im}(u)$.}
    Soit $y \in F$. Par définition de $F$, nous avons $u(y) = y$.
    L'expression $y = u(y)$ montre que $y$ est l'image d'un certain vecteur par $u$ (en l'occurrence, $y$ est l'image de lui-même par $u$).
    Par définition de l'image, tout vecteur de la forme $u(z)$ est dans $\text{Im}(u)$.
    Donc, $y \in \text{Im}(u)$.
    Ainsi, tout élément de $F$ est un élément de $\text{Im}(u)$, d'où $F \subset \text{Im}(u)$.

    Des inclusions a) et b), nous concluons que $\text{Im}(u) = \{ y \in E \mid u(y) = y \}$.
    Cette condition $u(y) = y$ est équivalente à $u(y) - y = 0_E$, ce qui peut s'écrire $(u - \text{Id}_E)(y) = 0_E$. C'est précisément la définition du noyau de l'endomorphisme $(u - \text{Id}_E)$, donc $\text{Im}(u) = \text{Ker}(u - \text{Id}_E)$.

3.  \textbf{Soit $v \in \mathcal{L}(E)$ tel que $u \circ v = v \circ u$. Démontrer que $\text{Ker}(u)$ est stable par $v$, et que $\text{Im}(u)$ est stable par $v$.}

    a) \textbf{Stabilité de $\text{Ker}(u)$ par $v$.}
    Un sous-espace $F$ est stable par $v$ si pour tout $x \in F$, $v(x) \in F$.
    Soit $x \in \text{Ker}(u)$. Par définition, $u(x) = 0_E$.
    Nous voulons montrer que $v(x) \in \text{Ker}(u)$, c'est-à-dire que $u(v(x)) = 0_E$.
    Calculons $u(v(x))$. Puisque $u$ et $v$ commutent (c'est-à-dire $u \circ v = v \circ u$), nous avons :
    $u(v(x)) = (u \circ v)(x) = (v \circ u)(x) = v(u(x))$.
    Comme $x \in \text{Ker}(u)$, nous savons que $u(x) = 0_E$.
    Donc, $u(v(x)) = v(0_E)$.
    Puisque $v$ est un endomorphisme linéaire, il transforme le vecteur nul en le vecteur nul : $v(0_E) = 0_E$.
    Ainsi, $u(v(x)) = 0_E$.
    Cela signifie que $v(x)$ appartient à $\text{Ker}(u)$.
    Par conséquent, $\text{Ker}(u)$ est stable par $v$.

    b) \textbf{Stabilité de $\text{Im}(u)$ par $v$.}
    Un sous-espace $F$ est stable par $v$ si pour tout $y \in F$, $v(y) \in F$.
    Soit $y \in \text{Im}(u)$. Par définition, il existe un vecteur $x \in E$ tel que $y = u(x)$.
    Nous voulons montrer que $v(y) \in \text{Im}(u)$.
    Calculons $v(y)$ : $v(y) = v(u(x))$.
    Puisque $u$ et $v$ commutent ($v \circ u = u \circ v$), nous pouvons écrire :
    $v(u(x)) = (v \circ u)(x) = (u \circ v)(x) = u(v(x))$.
    Donc, $v(y) = u(v(x))$.
    Par définition, tout vecteur qui est l'image d'un autre vecteur (ici $v(x)$) par $u$ appartient à $\text{Im}(u)$.
    Ainsi, $v(y) \in \text{Im}(u)$.
    Par conséquent, $\text{Im}(u)$ est stable par $v$.

4.  \textbf{En déduire que si $u$ et $v$ commutent, alors la restriction de $v$ à $\text{Im}(u)$, notée $v|_{\text{Im}(u)}$, est un endomorphisme de $\text{Im}(u)$, et la restriction de $v$ à $\text{Ker}(u)$, notée $v|_{\text{Ker}(u)}$, est un endomorphisme de $\text{Ker}(u)$.}

    Par définition, une application linéaire $w: F \to E'$ est un endomorphisme de $F$ si $E'=F$. C'est-à-dire que $w$ doit être une application linéaire de $F$ dans $F$. Pour qu'une restriction $v|_F : F \to E$ devienne un endomorphisme de $F$, il est nécessaire et suffisant que $F$ soit stable par $v$ (car la linéarité de $v$ implique la linéarité de sa restriction).

    a) \textbf{Pour $F = \text{Ker}(u)$ :}
    D'après la question 3a), si $u$ et $v$ commutent, alors $\text{Ker}(u)$ est stable par $v$.
    Cela signifie que pour tout $x \in \text{Ker}(u)$, $v(x) \in \text{Ker}(u)$.
    Nous pouvons donc définir une application $v|_{\text{Ker}(u)} : \text{Ker}(u) \to \text{Ker}(u)$ qui associe à $x \in \text{Ker}(u)$ le vecteur $v(x)$.
    Puisque $v$ est un endomorphisme linéaire de $E$, sa restriction à un sous-espace, $v|_{\text{Ker}(u)}$, est également linéaire.
    Par conséquent, $v|_{\text{Ker}(u)}$ est un endomorphisme de $\text{Ker}(u)$.

    b) \textbf{Pour $F = \text{Im}(u)$ :}
    D'après la question 3b), si $u$ et $v$ commutent, alors $\text{Im}(u)$ est stable par $v$.
    Cela signifie que pour tout $y \in \text{Im}(u)$, $v(y) \in \text{Im}(u)$.
    Nous pouvons donc définir une application $v|_{\text{Im}(u)} : \text{Im}(u) \to \text{Im}(u)$ qui associe à $y \in \text{Im}(u)$ le vecteur $v(y)$.
    Puisque $v$ est un endomorphisme linéaire de $E$, sa restriction à un sous-espace, $v|_{\text{Im}(u)}$, est également linéaire.
    Par conséquent, $v|_{\text{Im}(u)}$ est un endomorphisme de $\text{Im}(u)$.

\newpage
\chapter*{Exercice 8 : Indice de Fitting et Décomposition de Fitting (Difficulté : \textbf{})}

Soit $\mathbb{K}$ un corps commutatif.
Soit $E$ un espace vectoriel sur $\mathbb{K}$ de dimension finie $n \ge 1$.
Soit $f \in \mathcal{L}(E)$ un endomorphisme de $E$.
On définit l'application $f^0 = \text{Id}_E$ (l'application identité sur $E$), et pour tout entier $k \in \mathbb{N}$, $f^{k+1} = f \circ f^k$.



\section*{Correction détaillée}

1.  \textbf{Analyse des Noyaux}

    a.  Pour montrer que $\ker(f^k) \subseteq \ker(f^{k+1})$ pour tout $k \in \mathbb{N}$ :
        Soit $x \in \ker(f^k)$. Par définition du noyau d'une application linéaire, cela signifie que $f^k(x) = 0_E$.
        Appliquons l'endomorphisme $f$ à cette égalité. Puisque $f$ est linéaire, $f(0_E) = 0_E$.
        Nous obtenons $f(f^k(x)) = f(0_E)$, ce qui se simplifie en $f^{k+1}(x) = 0_E$ par définition de la composition des applications.
        L'égalité $f^{k+1}(x) = 0_E$ signifie que $x$ appartient au noyau de $f^{k+1}$, c'est-à-dire $x \in \ker(f^{k+1})$.
        Par conséquent, tout élément de $\ker(f^k)$ est aussi un élément de $\ker(f^{k+1})$, d'où l'inclusion $\ker(f^k) \subseteq \ker(f^{k+1})$ pour tout $k \in \mathbb{N}$.

    b.  La suite $(\ker(f^k))_{k \in \mathbb{N}}$ est une suite croissante de sous-espaces vectoriels de $E$.
        Puisque $E$ est un espace vectoriel de dimension finie $n$, la dimension de chaque sous-espace vectoriel $\ker(f^k)$ est un entier naturel. Cette dimension est bornée supérieurement par la dimension de $E$, c'est-à-dire $n$.
        Nous avons la suite d'inégalités de dimensions :
        $0 \le \dim(\ker(f^0)) \le \dim(\ker(f^1)) \le \dim(\ker(f^2)) \le \dots \le \dim(E) = n$.
        Une suite croissante d'entiers naturels qui est bornée supérieurement converge et donc doit nécessairement stabiliser à partir d'un certain rang.
        Si la dimension $\dim(\ker(f^k))$ est égale à $\dim(\ker(f^{k+1}))$, et comme nous avons déjà établi l'inclusion $\ker(f^k) \subseteq \ker(f^{k+1})$, l'égalité des dimensions implique que les deux sous-espaces vectoriels sont identiques : $\ker(f^k) = \ker(f^{k+1})$.
        Ainsi, la suite des noyaux doit stabiliser. Par le principe du bon ordre des entiers naturels, il existe un plus petit entier naturel $d_0$ tel que $\ker(f^{d_0}) = \ker(f^{d_0+1})$.

    c.  Nous montrons par récurrence sur $k$ que pour tout $k \ge d_0$, on a $\ker(f^k) = \ker(f^{d_0})$.
        *   \textbf{Initialisation (k = d\_0) :} La propriété est vraie par définition de $d_0$, puisque $\ker(f^{d_0}) = \ker(f^{d_0+1})$.
        *   \textbf{Hypothèse de récurrence :} Supposons que pour un certain entier $k \ge d_0$, on ait $\ker(f^k) = \ker(f^{d_0})$.
        *   \textbf{Hérédité :} Démontrons que $\ker(f^{k+1}) = \ker(f^{d_0})$.
            Nous savons d'après 1.a que $\ker(f^k) \subseteq \ker(f^{k+1})$. Par transitivité, $\ker(f^{d_0}) = \ker(f^k) \subseteq \ker(f^{k+1})$.
            Il reste à prouver l'inclusion inverse : $\ker(f^{k+1}) \subseteq \ker(f^{d_0})$.
            Soit $x \in \ker(f^{k+1})$. Par définition, $f^{k+1}(x) = 0_E$.
            On peut écrire $f^{k+1}(x) = f^k(f(x))$. Donc $f^k(f(x)) = 0_E$.
            Ceci signifie que $f(x) \in \ker(f^k)$.
            D'après l'hypothèse de récurrence, $\ker(f^k) = \ker(f^{d_0})$.
            Par conséquent, $f(x) \in \ker(f^{d_0})$.
            Cela implique que $f^{d_0}(f(x)) = 0_E$, ce qui est équivalent à $f^{d_0+1}(x) = 0_E$.
            Puisque, par définition de $d_0$, nous avons $\ker(f^{d_0+1}) = \ker(f^{d_0})$, l'égalité $f^{d_0+1}(x) = 0_E$ signifie que $x \in \ker(f^{d_0})$.
            Ainsi, $\ker(f^{k+1}) \subseteq \ker(f^{d_0})$.
            En combinant les inclusions, nous avons $\ker(f^{k+1}) = \ker(f^{d_0})$.
        Par le principe de récurrence, pour tout $k \ge d_0$, l'égalité $\ker(f^k) = \ker(f^{d_0})$ est établie.

2.  \textbf{Analyse des Images}

    a.  Pour montrer que $\text{Im}(f^{k+1}) \subseteq \text{Im}(f^k)$ pour tout $k \in \mathbb{N}$ :
        Soit $y \in \text{Im}(f^{k+1})$. Par définition de l'image, il existe un vecteur $x \in E$ tel que $y = f^{k+1}(x)$.
        On peut réécrire $f^{k+1}(x)$ comme $f^k(f(x))$.
        Posons $x' = f(x)$. Alors $x'$ est un vecteur de $E$.
        L'égalité devient $y = f^k(x')$.
        Ceci signifie que $y$ appartient à l'image de $f^k$, c'est-à-dire $y \in \text{Im}(f^k)$.
        Par conséquent, $\text{Im}(f^{k+1}) \subseteq \text{Im}(f^k)$ pour tout $k \in \mathbb{N}$.

    b.  La suite $(\text{Im}(f^k))_{k \in \mathbb{N}}$ est une suite décroissante de sous-espaces vectoriels de $E$.
        Puisque $E$ est un espace vectoriel de dimension finie $n$, la dimension de chaque sous-espace vectoriel $\text{Im}(f^k)$ est un entier naturel, bornée inférieurement par 0.
        Nous avons la suite d'inégalités de dimensions :
        $n = \dim(E) \ge \dim(\text{Im}(f^0)) \ge \dim(\text{Im}(f^1)) \ge \dim(\text{Im}(f^2)) \ge \dots \ge 0$.
        Une suite décroissante d'entiers naturels qui est bornée inférieurement doit nécessairement stabiliser à partir d'un certain rang.
        Si la dimension $\dim(\text{Im}(f^k))$ est égale à $\dim(\text{Im}(f^{k+1}))$, et comme nous avons déjà établi l'inclusion $\text{Im}(f^{k+1}) \subseteq \text{Im}(f^k)$, l'égalité des dimensions implique que les deux sous-espaces vectoriels sont identiques : $\text{Im}(f^k) = \text{Im}(f^{k+1})$.
        Ainsi, la suite des images doit stabiliser. Par le principe du bon ordre des entiers naturels, il existe un plus petit entier naturel $p_0$ tel que $\text{Im}(f^{p_0}) = \text{Im}(f^{p_0+1})$.

    c.  Nous devons montrer que pour tout $k \ge p_0$, on a $\text{Im}(f^k) = \text{Im}(f^{p_0})$.
        Par définition de $p_0$, nous avons $\text{Im}(f^{p_0}) = \text{Im}(f^{p_0+1})$.
        Considérons la restriction de $f$ à $\text{Im}(f^{p_0})$. Notons cette application $g = f|_{\text{Im}(f^{p_0})} : \text{Im}(f^{p_0}) \to E$.
        L'image de cette application est $g(\text{Im}(f^{p_0})) = f(\text{Im}(f^{p_0})) = f(f^{p_0}(E)) = f^{p_0+1}(E) = \text{Im}(f^{p_0+1})$.
        Puisque $\text{Im}(f^{p_0}) = \text{Im}(f^{p_0+1})$, cela signifie que l'image de $g$ est $\text{Im}(f^{p_0})$.
        Ainsi, $g: \text{Im}(f^{p_0}) \to \text{Im}(f^{p_0})$ est un endomorphisme de l'espace $\text{Im}(f^{p_0})$.
        De plus, puisque l'image de $g$ est $\text{Im}(f^{p_0})$ lui-même, cet endomorphisme $g$ est surjectif.
        Comme $\text{Im}(f^{p_0})$ est un sous-espace vectoriel de $E$, il est de dimension finie. Pour un espace vectoriel de dimension finie, un endomorphisme surjectif est aussi injectif, et donc un automorphisme.
        Par conséquent, $f|_{\text{Im}(f^{p_0})} : \text{Im}(f^{p_0}) \to \text{Im}(f^{p_0})$ est un automorphisme.
        Cela signifie que pour tout entier $k \ge 1$, la $k$-ième composition de $f|_{\text{Im}(f^{p_0})}$ sur lui-même, notée $(f|_{\text{Im}(f^{p_0})})^k$, est également un automorphisme de $\text{Im}(f^{p_0})$.
        En particulier, $(f|_{\text{Im}(f^{p_0})})^k$ est surjective, ce qui implique que $f^k(\text{Im}(f^{p_0})) = \text{Im}(f^{p_0})$ pour tout $k \ge 1$.
        Par définition, $f^k(\text{Im}(f^{p_0})) = f^k(f^{p_0}(E)) = f^{k+p_0}(E) = \text{Im}(f^{k+p_0})$.
        En combinant ces résultats, nous obtenons $\text{Im}(f^{p_0+k}) = \text{Im}(f^{p_0})$ pour tout $k \ge 1$.
        Ceci démontre que pour tout $k \ge p_0$, on a $\text{Im}(f^k) = \text{Im}(f^{p_0})$.

3.  \textbf{Lien entre Noyau et Image (Théorème du Rang)}

    a.  Le théorème du rang affirme que pour toute application linéaire $L: V \to W$ où $V$ est de dimension finie, on a $\dim(V) = \dim(\ker(L)) + \dim(\text{Im}(L))$.
        Appliquons ce théorème à l'endomorphisme $f^k: E \to E$ pour tout $k \in \mathbb{N}$. L'espace $E$ est de dimension finie $n$.
        Donc, $n = \dim(\ker(f^k)) + \dim(\text{Im}(f^k))$.

        *   Montrons que $p_0 \le d_0$ :
            Par définition de $d_0$, nous avons $\ker(f^{d_0}) = \ker(f^{d_0+1})$.
            Ceci implique que $\dim(\ker(f^{d_0})) = \dim(\ker(f^{d_0+1}))$.
            En utilisant le théorème du rang pour $f^{d_0}$ et $f^{d_0+1}$ :
            $\dim(\text{Im}(f^{d_0})) = n - \dim(\ker(f^{d_0}))$
            $\dim(\text{Im}(f^{d_0+1})) = n - \dim(\ker(f^{d_0+1}))$
            Puisque $\dim(\ker(f^{d_0})) = \dim(\ker(f^{d_0+1}))$, il s'ensuit que $\dim(\text{Im}(f^{d_0})) = \dim(\text{Im}(f^{d_0+1}))$.
            Comme nous savons que $\text{Im}(f^{d_0+1}) \subseteq \text{Im}(f^{d_0})$ (d'après 2.a), l'égalité des dimensions implique que $\text{Im}(f^{d_0}) = \text{Im}(f^{d_0+1})$.
            Par définition, $p_0$ est le plus petit entier tel que $\text{Im}(f^{p_0}) = \text{Im}(f^{p_0+1})$. Puisque cette égalité est vérifiée pour $d_0$, nous devons avoir $p_0 \le d_0$.

        *   Montrons que $d_0 \le p_0$ :
            Par définition de $p_0$, nous avons $\text{Im}(f^{p_0}) = \text{Im}(f^{p_0+1})$.
            Ceci implique que $\dim(\text{Im}(f^{p_0})) = \dim(\text{Im}(f^{p_0+1}))$.
            En utilisant le théorème du rang pour $f^{p_0}$ et $f^{p_0+1}$ :
            $\dim(\ker(f^{p_0})) = n - \dim(\text{Im}(f^{p_0}))$
            $\dim(\ker(f^{p_0+1})) = n - \dim(\text{Im}(f^{p_0+1}))$
            Puisque $\dim(\text{Im}(f^{p_0})) = \dim(\text{Im}(f^{p_0+1}))$, il s'ensuit que $\dim(\ker(f^{p_0})) = \dim(\ker(f^{p_0+1}))$.
            Comme nous savons que $\ker(f^{p_0}) \subseteq \ker(f^{p_0+1})$ (d'après 1.a), l'égalité des dimensions implique que $\ker(f^{p_0}) = \ker(f^{p_0+1})$.
            Par définition, $d_0$ est le plus petit entier tel que $\ker(f^{d_0}) = \ker(f^{d_0+1})$. Puisque cette égalité est vérifiée pour $p_0$, nous devons avoir $d_0 \le p_0$.

        Des deux inégalités $p_0 \le d_0$ et $d_0 \le p_0$, nous concluons que $d_0 = p_0$. Nous noterons cette valeur commune $k_0$.

    b.  D'après la question 2.c, la condition $\text{Im}(f^{p_0}) = \text{Im}(f^{p_0+1})$ implique que la restriction de $f$ à $\text{Im}(f^{p_0})$, soit $f|_{\text{Im}(f^{p_0})}$, est un automorphisme de $\text{Im}(f^{p_0})$.
        En utilisant $k_0 = p_0$, cela signifie que $f|_{\text{Im}(f^{k_0})} : \text{Im}(f^{k_0}) \to \text{Im}(f^{k_0})$ est un endomorphisme dont l'image est $\text{Im}(f^{k_0+1})$, qui est égale à $\text{Im}(f^{k_0})$.
        Par conséquent, $f|_{\text{Im}(f^{k_0})}$ est un endomorphisme surjectif de $\text{Im}(f^{k_0})$.
        Puisque $\text{Im}(f^{k_0})$ est un sous-espace vectoriel de $E$, il est de dimension finie.
        Dans un espace vectoriel de dimension finie, un endomorphisme surjectif est nécessairement bijectif (et donc injectif).
        Ainsi, $f|_{\text{Im}(f^{k_0})}$ est un isomorphisme (un automorphisme) de $\text{Im}(f^{k_0})$ sur lui-même.

4.  \textbf{Décomposition de Fitting}

    a.  Nous devons montrer que $\ker(f^{k_0}) \cap \text{Im}(f^{k_0}) = \{0_E\}$.
        Soit $x \in \ker(f^{k_0}) \cap \text{Im}(f^{k_0})$.
        Puisque $x \in \ker(f^{k_0})$, par définition, $f^{k_0}(x) = 0_E$.
        Puisque $x \in \text{Im}(f^{k_0})$, par définition, il existe un vecteur $y \in E$ tel que $x = f^{k_0}(y)$.
        Substituons cette expression de $x$ dans la première égalité :
        $f^{k_0}(f^{k_0}(y)) = 0_E$.
        Ceci s'écrit $f^{2k_0}(y) = 0_E$.
        Cela implique que $y \in \ker(f^{2k_0})$.
        D'après la question 1.c, nous savons que pour tout $k \ge d_0$, $\ker(f^k) = \ker(f^{d_0})$. Puisque $k_0 = d_0$, cette propriété s'applique à $k_0$.
        En particulier, puisque $2k_0 \ge k_0$ (tout en supposant $k_0 \ge 0$; si $k_0=0$, $\ker(f^0)=\{0_E\}$, alors $\ker(f^{2\cdot 0})=\ker(f^0)$), nous avons $\ker(f^{2k_0}) = \ker(f^{k_0})$.
        Par conséquent, $y \in \ker(f^{k_0})$.
        Cela signifie que $f^{k_0}(y) = 0_E$.
        Puisque nous avions posé $x = f^{k_0}(y)$, nous en déduisons que $x = 0_E$.
        Ainsi, le seul élément de l'intersection de $\ker(f^{k_0})$ et $\text{Im}(f^{k_0})$ est le vecteur nul, donc $\ker(f^{k_0}) \cap \text{Im}(f^{k_0}) = \{0_E\}$.

    b.  Pour prouver que $E = \ker(f^{k_0}) \oplus \text{Im}(f^{k_0})$, il suffit de montrer deux conditions pour les sous-espaces vectoriels $U = \ker(f^{k_0})$ et $V = \text{Im}(f^{k_0})$ :
        i.  Leur intersection est immédiate d'après la définition : $U \cap V = \{0_E\}$. Cette condition a été démontrée à la question 4.a.
        ii. La somme de leurs dimensions est égale à la dimension de l'espace ambiant : $\dim(U) + \dim(V) = \dim(E)$.
            Appliquons le théorème du rang à l'endomorphisme $f^{k_0} : E \to E$:
            $\dim(E) = \dim(\ker(f^{k_0})) + \dim(\text{Im}(f^{k_0}))$.
            Cette condition est également vérifiée.

        Étant donné que ces deux conditions sont satisfaites, nous pouvons conclure que l'espace vectoriel $E$ est la somme directe de $\ker(f^{k_0})$ et $\text{Im}(f^{k_0})$.
        Ainsi, $E = \ker(f^{k_0}) \oplus \text{Im}(f^{k_0})$.

\newpage
\chapter*{Exercice 9 : Décomposition de Fitting et endomorphismes nilpotents/inversibles (Difficulté : \textbf{}*)}

Soit $K$ un corps commutatif quelconque. Soit $E$ un $K$-espace vectoriel de dimension finie $n \ge 1$.
Soit $u \in L(E)$ un endomorphisme de $E$.

Pour tout entier naturel $k \ge 0$, on définit les sous-espaces vectoriels :
- $K_k = \ker(u^k)$, le noyau de $u^k$.
- $I_k = \mathrm{Im}(u^k)$, l'image de $u^k$.
Par convention, $u^0 = \mathrm{id}_E$, donc $K_0 = \{0\}$ et $I_0 = E$.

1.  Démontrer que la suite $(K_k)_{k \ge 0}$ est une suite croissante de sous-espaces vectoriels de $E$ (pour l'inclusion), et que la suite $(I_k)_{k \ge 0}$ est une suite décroissante de sous-espaces vectoriels de $E$.
2.  Démontrer qu'il existe un plus petit entier $p \in \{0, \dots, n\}$ tel que $K_p = K_{p+1}$. Démontrer alors que pour tout $k \ge p$, on a $K_k = K_p$. Démontrer un résultat analogue pour la suite $(I_k)_{k \ge 0}$, c'est-à-dire $I_p = I_{p+1}$ et $I_k = I_p$ pour tout $k \ge p$. Cet entier $p$ est appelé l'indice de Fitting de l'endomorphisme $u$.
3.  On se place avec l'entier $p$ défini à la question précédente. Démontrer que $E = K_p \oplus I_p$.
4.  Démontrer que les sous-espaces $K_p$ et $I_p$ sont stables par $u$. On note $u_K$ la restriction de $u$ à $K_p$ (vue comme endomorphisme de $K_p$) et $u_I$ la restriction de $u$ à $I_p$ (vue comme endomorphisme de $I_p$). Démontrer que $u_K$ est un endomorphisme nilpotent et que $u_I$ est un automorphisme.

---

\section*{Correction détaillée}

1.  \textbf{Monotonie des suites $(K_k)_{k \ge 0}$ et $(I_k)_{k \ge 0}$}

    *   \textbf{Pour $(K_k)_{k \ge 0}$ :}
        Soit $k \ge 0$. Nous voulons montrer que $K_k \subseteq K_{k+1}$.
        Soit $x \in K_k$. Par définition, $u^k(x) = 0$.
        Alors $u^{k+1}(x) = u(u^k(x)) = u(0) = 0$.
        Donc $x \in K_{k+1}$.
        Ceci prouve que $K_k \subseteq K_{k+1}$ pour tout $k \ge 0$. La suite $(K_k)_{k \ge 0}$ est donc croissante pour l'inclusion.

    *   \textbf{Pour $(I_k)_{k \ge 0}$ :}
        Soit $k \ge 0$. Nous voulons montrer que $I_{k+1} \subseteq I_k$.
        Soit $y \in I_{k+1}$. Par définition, il existe $x \in E$ tel que $y = u^{k+1}(x)$.
        On peut écrire $y = u^k(u(x))$.
        Puisque $u(x) \in E$, $y$ est l'image d'un élément de $E$ par $u^k$.
        Donc $y \in I_k$.
        Ceci prouve que $I_{k+1} \subseteq I_k$ pour tout $k \ge 0$. La suite $(I_k)_{k \ge 0}$ est donc décroissante pour l'inclusion.

2.  \textbf{Stabilisation des suites $(K_k)_{k \ge 0}$ et $(I_k)_{k \ge 0}$}

    *   \textbf{Pour $(K_k)_{k \ge 0}$ :}
        La suite $(K_k)_{k \ge 0}$ est une suite croissante de sous-espaces vectoriels de $E$.
        De plus, $E$ est de dimension finie $n$.
        Cela implique que la suite des dimensions $\mathrm{dim}(K_k)$ est une suite croissante d'entiers naturels, majorée par $n$.
        Une telle suite doit nécessairement stationner. C'est-à-dire qu'il existe un entier $p_0$ tel que $\mathrm{dim}(K_{p_0}) = \mathrm{dim}(K_{p_0+1})$.
        Puisque $K_{p_0} \subseteq K_{p_0+1}$ et qu'ils ont la même dimension, on en déduit que $K_{p_0} = K_{p_0+1}$.
        Soit $p$ le plus petit entier tel que $K_p = K_{p+1}$. Un tel $p$ existe car la suite stationne et $K_0 \subseteq K_1 \subseteq \dots \subseteq K_n = E$ est la chaîne maximale de sous-espaces strictement croissants. Si $K_i \subsetneq K_{i+1}$ pour tout $i$, alors on aurait $\mathrm{dim}(K_{i+1}) \ge \mathrm{dim}(K_i) + 1$, ce qui impliquerait $\mathrm{dim}(K_n) \ge \mathrm{dim}(K_0) + n = n$, atteignant la dimension de $E$ au plus en $n$ étapes.
        Nous allons montrer par récurrence que pour tout $k \ge p$, $K_k = K_p$.
        L'initialisation ($k=p$) est immédiate d'après la définition : $K_p = K_p$.
        L'étape de récurrence : Supposons que $K_k = K_p$ pour un $k \ge p$. Nous voulons montrer que $K_{k+1} = K_p$.
        Nous savons déjà que $K_p \subseteq K_{k+1}$ (car la suite est croissante).
        Il reste à montrer que $K_{k+1} \subseteq K_p$. Soit $x \in K_{k+1}$. Alors $u^{k+1}(x) = 0$.
        Nous avons $u^p(x) \in E$. Alors $u^{k+1-p}(u^p(x)) = u^{k+1}(x) = 0$.
        Puisque $k \ge p$, $k+1-p \ge 1$.
        Considérons l'application $u$ restreinte à $K_{j+1}$ pour un certain $j$.
        Si $y \in K_{k+1}$, alors $u^{k+1}(y)=0$.
        Si $K_p=K_{p+1}$, montrons $K_{p+1}=K_{p+2}$.
        Soit $x \in K_{p+2}$. Alors $u^{p+2}(x)=0$. On peut écrire $u^{p+1}(u(x))=0$. Donc $u(x) \in K_{p+1}$.
        Puisque $K_{p+1} = K_p$, on a $u(x) \in K_p$, ce qui signifie $u^p(u(x)) = 0$, donc $u^{p+1}(x) = 0$.
        Par conséquent $x \in K_{p+1}$. Comme $K_{p+1} = K_p$, on a $x \in K_p$.
        Donc $K_{p+2} \subseteq K_p$. Puisque $K_p \subseteq K_{p+2}$ (par croissance), on a $K_{p+2}=K_p$.
        En réitérant ce raisonnement, on montre par récurrence que $K_k=K_p$ pour tout $k \ge p$.

    *   \textbf{Pour $(I_k)_{k \ge 0}$ :}
        La suite $(I_k)_{k \ge 0}$ est une suite décroissante de sous-espaces vectoriels de $E$.
        Comme précédemment, la suite des dimensions $\mathrm{dim}(I_k)$ est une suite décroissante d'entiers naturels, minorée par $0$.
        Elle doit donc stationner. Il existe un entier $p'$ tel que $\mathrm{dim}(I_{p'}) = \mathrm{dim}(I_{p'+1})$.
        Puisque $I_{p'+1} \subseteq I_{p'}$ et qu'ils ont la même dimension, on en déduit que $I_{p'} = I_{p'+1}$.
        Le théorème du rang stipule que $\mathrm{dim}(E) = \mathrm{dim}(\ker(u^k)) + \mathrm{dim}(\mathrm{Im}(u^k))$, soit $n = \mathrm{dim}(K_k) + \mathrm{dim}(I_k)$.
        Si $K_p = K_{p+1}$, alors $\mathrm{dim}(K_p) = \mathrm{dim}(K_{p+1})$.
        Donc $n - \mathrm{dim}(I_p) = n - \mathrm{dim}(I_{p+1})$, ce qui implique $\mathrm{dim}(I_p) = \mathrm{dim}(I_{p+1})$.
        Puisque $I_{p+1} \subseteq I_p$ et qu'ils ont la même dimension, on a $I_p = I_{p+1}$.
        Ainsi, l'entier $p$ est le même pour la stabilisation de $(K_k)$ et $(I_k)$.
        Pour montrer $I_k = I_p$ pour tout $k \ge p$:
        Nous savons déjà que $I_k \subseteq I_p$ (car la suite est décroissante).
        Il reste à montrer $I_p \subseteq I_k$.
        Puisque $I_p = I_{p+1}$, cela signifie que pour tout $y \in I_p$, $y \in I_{p+1}$, donc $y = u^{p+1}(x')$ pour un certain $x' \in E$.
        Plus généralement, si $I_j = I_{j+1}$, alors l'application $u$ restreinte à $I_j$ est surjective de $I_j$ dans $I_j$. En effet, $u(I_j) = u(\mathrm{Im}(u^j)) = \mathrm{Im}(u^{j+1}) = I_{j+1} = I_j$. Étant surjective sur un espace de dimension finie, elle est un automorphisme.
        Donc $u|_{I_p} : I_p \to I_p$ est un automorphisme.
        Ceci signifie que $u(I_p) = I_p$. Par conséquent $u^2(I_p) = u(I_p) = I_p$, et par récurrence $u^k(I_p) = I_p$ pour tout $k \ge 0$.
        En particulier, $u^k(E) = u^k(I_0) \supseteq u^k(I_p)$.
        Soit $y \in I_p$. Alors $y = u^p(x)$ pour un certain $x \in E$.
        Puisque $I_p = I_k$ pour $k \ge p$, et on sait que $I_k \subseteq I_p$.
        Si $y \in I_p$, alors $y = u^p(x)$. Mais on a montré que $u|_{I_p}$ est un automorphisme.
        Donc $u^k|_{I_p}$ est aussi un automorphisme pour tout $k \ge 0$.
        Conséquence : $\mathrm{Im}(u^k|_{I_p}) = I_p$.
        Mais $\mathrm{Im}(u^k|_{I_p})$ est l'ensemble des éléments de la forme $u^k(z)$ où $z \in I_p$.
        C'est donc un sous-ensemble de $I_k$.
        Non, c'est $u^k(I_p)$. Et nous avons montré que $u^k(I_p) = I_p$.
        Donc $I_p \subseteq I_k$ pour $k \ge p$.
        Puisque nous avons aussi $I_k \subseteq I_p$, il s'ensuit que $I_k = I_p$ pour tout $k \ge p$.

3.  \textbf{Décomposition $E = K_p \oplus I_p$}

    Pour montrer que $E = K_p \oplus I_p$, nous devons montrer deux choses :
    a. $K_p \cap I_p = \{0\}$
    b. $K_p + I_p = E$ (ou de manière équivalente $\mathrm{dim}(K_p) + \mathrm{dim}(I_p) = \mathrm{dim}(E)$ si $K_p \cap I_p = \{0\}$)

    *   \textbf{a. $K_p \cap I_p = \{0\}$ :}
        Soit $y \in K_p \cap I_p$.
        Puisque $y \in K_p$, on a $u^p(y) = 0$.
        Puisque $y \in I_p$, il existe un $x \in E$ tel que $y = u^p(x)$.
        En substituant, nous obtenons $u^p(u^p(x)) = 0$, c'est-à-dire $u^{2p}(x) = 0$.
        Cela signifie que $x \in K_{2p}$.
        Or, d'après la question 2, puisque $p$ est l'indice de stabilisation, nous avons $K_{2p} = K_p$ (car $2p \ge p$).
        Donc $x \in K_p$.
        Par conséquent, $u^p(x) = 0$.
        Puisque $y = u^p(x)$, on en déduit $y = 0$.
        Donc $K_p \cap I_p = \{0\}$.

    *   \textbf{b. $K_p + I_p = E$ :}
        D'après le théorème du rang appliqué à $u^p$:
        $\mathrm{dim}(E) = \mathrm{dim}(\ker(u^p)) + \mathrm{dim}(\mathrm{Im}(u^p))$
        $\mathrm{dim}(E) = \mathrm{dim}(K_p) + \mathrm{dim}(I_p)$.
        Puisque $K_p \cap I_p = \{0\}$, la somme $K_p + I_p$ est directe, et sa dimension est $\mathrm{dim}(K_p \oplus I_p) = \mathrm{dim}(K_p) + \mathrm{dim}(I_p)$.
        Donc $\mathrm{dim}(K_p \oplus I_p) = \mathrm{dim}(E)$.
        Puisque $K_p \oplus I_p$ est un sous-espace vectoriel de $E$ et qu'il a la même dimension que $E$, il est égal à $E$.
        Ainsi, $E = K_p \oplus I_p$.

4.  \textbf{Stabilité par $u$, nilpotence et automorphisme}

    *   \textbf{Stabilité des sous-espaces $K_p$ et $I_p$ par $u$ :}
        *   \textbf{Pour $K_p$ :}
            Soit $x \in K_p$. Alors $u^p(x) = 0$.
            Nous devons montrer que $u(x) \in K_p$.
            Calculons $u^p(u(x))$. C'est $u^{p+1}(x)$.
            D'après la question 2, $K_p = K_{p+1}$ (par définition de $p$).
            Puisque $x \in K_p$, $u^p(x)=0$. On peut dire que $x \in K_{p+1}$ également.
            Or $K_p = K_{p+1}$ signifie que si $u^{p+1}(x) = 0$, alors $u^p(x) = 0$.
            Ah, wait. $u(x) \in K_p$ requires $u^p(u(x))=0$.
            We know $u^p(x)=0$. So $x \in K_p$.
            We want to check if $u(x) \in K_p$. This means $u^p(u(x))=0$, i.e. $u^{p+1}(x)=0$.
            Since $x \in K_p$ and $K_p = K_{p+1}$, it implies that $x \in K_{p+1}$, so $u^{p+1}(x)=0$ holds.
            Therefore, $u(x) \in K_p$.
            Ceci prouve que $K_p$ est stable par $u$.

        *   \textbf{Pour $I_p$ :}
            Soit $y \in I_p$. Alors il existe $x \in E$ tel que $y = u^p(x)$.
            Nous devons montrer que $u(y) \in I_p$.
            Calculons $u(y) = u(u^p(x)) = u^{p+1}(x)$.
            Par définition, $u^{p+1}(x) \in \mathrm{Im}(u^{p+1}) = I_{p+1}$.
            D'après la question 2, $I_p = I_{p+1}$ (car $p$ est l'indice de stabilisation).
            Donc $u(y) \in I_p$.
            Ceci prouve que $I_p$ est stable par $u$.

    *   \textbf{Nature des endomorphismes $u_K$ et $u_I$ :}
        *   \textbf{Pour $u_K : K_p \to K_p$ :}
            Par définition, $u_K(x) = u(x)$ pour $x \in K_p$.
            Pour tout $x \in K_p$, nous avons $u^p(x) = 0$.
            Donc $u_K^p(x) = u^p(x) = 0$.
            Ceci signifie que $u_K^p$ est l'endomorphisme nul sur $K_p$.
            Par conséquent, $u_K$ est un endomorphisme nilpotent.

        *   \textbf{Pour $u_I : I_p \to I_p$ :}
            Par définition, $u_I(y) = u(y)$ pour $y \in I_p$.
            Nous savons que $u_I$ est un endomorphisme de $I_p$.
            Pour montrer que $u_I$ est un automorphisme, il suffit de montrer qu'il est injectif ou surjectif, car $I_p$ est de dimension finie.
            Montrons que $u_I$ est surjectif.
            L'image de $u_I$ est $\mathrm{Im}(u_I) = \{u(y) \mid y \in I_p\}$.
            Puisque $y \in I_p = \mathrm{Im}(u^p)$, il existe $x \in E$ tel que $y = u^p(x)$.
            Alors $u(y) = u(u^p(x)) = u^{p+1}(x)$.
            Ainsi, $\mathrm{Im}(u_I) = \{u^{p+1}(x) \mid x \in E\} = \mathrm{Im}(u^{p+1}) = I_{p+1}$.
            D'après la question 2, nous savons que $I_{p+1} = I_p$.
            Donc $\mathrm{Im}(u_I) = I_p$.
            Ceci signifie que $u_I$ est surjectif.
            Puisque $u_I$ est un endomorphisme d'un espace de dimension finie $I_p$ et qu'il est surjectif, il est également injectif.
            Par conséquent, $u_I$ est un automorphisme de $I_p$.

Cette décomposition $E = K_p \oplus I_p$ est connue sous le nom de \textbf{Décomposition de Fitting}, et elle est fondamentale dans l'étude des endomorphismes en dimension finie, notamment pour la forme de Jordan ou la décomposition de Dunford.

\newpage
\chapter*{Exercice 10 : Noyau et Image d'une Composée d'Applications Linéaires (Difficulté : \textbf{}*)}

\textbf{Énoncé}

Soient $\mathbb{K}$ un corps commutatif et $E, F, G$ des $\mathbb{K}$-espaces vectoriels.
Soient $f \in \mathcal{L}(E, F)$ et $g \in \mathcal{L}(F, G)$ deux applications linéaires.
On note $\ker(\cdot)$ le noyau et $\text{Im}(\cdot)$ l'image d'une application linéaire.

1.  Démontrer que $\ker(g \circ f) = f^{-1}(\ker g \cap \text{Im } f)$.
    (Rappel : $f^{-1}(A) = \{x \in E \mid f(x) \in A\}$ pour un sous-ensemble $A \subseteq F$).

2.  Prouver que l'application $\hat{f}: \ker(g \circ f) \to \ker g \cap \text{Im } f$ définie par $\hat{f}(x) = f(x)$ est une application linéaire surjective dont le noyau est $\ker f$. En déduire l'isomorphisme de $\mathbb{K}$-espaces vectoriels : $\ker(g \circ f) / \ker f \cong \ker g \cap \text{Im } f$.

3.  Prouver que l'application $\bar{g}: \text{Im } f / (\ker g \cap \text{Im } f) \to \text{Im}(g \circ f)$ définie par $\bar{g}(y + (\ker g \cap \text{Im } f)) = g(y)$ pour tout $y \in \text{Im } f$ est un isomorphisme de $\mathbb{K}$-espaces vectoriels.
    (On vérifiera d'abord que le sous-espace $\ker g \cap \text{Im } f$ est bien un sous-espace de $\text{Im } f$ et que $\bar{g}$ est bien définie, linéaire et injective/surjective).

4.  On suppose désormais que $E, F, G$ sont des $\mathbb{K}$-espaces vectoriels de dimensions finies. On note $\dim(V)$ la dimension de $V$ et $\text{rang}(h) = \dim(\text{Im } h)$ le rang d'une application linéaire $h$.
    À partir des résultats des questions 2 et 3, établir les égalités suivantes :
    a) $\dim(\ker(g \circ f)) = \dim(\ker f) + \dim(\ker g \cap \text{Im } f)$.
    b) $\text{rang}(g \circ f) = \text{rang}(f) - \dim(\ker g \cap \text{Im } f)$.

5.  Utiliser les résultats précédents et le théorème du rang pour $f$ et $g$ pour démontrer l'inégalité de Sylvester :
    $\text{rang}(f) + \text{rang}(g) - \dim(F) \le \text{rang}(g \circ f) \le \min(\text{rang } f, \text{rang } g)$.

\section*{Correction détaillée}

1.  \textbf{Démonstration de $\ker(g \circ f) = f^{-1}(\ker g \cap \text{Im } f)$}

    Pour prouver l'égalité de deux ensembles, nous allons démontrer l'inclusion dans les deux sens.

    *   \textbf{Inclusion $\ker(g \circ f) \subseteq f^{-1}(\ker g \cap \text{Im } f)$} :
        Soit $x \in \ker(g \circ f)$. Par définition du noyau, cela signifie que $(g \circ f)(x) = 0_G$.
        Par définition de la composition, $g(f(x)) = 0_G$.
        Cela implique que $f(x) \in \ker g$.
        Par ailleurs, par définition de l'image, $f(x)$ est un élément de $\text{Im } f$.
        Donc, $f(x)$ appartient à l'intersection $\ker g \cap \text{Im } f$.
        Par définition de l'image réciproque, $x \in f^{-1}(\ker g \cap \text{Im } f)$.
        L'inclusion est démontrée.

    *   \textbf{Inclusion $f^{-1}(\ker g \cap \text{Im } f) \subseteq \ker(g \circ f)$} :
        Soit $x \in f^{-1}(\ker g \cap \text{Im } f)$. Par définition de l'image réciproque, cela signifie que $f(x) \in \ker g \cap \text{Im } f$.
        Puisque $f(x) \in \ker g$, par définition du noyau, nous avons $g(f(x)) = 0_G$.
        Par définition de la composition, $(g \circ f)(x) = 0_G$.
        Cela signifie que $x \in \ker(g \circ f)$.
        L'inclusion est démontrée.

    Ayant démontré les deux inclusions, nous concluons que $\ker(g \circ f) = f^{-1}(\ker g \cap \text{Im } f)$.

2.  \textbf{Analyse de l'application $\hat{f}$ et isomorphisme de quotient}

    Considérons l'application $\hat{f}: \ker(g \circ f) \to \ker g \cap \text{Im } f$ définie par $\hat{f}(x) = f(x)$.

    *   \textbf{$\hat{f}$ est bien définie et linéaire} :
        Le domaine de $\hat{f}$ est $\ker(g \circ f)$, qui est un sous-espace vectoriel de $E$.
        Pour tout $x \in \ker(g \circ f)$, nous avons montré à la question 1 que $f(x) \in \ker g \cap \text{Im } f$. Donc le codomaine est correct.
        L'application $f$ est linéaire, et $\hat{f}$ est simplement la restriction de $f$ à un sous-espace vectoriel, avec un codomaine restreint. Par conséquent, $\hat{f}$ est linéaire.

    *   \textbf{$\hat{f}$ est surjective} :
        Soit $y \in \ker g \cap \text{Im } f$. Puisque $y \in \text{Im } f$, il existe un $x \in E$ tel que $f(x) = y$.
        Puisque $y \in \ker g$, nous avons $g(y) = 0_G$.
        Donc, $g(f(x)) = 0_G$, ce qui signifie $(g \circ f)(x) = 0_G$.
        Par conséquent, $x \in \ker(g \circ f)$.
        Nous avons trouvé un élément $x$ dans le domaine de $\hat{f}$ tel que $\hat{f}(x) = f(x) = y$.
        Donc, $\hat{f}$ est surjective sur son codomaine $\ker g \cap \text{Im } f$.

    *   \textbf{Détermination du noyau de $\hat{f}$} :
        Par définition, $\ker \hat{f} = \{x \in \ker(g \circ f) \mid \hat{f}(x) = 0_F\}$.
        Cela signifie $\ker \hat{f} = \{x \in \ker(g \circ f) \mid f(x) = 0_F\}$.
        L'ensemble $\{x \in E \mid f(x) = 0_F\}$ est $\ker f$.
        Donc, $\ker \hat{f} = \ker(g \circ f) \cap \ker f$.
        Nous savons que si $x \in \ker f$, alors $f(x) = 0_F$. Par suite, $g(f(x)) = g(0_F) = 0_G$, ce qui implique $x \in \ker(g \circ f)$.
        Ainsi, $\ker f \subseteq \ker(g \circ f)$.
        Par conséquent, l'intersection $\ker(g \circ f) \cap \ker f$ se réduit à $\ker f$.
        Donc, $\ker \hat{f} = \ker f$.

    *   \textbf{Application du premier théorème d'isomorphisme} :
        Puisque $\hat{f}: \ker(g \circ f) \to \ker g \cap \text{Im } f$ est une application linéaire surjective de noyau $\ker f$, le premier théorème d'isomorphisme établit que
        $\ker(g \circ f) / \ker f \cong \text{Im}(\hat{f})$.
        Comme $\hat{f}$ est surjective sur son codomaine, $\text{Im}(\hat{f}) = \ker g \cap \text{Im } f$.
        Nous en déduisons l'isomorphisme : $\ker(g \circ f) / \ker f \cong \ker g \cap \text{Im } f$.

3.  \textbf{Analyse de l'application $\bar{g}$ et isomorphisme de quotient}

    Considérons l'application $\bar{g}: \text{Im } f / (\ker g \cap \text{Im } f) \to \text{Im}(g \circ f)$ définie par $\bar{g}(y + (\ker g \cap \text{Im } f)) = g(y)$.

    *   \textbf{$\ker g \cap \text{Im } f$ est un sous-espace de $\text{Im } f$} :
        $\ker g$ est un sous-espace de $F$ car c'est le noyau d'une application linéaire.
        $\text{Im } f$ est un sous-espace de $F$ car c'est l'image d'une application linéaire.
        L'intersection de deux sous-espaces vectoriels est toujours un sous-espace vectoriel.
        Puisque $\ker g \cap \text{Im } f$ est contenu dans $\text{Im } f$, il est bien un sous-espace de $\text{Im } f$.
        Le quotient $\text{Im } f / (\ker g \cap \text{Im } f)$ est donc bien défini.

    *   \textbf{$\bar{g}$ est bien définie} :
        Soient $Y_1 = y_1 + (\ker g \cap \text{Im } f)$ et $Y_2 = y_2 + (\ker g \cap \text{Im } f)$ deux représentants de la même classe dans le quotient.
        Cela signifie que $y_1 - y_2 \in \ker g \cap \text{Im } f$.
        Puisque $y_1 - y_2 \in \ker g$, nous avons $g(y_1 - y_2) = 0_G$.
        Par linéarité de $g$, $g(y_1) - g(y_2) = 0_G$, d'où $g(y_1) = g(y_2)$.
        Ainsi, $\bar{g}(Y_1) = g(y_1) = g(y_2) = \bar{g}(Y_2)$. L'application $\bar{g}$ est donc bien définie.

    *   \textbf{$\bar{g}$ est linéaire} :
        Soient $Y_1 = y_1 + (\ker g \cap \text{Im } f)$ et $Y_2 = y_2 + (\ker g \cap \text{Im } f)$ deux éléments du quotient, et $\lambda \in \mathbb{K}$.
        $\bar{g}(Y_1 + Y_2) = \bar{g}((y_1+y_2) + (\ker g \cap \text{Im } f)) = g(y_1+y_2)$.
        Par linéarité de $g$, $g(y_1+y_2) = g(y_1) + g(y_2) = \bar{g}(Y_1) + \bar{g}(Y_2)$.
        $\bar{g}(\lambda Y_1) = \bar{g}(\lambda y_1 + (\ker g \cap \text{Im } f)) = g(\lambda y_1)$.
        Par linéarité de $g$, $g(\lambda y_1) = \lambda g(y_1) = \lambda \bar{g}(Y_1)$.
        L'application $\bar{g}$ est donc linéaire.

    *   \textbf{$\bar{g}$ est surjective} :
        Soit $z \in \text{Im}(g \circ f)$. Par définition, il existe $x \in E$ tel que $z = (g \circ f)(x)$.
        Cela signifie $z = g(f(x))$.
        Posons $y = f(x)$. Alors $y \in \text{Im } f$.
        L'élément $y + (\ker g \cap \text{Im } f)$ est un élément du domaine de $\bar{g}$.
        Alors $\bar{g}(y + (\ker g \cap \text{Im } f)) = g(y) = g(f(x)) = z$.
        Donc, pour tout $z \in \text{Im}(g \circ f)$, il existe un antécédent dans le domaine de $\bar{g}$.
        L'application $\bar{g}$ est surjective.

    *   \textbf{$\bar{g}$ est injective} :
        Soit $Y = y + (\ker g \cap \text{Im } f)$ un élément du noyau de $\bar{g}$.
        Cela signifie $\bar{g}(Y) = 0_G$, c'est-à-dire $g(y) = 0_G$.
        Puisque $g(y) = 0_G$, $y$ appartient à $\ker g$.
        Par ailleurs, par la définition du domaine du quotient, $y$ appartient à $\text{Im } f$.
        Donc, $y \in \ker g \cap \text{Im } f$.
        Par conséquent, la classe $Y = y + (\ker g \cap \text{Im } f)$ est la classe zéro du quotient, i.e., $Y = \ker g \cap \text{Im } f$.
        Le noyau de $\bar{g}$ est réduit à l'élément neutre du quotient.
        L'application $\bar{g}$ est injective.

    *   \textbf{Conclusion} :
        Puisque $\bar{g}$ est linéaire, bien définie, surjective et injective, c'est un isomorphisme de $\mathbb{K}$-espaces vectoriels.
        Ainsi, $\text{Im } f / (\ker g \cap \text{Im } f) \cong \text{Im}(g \circ f)$.

4.  \textbf{Application aux dimensions finies}

    Nous supposons $E, F, G$ de dimensions finies.

    a) \textbf{Égalité pour le noyau} :
        D'après la question 2, nous avons l'isomorphisme $\ker(g \circ f) / \ker f \cong \ker g \cap \text{Im } f$.
        Pour des espaces vectoriels de dimensions finies, si $V_1 \cong V_2$, alors $\dim(V_1) = \dim(V_2)$.
        De plus, pour un quotient $V/W$, $\dim(V/W) = \dim(V) - \dim(W)$.
        Donc, $\dim(\ker(g \circ f) / \ker f) = \dim(\ker(g \circ f)) - \dim(\ker f)$.
        Par conséquent, $\dim(\ker(g \circ f)) - \dim(\ker f) = \dim(\ker g \cap \text{Im } f)$.
        Ce qui donne l'égalité souhaitée :
        $\dim(\ker(g \circ f)) = \dim(\ker f) + \dim(\ker g \cap \text{Im } f)$.

    b) \textbf{Égalité pour l'image (rang)} :
        D'après la question 3, nous avons l'isomorphisme $\text{Im } f / (\ker g \cap \text{Im } f) \cong \text{Im}(g \circ f)$.
        En appliquant le même principe des dimensions finies :
        $\dim(\text{Im } f / (\ker g \cap \text{Im } f)) = \dim(\text{Im } f) - \dim(\ker g \cap \text{Im } f)$.
        Et $\dim(\text{Im}(g \circ f)) = \text{rang}(g \circ f)$.
        De plus, $\dim(\text{Im } f) = \text{rang}(f)$.
        Par conséquent, $\text{rang}(g \circ f) = \text{rang}(f) - \dim(\ker g \cap \text{Im } f)$.

5.  \textbf{Démonstration de l'inégalité de Sylvester}

    L'inégalité de Sylvester est de la forme $\text{rang}(f) + \text{rang}(g) - \dim(F) \le \text{rang}(g \circ f) \le \min(\text{rang } f, \text{rang } g)$.
    Nous allons prouver les deux parties de l'inégalité.

    *   \textbf{Borne supérieure : $\text{rang}(g \circ f) \le \min(\text{rang } f, \text{rang } g)$} :
        1.  \textbf{$\text{rang}(g \circ f) \le \text{rang}(f)$} :
            D'après la question 4b, $\text{rang}(g \circ f) = \text{rang}(f) - \dim(\ker g \cap \text{Im } f)$.
            Comme $\ker g \cap \text{Im } f$ est un sous-espace vectoriel, sa dimension est non-négative : $\dim(\ker g \cap \text{Im } f) \ge 0$.
            Par conséquent, $\text{rang}(g \circ f) \le \text{rang}(f)$.
            On peut aussi observer que $\text{Im}(g \circ f) = g(\text{Im } f)$. L'image d'un sous-espace par une application linéaire a une dimension inférieure ou égale à la dimension de ce sous-espace. Ainsi, $\dim(g(\text{Im } f)) \le \dim(\text{Im } f)$, ce qui signifie $\text{rang}(g \circ f) \le \text{rang}(f)$.

        2.  \textbf{$\text{rang}(g \circ f) \le \text{rang}(g)$} :
            Comme précédemment, $\text{Im}(g \circ f) = g(\text{Im } f)$.
            Puisque $\text{Im } f \subseteq F$, tout vecteur dans $g(\text{Im } f)$ est aussi un vecteur dans $g(F)$, c'est-à-dire dans $\text{Im } g$.
            Donc $\text{Im}(g \circ f)$ est un sous-espace de $\text{Im } g$.
            Par conséquent, $\dim(\text{Im}(g \circ f)) \le \dim(\text{Im } g)$, ce qui signifie $\text{rang}(g \circ f) \le \text{rang}(g)$.

        Ces deux inégalités prouvent que $\text{rang}(g \circ f) \le \min(\text{rang } f, \text{rang } g)$.

    *   \textbf{Borne inférieure : $\text{rang}(f) + \text{rang}(g) - \dim(F) \le \text{rang}(g \circ f)$} :
        D'après la question 4b, nous avons $\text{rang}(g \circ f) = \text{rang}(f) - \dim(\ker g \cap \text{Im } f)$.
        Nous devons maintenant borner $\dim(\ker g \cap \text{Im } f)$.
        Le sous-espace $\ker g \cap \text{Im } f$ est un sous-espace de $\ker g$.
        Donc, $\dim(\ker g \cap \text{Im } f) \le \dim(\ker g)$.
        En utilisant cette inégalité dans l'expression de $\text{rang}(g \circ f)$ :
        $\text{rang}(g \circ f) \ge \text{rang}(f) - \dim(\ker g)$.

        Maintenant, nous utilisons le théorème du rang pour l'application linéaire $g: F \to G$.
        Le théorème du rang stipule que $\dim(F) = \dim(\ker g) + \dim(\text{Im } g)$.
        Puisque $\dim(\text{Im } g) = \text{rang}(g)$, nous avons $\dim(F) = \dim(\ker g) + \text{rang}(g)$.
        D'où $\dim(\ker g) = \dim(F) - \text{rang}(g)$.

        Substituons cette expression de $\dim(\ker g)$ dans l'inégalité pour $\text{rang}(g \circ f)$ :
        $\text{rang}(g \circ f) \ge \text{rang}(f) - (\dim(F) - \text{rang}(g))$.
        En réarrangeant les termes, nous obtenons :
        $\text{rang}(g \circ f) \ge \text{rang}(f) + \text{rang}(g) - \dim(F)$.

    Nous avons ainsi démontré les deux parties de l'inégalité de Sylvester.

\newpage
\chapter*{Travaux Pratiques}
\addcontentsline{toc}{chapter}{Travaux Pratiques}
\chapter*{TP 1 : Fondamentaux de l'Algèbre Linéaire pour l'IA : Applications, Noyau, Image et Théorème du Rang "From Scratch"}

\section*{Objectif}

Ce TP a pour but de vous faire implémenter, à partir de zéro et sans aucune librairie externe (comme NumPy ou SciPy), les concepts fondamentaux de l'algèbre linéaire que sont les applications linéaires, leur noyau (kernel), leur image (range) et le théorème du rang.

\subsection*{Objectif Mathématique}
*   Comprendre et représenter les vecteurs et matrices en Python pur (listes de nombres et listes de listes).
*   Implémenter les opérations de base de l'algèbre linéaire : addition de vecteurs, multiplication scalaire-vecteur, produit scalaire, et surtout la multiplication matrice-vecteur.
*   Saisir la définition formelle et l'implémentation d'une application linéaire comme une transformation représentée par une multiplication matricielle (\_\_INLINE\_CODE\_0\_\_).
*   Définir et vérifier l'appartenance d'un vecteur au noyau d'une application linéaire.
*   Définir et vérifier l'appartenance d'un vecteur à l'image d'une application linéaire (avec une méthode adaptée aux contraintes "from scratch").
*   Comprendre et illustrer le Théorème du Rang qui établit une relation fondamentale entre la dimension de l'espace de départ, la dimension du noyau (nullité) et la dimension de l'image (rang).

\subsection*{Objectif IA}
Les concepts d'algèbre linéaire sont des piliers de nombreux algorithmes en Intelligence Artificielle et en Machine Learning.
*   \textbf{Applications Linéaires}: Elles sont le cœur des couches "Dense" ou "Linéaires" des réseaux de neurones. Comprendre \_\_INLINE\_CODE\_1\_\_ est crucial pour saisir comment les données d'entrée sont transformées et projetées dans des espaces de caractéristiques de plus grande ou plus petite dimension.
*   \textbf{Noyau (\_\_INLINE\_CODE\_2\_\_)}: Le noyau représente l'ensemble des entrées qui sont "écrasées" ou projetées sur le vecteur nul par l'application linéaire. En IA, cela peut être interprété comme une perte d'information, des caractéristiques redondantes, ou des directions dans l'espace d'entrée qui n'influencent pas la sortie. Il est lié aux concepts de réduction de dimension et de compression de données.
*   \textbf{Image (\_\_INLINE\_CODE\_3\_\_)}: L'image est l'ensemble de toutes les sorties possibles de l'application linéaire. Elle définit l'espace de "capacité expressive" d'un modèle linéaire. Comprendre l'image permet de savoir quelles transformations et quelles sorties un modèle peut effectivement produire, et quelle est la "richesse" de son espace de sortie.
*   \textbf{Théorème du Rang}: Ce théorème (\_\_INLINE\_CODE\_4\_\_) fournit une relation fondamentale entre la dimension de l'espace d'entrée, la dimension de l'information "perdue" (noyau) et la dimension de l'information "transformée" (image). Il est essentiel pour l'analyse de la complexité, de la redondance et de la capacité de représentation des modèles d'apprentissage automatique. Une faible nullité (grand rang) signifie que l'application préserve beaucoup d'information, tandis qu'une grande nullité (petit rang) indique une forte compression ou perte d'information.

\section*{Implémentation en Python "From Scratch"}

Nous allons implémenter les fonctions nécessaires et les valider avec des assertions mathématiques pour prouver leur correction.

\begin{verbatim}
import math

# --- Fonctions Utilitaires de Base pour l'Algebre Lineaire ---

def are_close(val1, val2, tolerance=1e-9):
    """
    Verifie si deux valeurs flottantes sont considerees comme egales
    compte tenu des imprecisions de la virgule flottante.
    """
    return abs(val1 - val2) < tolerance

def is_zero_vector(v, tolerance=1e-9):
    """
    Verifie si un vecteur est le vecteur nul (toutes ses composantes sont proches de zero).
    """
    return all(are_close(x, 0.0, tolerance) for x in v)

def vector_add(v1, v2):
    """
    Additionne deux vecteurs composante par composante.
    Les vecteurs doivent avoir la meme dimension.
    """
    if len(v1) != len(v2):
        raise ValueError("Les vecteurs doivent avoir la meme dimension pour etre additionnes.")
    return [v1[i] + v2[i] for i in range(len(v1))]

def scalar_vector_multiply(scalar, v):
    """
    Multiplie un vecteur par un scalaire.
    """
    return [scalar * x for x in v]

def dot_product(v1, v2):
    """
    Calcule le produit scalaire (produit interne) de deux vecteurs.
    Les vecteurs doivent avoir la meme dimension.
    """
    if len(v1) != len(v2):
        raise ValueError("Les vecteurs doivent avoir la meme dimension pour le produit scalaire.")
    return sum(v1[i] * v2[i] for i in range(len(v1)))

def matrix_vector_multiply(matrix, vector):
    """
    Multiplie une matrice par un vecteur.
    Cette operation est la representation standard d'une application lineaire f(x) = Ax.

    Args:
        matrix (list of list of float): La matrice A, representee comme une liste de lignes.
        vector (list of float): Le vecteur x.

    Returns:
        list of float: Le vecteur resultant y = Ax.
    """
    num_rows = len(matrix)
    if num_rows == 0:
        if len(vector) != 0:
            raise ValueError("La matrice est vide mais le vecteur ne l'est pas.")
        return [] # Matrice vide, vecteur vide -> resultat vide

    num_cols = len(matrix[0])
    if num_cols == 0 and num_rows > 0:
        raise ValueError("La matrice a des lignes mais pas de colonnes (vide).")


    if len(vector) != num_cols:
        raise ValueError(
            f"Le nombre de colonnes de la matrice ({num_cols}) "
            f"doit etre egal a la dimension du vecteur ({len(vector)})."
        )

    result_vector = [0.0] * num_rows
    for i in range(num_rows):
        # Chaque composante du vecteur resultat est le produit scalaire de la i-eme ligne de la matrice
        # avec le vecteur d'entree.
        result_vector[i] = dot_product(matrix[i], vector)
    return result_vector

# --- Verification de l'implementation des utilitaires ---
print("--- Test des Fonctions Utilitaires de Base ---")
v_a = [1.0, 2.0, 3.0]
v_b = [4.0, 5.0, 6.0]
s = 2.0
matrix_A_test = [[1.0, 2.0, 3.0],
                 [4.0, 5.0, 6.0]]
vector_x_test = [1.0, 0.5, 0.1]

assert vector_add(v_a, v_b) == [5.0, 7.0, 9.0], "Test vector_add echoue"
assert scalar_vector_multiply(s, v_a) == [2.0, 4.0, 6.0], "Test scalar_vector_multiply echoue"
assert are_close(dot_product(v_a, v_b), 32.0), "Test dot_product echoue"

expected_mv_result = [1.0*1.0 + 2.0*0.5 + 3.0*0.1, 4.0*1.0 + 5.0*0.5 + 6.0*0.1] # [1 + 1 + 0.3, 4 + 2.5 + 0.6] = [2.3, 7.1]
mv_result = matrix_vector_multiply(matrix_A_test, vector_x_test)
assert len(mv_result) == len(expected_mv_result) and \
       all(are_close(val, expected) for val, expected in zip(mv_result, expected_mv_result)), \
       f"Test matrix_vector_multiply echoue. Attendu {expected_mv_result}, Obtenu {mv_result}"

assert is_zero_vector([0.0, 0.0, 0.0]), "Test is_zero_vector pour zero echoue"
assert not is_zero_vector([0.1, 0.0, 0.0]), "Test is_zero_vector pour non-zero echoue"
print("Tests des Fonctions Utilitaires reussis !")

# --- 1. Application Lineaire ---

def linear_application(matrix, vector):
    """
    Definit une application lineaire f(x) = Ax.
    C'est conceptuellement un alias pour matrix_vector_multiply pour la clarte.
    """
    return matrix_vector_multiply(matrix, vector)

print("\n--- 1. Application Lineaire ---")

# Matrice pour notre application lineaire, R^2 -> R^2
A_linear = [[1.0, 2.0],
            [3.0, 4.0]]

# Vecteurs d'entree pour tester les proprietes de linearite
x1 = [1.0, 0.0]
x2 = [0.0, 1.0]
x_sum = vector_add(x1, x2) # [1.0, 1.0]
c = 2.0
x_scaled = scalar_vector_multiply(c, x1) # [2.0, 0.0]

# Appliquer la transformation
y1 = linear_application(A_linear, x1) # [1.0, 3.0]
y2 = linear_application(A_linear, x2) # [2.0, 4.0]
y_sum = linear_application(A_linear, x_sum) # [3.0, 7.0]
y_scaled = linear_application(A_linear, x_scaled) # [2.0, 6.0]

# Test de la propriete d'additivite : f(x1 + x2) = f(x1) + f(x2)
expected_y_sum = vector_add(y1, y2) # [1+2, 3+4] = [3.0, 7.0]
assert len(y_sum) == len(expected_y_sum) and \
       all(are_close(val, expected) for val, expected in zip(y_sum, expected_y_sum)), \
    f"Test d'additivite de l'application lineaire echoue. f(x1+x2)={y_sum}, f(x1)+f(x2)={expected_y_sum}"
print(f"f({x1} + {x2}) = {y_sum}")
print(f"f({x1}) + f({x2}) = {expected_y_sum}")
print("Test d'additivite de l'application lineaire reussi.")

# Test de la propriete d'homogeneite : f(c * x1) = c * f(x1)
expected_y_scaled = scalar_vector_multiply(c, y1) # 2 * [1.0, 3.0] = [2.0, 6.0]
assert len(y_scaled) == len(expected_y_scaled) and \
       all(are_close(val, expected) for val, expected in zip(y_scaled, expected_y_scaled)), \
    f"Test d'homogeneite de l'application lineaire echoue. f(c*x1)={y_scaled}, c*f(x1)={expected_y_scaled}"
print(f"f({c} * {x1}) = {y_scaled}")
print(f"{c} * f({x1}) = {expected_y_scaled}")
print("Test d'homogeneite de l'application lineaire reussi.")

print("L'implementation de l'application lineaire est correcte et respecte les proprietes de linearite.")

# --- 2. Noyau (Kernel) ---

def is_in_kernel(matrix, vector, tolerance=1e-9):
    """
    Verifie si un vecteur appartient au noyau d'une application lineaire definie par une matrice.
    Un vecteur v est dans le noyau de A si Av = 0.
    """
    result = matrix_vector_multiply(matrix, vector)
    return is_zero_vector(result, tolerance)

print("\n--- 2. Noyau (Kernel) ---")

# Exemple 1 : Matrice identite (le noyau est reduit au vecteur nul : seulement le vecteur nul)
I_2x2 = [[1.0, 0.0],
         [0.0, 1.0]]
v_zero = [0.0, 0.0]
v_non_zero = [1.0, 2.0]

assert is_in_kernel(I_2x2, v_zero), "Le vecteur nul doit etre dans le noyau de la matrice identite."
assert not is_in_kernel(I_2x2, v_non_zero), "Un vecteur non nul ne doit pas etre dans le noyau de la matrice identite."
print("Tests pour la matrice identite reussis : le noyau est reduit au vecteur nul.")

# Exemple 2 : Matrice avec un noyau non reduit au vecteur nul
# Matrice A_kernel : R^2 -> R^2, ou la deuxieme colonne est un multiple de la premiere.
# Cela signifie que les colonnes ne sont pas lineairement independantes, donc il y a un noyau non reduit au vecteur nul.
A_kernel = [[1.0, 2.0],
            [2.0, 4.0]]
# Pour cette matrice, tout vecteur de la forme k * [-2, 1] est dans le noyau.
# Par exemple, si k=1, x_k = [-2, 1].
# A_kernel @ x_k = [1*(-2) + 2*1, 2*(-2) + 4*1] = [-2+2, -4+4] = [0, 0]
x_kernel_candidate = [-2.0, 1.0]

assert is_in_kernel(A_kernel, x_kernel_candidate), \
    f"Le vecteur {x_kernel_candidate} devrait etre dans le noyau de A_kernel."
print(f"Verification du noyau pour A_kernel: {A_kernel}")
print(f"A_kernel @ {x_kernel_candidate} = {matrix_vector_multiply(A_kernel, x_kernel_candidate)}")
print(f"Le vecteur {x_kernel_candidate} est bien dans le noyau de la matrice A_kernel.")

# Un autre vecteur dans le noyau (un multiple du premier)
x_kernel_candidate_2 = [-4.0, 2.0] # 2 * [-2, 1]
assert is_in_kernel(A_kernel, x_kernel_candidate_2), \
    f"Le vecteur {x_kernel_candidate_2} (multiple) devrait etre dans le noyau de A_kernel."
print(f"Le vecteur {x_kernel_candidate_2} est bien dans le noyau de la matrice A_kernel.")

# Un vecteur qui n'est PAS dans le noyau
x_not_kernel = [1.0, 1.0]
assert not is_in_kernel(A_kernel, x_not_kernel), \
    f"Le vecteur {x_not_kernel} ne devrait PAS etre dans le noyau de A_kernel."
print(f"A_kernel @ {x_not_kernel} = {matrix_vector_multiply(A_kernel, x_not_kernel)}")
print(f"Le vecteur {x_not_kernel} n'est PAS dans le noyau de la matrice A_kernel.")

print("Implementation et tests du noyau reussis.")

# --- 3. Image (Range) ---

def is_in_image(matrix, vector_y, input_vector_x=None, tolerance=1e-9):
    """
    Verifie si un vecteur y appartient a l'image d'une application lineaire A.
    Ceci est verifie si y = Ax pour un certain vecteur x.

    Pour une implementation "from scratch" sans resolution de systeme lineaire (Ax = y),
    nous testons si un 'input_vector_x' (qui est cense produire 'vector_y') est fourni.
    Si 'input_vector_x' n'est pas fourni, cette fonction ne peut pas determiner l'appartenance
    pour un vecteur 'y' arbitraire sans resoudre Ax = y, ce qui est une tache plus complexe
    (necessitant l'elimination de Gauss ou des factorisations matricielles) et depasse
    le cadre de ce TP 1.

    Args:
        matrix (list of list of float): La matrice A de l'application lineaire.
        vector_y (list of float): Le vecteur a verifier s'il est dans l'image.
        input_vector_x (list of float, optional): Un vecteur x tel que A @ x = y.
                                                  Obligatoire pour cette implementation simple.
        tolerance (float): La tolerance pour la comparaison des flottants.

    Returns:
        bool: True si A @ input_vector_x est egal a vector_y, False sinon.
    """
    if input_vector_x is None:
        raise ValueError(
            "Pour verifier si un vecteur est dans l'image 'from scratch' sans solveur, "
            "vous devez fournir un 'input_vector_x' tel que A @ x = y. "
            "La resolution generale de Ax=y n'est pas implementee ici."
        )

    result_y_computed = matrix_vector_multiply(matrix, input_vector_x)

    if len(result_y_computed) != len(vector_y):
        return False # Les dimensions ne correspondent pas

    return all(are_close(val_res, val_expected, tolerance) for val_res, val_expected in zip(result_y_computed, vector_y))


print("\n--- 3. Image (Range) ---")

# Matrice pour l'image : A_image : R^2 -> R^3
# Les colonnes de A_image sont [1,0,1] et [0,1,1]. L'image est le plan engendre par ces deux vecteurs dans R^3.
A_image = [[1.0, 0.0],
           [0.0, 1.0],
           [1.0, 1.0]]

# Vecteur d'entree qui genere un vecteur dans l'image
x_image_input = [1.0, 2.0] # Un vecteur de R^2
y_in_image = linear_application(A_image, x_image_input) # y = A_image @ x_image_input = [1, 2, 3]

# Verifier que y_in_image est bien dans l'image en utilisant l'input_vector_x
assert is_in_image(A_image, y_in_image, x_image_input), \
    f"Le vecteur {y_in_image} devrait etre dans l'image de A_image avec l'entree {x_image_input}."
print(f"Matrice A_image: {A_image}")
print(f"Vecteur d'entree x_image_input: {x_image_input}")
print(f"Vecteur de sortie y_in_image (A_image @ x_image_input): {y_in_image}")
print(f"Le vecteur {y_in_image} est bien dans l'image de A_image.")

# Un autre vecteur dans l'image (obtenu avec une autre entree)
x_image_input_2 = [3.0, -1.0]
y_in_image_2 = linear_application(A_image, x_image_input_2) # y = [3, -1, 2]
assert is_in_image(A_image, y_in_image_2, x_image_input_2), \
    f"Le vecteur {y_in_image_2} devrait etre dans l'image de A_image avec l'entree {x_image_input_2}."
print(f"Un autre vecteur {y_in_image_2} est aussi dans l'image de A_image.")

# Un vecteur qui n'est PAS dans l'image (par exemple, non sur le plan genere par les colonnes de A_image)
# L'image de A_image est le plan defini par les vecteurs [1,0,1] et [0,1,1].
# Un vecteur [a,b,c] est dans l'image si c = a + b.
y_not_in_image = [1.0, 1.0, 3.0] # Ici, 3 != 1 + 1. Ce vecteur n'est pas sur le plan.

print(f"\nVecteur y_not_in_image (conceptuellement non dans l'image): {y_not_in_image}")
print("Sans un solveur pour Ax=y, il est difficile de verifier qu'un vecteur N'EST PAS dans l'image 'from scratch'.")
print("Notre fonction `is_in_image` est concue pour la *verification* etant donne un x.")
try:
    is_in_image(A_image, y_not_in_image) # Cela doit lever une erreur car input_vector_x est manquant
except ValueError as e:
    print(f"Erreur attendue (pas d'input_vector_x fourni pour la verification negative) : {e}")

print("Implementation et tests de l'image (avec x fourni) reussis.")


# --- 4. Theoreme du Rang ---

# Le Theoreme du Rang : dim(Espace de Depart) = dim(Noyau) + dim(Image)
# (ou dim(Noyau) est la nullite et dim(Image) est le rang)

# Comme l'implementation robuste de la determination du rang et de la nullite
# pour une matrice arbitraire (qui necessite la reduction de Gauss-Jordan)
# est tres complexe "from scratch" et depasse le cadre de ce TP 1,
# nous allons illustrer le theoreme avec des matrices ou le rang et la nullite
# sont connus par inspection ou par construction simple.

print("\n--- 4. Theoreme du Rang ---")

# Exemple 1 : Matrice identite 2x2
# Application de R^2 -> R^2
I_2x2 = [[1.0, 0.0],
         [0.0, 1.0]]
dim_domain_I = 2        # Dimension de l'espace de depart (R^2)
rank_I = 2              # Deux colonnes lineairement independantes => image de dimension 2
nullity_I = 0           # Seul le vecteur nul est dans le noyau => noyau de dimension 0

print("\nExemple 1 : Matrice Identite 2x2 (Transformation qui ne perd pas d'information)")
print(f"Matrice: {I_2x2}")
print(f"Dimension de l'espace de depart (R^2): {dim_domain_I}")
print(f"Rang (dimension de l'image): {rank_I}")
print(f"Nullite (dimension du noyau): {nullity_I}")

assert dim_domain_I == rank_I + nullity_I, "Le Theoreme du Rang ne tient pas pour la matrice identite."
print(f"Verification du Theoreme du Rang : {dim_domain_I} == {rank_I} + {nullity_I} (Vrai)")


# Exemple 2 : Matrice de projection 2x2 sur un axe (image de dimension 1, perte d'information)
# Application de R^2 -> R^2
# Projette sur l'axe des x
P_x = [[1.0, 0.0],
       [0.0, 0.0]]
dim_domain_P = 2        # Dimension de l'espace de depart (R^2)
rank_P = 1              # L'image est l'axe des x (engendree par [1,0]), dimension 1
# Le noyau est l'axe des y (vecteurs de la forme [0, k]), dimension 1.
# Car P_x @ [0, k] = [0, 0]
nullity_P = 1

print("\nExemple 2 : Matrice de Projection 2x2 (Perte d'information le long de l'axe y)")
print(f"Matrice: {P_x}")
print(f"Dimension de l'espace de depart (R^2): {dim_domain_P}")
print(f"Rang (dimension de l'image): {rank_P}")
print(f"Nullite (dimension du noyau): {nullity_P}")

assert dim_domain_P == rank_P + nullity_P, "Le Theoreme du Rang ne tient pas pour la matrice de projection."
print(f"Verification du Theoreme du Rang : {dim_domain_P} == {rank_P} + {nullity_P} (Vrai)")


# Exemple 3 : Matrice 3x2 avec rang 2 (pas de perte de dimension si l'espace de depart est R^2)
# Application de R^2 -> R^3 (injection si le noyau est reduit au vecteur nul)
# C'est notre A_image precedente
A_rank_3x2 = [[1.0, 0.0],
              [0.0, 1.0],
              [1.0, 1.0]]
dim_domain_A_3x2 = 2    # L'espace de depart est R^2
rank_A_3x2 = 2          # Les deux colonnes sont lineairement independantes dans R^3. L'image est un plan dans R^3.
nullity_A_3x2 = 0       # Le noyau est reduit au vecteur nul, seul le vecteur nul de R^2 donne 0 en R^3.

print("\nExemple 3 : Matrice 3x2 (R^2 -> R^3, pas de perte d'information de l'espace de depart)")
print(f"Matrice: {A_rank_3x2}")
print(f"Dimension de l'espace de depart (R^2): {dim_domain_A_3x2}")
print(f"Rang (dimension de l'image, un plan dans R^3): {rank_A_3x2}")
print(f"Nullite (dimension du noyau, seulement le vecteur nul de R^2): {nullity_A_3x2}")

assert dim_domain_A_3x2 == rank_A_3x2 + nullity_A_3x2, "Le Theoreme du Rang ne tient pas pour la matrice 3x2."
print(f"Verification du Theoreme du Rang : {dim_domain_A_3x2} == {rank_A_3x2} + {nullity_A_3x2} (Vrai)")


# Exemple 4 : Matrice 2x3 de rang 1 (reduit fortement la dimension, compression)
# Application de R^3 -> R^2
# Toutes les lignes sont des multiples l'une de l'autre
A_rank_2x3 = [[1.0, 2.0, 3.0],
              [2.0, 4.0, 6.0]]
dim_domain_A_2x3 = 3    # L'espace de depart est R^3
rank_A_2x3 = 1          # L'image est une droite dans R^2 (engendree par [1,2]). Toutes les colonnes sont multiples de la premiere.
# La nullite est dim(domain) - rank = 3 - 1 = 2.
# Le noyau est un plan dans R^3. Par exemple, [-2, 1, 0] et [-3, 0, 1] forment une base du noyau.
nullity_A_2x3 = 2

print("\nExemple 4 : Matrice 2x3 (R^3 -> R^2, compression forte, perte d'information)")
print(f"Matrice: {A_rank_2x3}")
print(f"Dimension de l'espace de depart (R^3): {dim_domain_A_2x3}")
print(f"Rang (dimension de l'image, une droite dans R^2): {rank_A_2x3}")
print(f"Nullite (dimension du noyau, un plan dans R^3): {nullity_A_2x3}")

assert dim_domain_A_2x3 == rank_A_2x3 + nullity_A_2x3, "Le Theoreme du Rang ne tient pas pour la matrice 2x3."
print(f"Verification du Theoreme du Rang : {dim_domain_A_2x3} == {rank_A_2x3} + {nullity_A_2x3} (Vrai)")


print("\nConclusion des tests du Theoreme du Rang : les relations dimensionnelles sont verifiees pour ces exemples.")
print("\nFelicitations ! Tous les tests pour le TP 1 ont ete reussis. Vous avez implemente des concepts cles de l'algebre lineaire \"from scratch\"!")
\end{verbatim}

\newpage
\chapter*{TP 2 : Applications Linéaires, Noyau, Image et Théorème du Rang}

\section*{Objectifs du TP}

Ce TP a pour objectif d'implémenter de manière "from scratch" (sans l'aide de bibliothèques tierces comme NumPy) les concepts fondamentaux d'algèbre linéaire que sont les applications linéaires, leur noyau (null space), leur image (column space) et le théorème du rang.

\subsection*{Objectifs Mathématiques}

*   Comprendre et coder les opérations de base sur les vecteurs et les matrices (addition, multiplication scalaire, multiplication matrice-vecteur).
*   Implémenter l'algorithme d'élimination de Gauss-Jordan pour réduire une matrice à sa forme échelonnée réduite par lignes (RREF).
*   À partir de la RREF, déterminer une base pour le noyau d'une application linéaire (espace nul de la matrice associée).
*   À partir de la RREF, déterminer une base pour l'image d'une application linéaire (espace colonne de la matrice associée).
*   Calculer le rang d'une matrice (dimension de l'image) et la nullité (dimension du noyau).
*   Vérifier le théorème du rang (\_\_INLINE\_CODE\_0\_\_).

\subsection*{Objectifs IA}

Les concepts d'applications linéaires, de noyau et d'image sont omniprésents en IA :
*   \textbf{Applications linéaires} modélisent les couches fully-connected des réseaux de neurones, les transformations de caractéristiques, les modèles de régression linéaire, ou les opérateurs de convolution.
*   \textbf{L'image} d'une application linéaire représente l'espace des sorties possibles, c'est-à-dire les "informations transformées" ou les "caractéristiques extraites" par un modèle. Sa dimension (le rang) indique la complexité, la capacité de transformation, ou les degrés de liberté d'un modèle. En PCA, l'espace des composantes principales est l'image d'une projection linéaire.
*   \textbf{Le noyau} représente les entrées qui sont "écrasées" ou "rendues invisibles" par la transformation, c'est-à-dire les "informations perdues" ou les "invariances" du modèle. Comprendre le noyau aide à déceler les limitations, les redondances ou les biais d'un modèle. Par exemple, dans une tâche de classification d'images, le noyau pourrait représenter des transformations de l'image qui ne changent pas la prédiction du modèle.
*   Le \textbf{théorème du rang} est fondamental pour la réduction de dimensionnalité, l'analyse des architectures de réseaux de neurones (où un rang faible peut indiquer une capacité limitée), et la compréhension des relations entre les données d'entrée et de sortie d'un système.

---

\section*{Implémentation en Python Pur}

Nous allons définir des fonctions pour gérer les vecteurs et les matrices, et implémenter toutes les opérations nécessaires "from scratch".

\begin{verbatim}
import math

# --- Configuration pour les comparaisons de flottants ---
EPSILON = 1e-9 # Petite valeur pour gerer les imprecisions des flottants

# --- Fonctions utilitaires de base sur les vecteurs et matrices ---

def are_almost_equal(a: float, b: float) -> bool:
    """Verifie si deux flottants sont presque egaux."""
    return abs(a - b) < EPSILON

def vec_add(v1: list[float], v2: list[float]) -> list[float]:
    """Additionne deux vecteurs."""
    assert len(v1) == len(v2), "Les vecteurs doivent avoir la meme dimension pour l'addition."
    return [v1[i] + v2[i] for i in range(len(v1))]

def vec_sub(v1: list[float], v2: list[float]) -> list[float]:
    """Soustrait deux vecteurs."""
    assert len(v1) == len(v2), "Les vecteurs doivent avoir la meme dimension pour la soustraction."
    return [v1[i] - v2[i] for i in range(len(v1))]

def scalar_vec_mul(scalar: float, v: list[float]) -> list[float]:
    """Multiplie un vecteur par un scalaire."""
    return [scalar * x for x in v]

def dot_product(v1: list[float], v2: list[float]) -> float:
    """Calcule le produit scalaire de deux vecteurs."""
    assert len(v1) == len(v2), "Les vecteurs doivent avoir la meme dimension pour le produit scalaire."
    return sum(v1[i] * v2[i] for i in range(len(v1)))

def vec_norm_squared(v: list[float]) -> float:
    """Calcule la norme euclidienne au carre d'un vecteur."""
    return sum(x*x for x in v)

def vec_norm(v: list[float]) -> float:
    """Calcule la norme euclidienne d'un vecteur."""
    return math.sqrt(vec_norm_squared(v))

def is_zero_vector(v: list[float]) -> bool:
    """Verifie si un vecteur est le vecteur nul."""
    return all(are_almost_equal(x, 0.0) for x in v)

def get_row(matrix: list[list[float]], row_idx: int) -> list[float]:
    """Recupere une ligne de la matrice."""
    return matrix[row_idx]

def get_col(matrix: list[list[float]], col_idx: int) -> list[float]:
    """Recupere une colonne de la matrice."""
    return [matrix[i][col_idx] for i in range(len(matrix))]

def matrix_transpose(matrix: list[list[float]]) -> list[list[float]]:
    """Calcule la transposee d'une matrice."""
    if not matrix or not matrix[0]:
        return []
    num_rows = len(matrix)
    num_cols = len(matrix[0])
    transposed = [[0.0 for _ in range(num_rows)] for _ in range(num_cols)]
    for i in range(num_rows):
        for j in range(num_cols):
            transposed[j][i] = matrix[i][j]
    return transposed

def mat_vec_mul(matrix: list[list[float]], vector: list[float]) -> list[float]:
    """Multiplie une matrice par un vecteur."""
    if not matrix:
        return []
    num_rows = len(matrix)
    num_cols = len(matrix[0])
    assert num_cols == len(vector), "Le nombre de colonnes de la matrice doit correspondre a la dimension du vecteur."
    result = [0.0] * num_rows
    for i in range(num_rows):
        result[i] = dot_product(matrix[i], vector)
    return result

# --- Implementation de l'elimination de Gauss-Jordan (RREF) ---

def swap_rows(matrix: list[list[float]], r1: int, r2: int):
    """Echange deux lignes d'une matrice."""
    matrix[r1], matrix[r2] = matrix[r2], matrix[r1]

def multiply_row(matrix: list[list[float]], r: int, scalar: float):
    """Multiplie une ligne d'une matrice par un scalaire."""
    for j in range(len(matrix[0])):
        matrix[r][j] *= scalar

def add_multiple_of_row(matrix: list[list[float]], target_row: int, source_row: int, scalar: float):
    """Ajoute un multiple d'une ligne a une autre ligne."""
    for j in range(len(matrix[0])):
        matrix[target_row][j] += scalar * matrix[source_row][j]

def get_rref(matrix_in: list[list[float]]) -> list[list[float]]:
    """
    Convertit une matrice en sa forme echelonnee reduite par lignes (RREF).
    Implemente l'algorithme de Gauss-Jordan.
    """
    matrix = [row[:] for row in matrix_in] # Creer une copie pour ne pas modifier l'original

    if not matrix or not matrix[0]:
        return []

    num_rows = len(matrix)
    num_cols = len(matrix[0])

    pivot_row = 0
    for col in range(num_cols):
        if pivot_row >= num_rows:
            break

        # 1. Trouver une ligne avec un element non nul dans la colonne courante (pivot)
        pivot_candidate_row = pivot_row
        while pivot_candidate_row < num_rows and are_almost_equal(matrix[pivot_candidate_row][col], 0.0):
            pivot_candidate_row += 1

        if pivot_candidate_row == num_rows:
            # Aucune ligne non nulle trouvee pour cette colonne, passer a la suivante
            continue

        # 2. Echanger la ligne trouvee avec la ligne pivot_row
        swap_rows(matrix, pivot_row, pivot_candidate_row)

        # 3. Normaliser la ligne pivot pour que le pivot soit 1
        pivot_val = matrix[pivot_row][col]
        # Eviter la division par zero si un pivot devient ~0 a cause d'erreurs flottantes intermediaires
        if are_almost_equal(pivot_val, 0.0):
            continue # Si le pivot est presque zero, il ne peut pas etre un pivot valide pour cette colonne.
        multiply_row(matrix, pivot_row, 1.0 / pivot_val)

        # 4. Eliminer les autres elements de la colonne (rendre les autres elements 0)
        for r in range(num_rows):
            if r != pivot_row:
                factor = -matrix[r][col]
                add_multiple_of_row(matrix, r, pivot_row, factor)

        pivot_row += 1

    return matrix

# --- Fonctions pour le noyau, l'image et le rang ---

def get_pivot_columns(rref_matrix: list[list[float]]) -> list[int]:
    """
    Determine les indices des colonnes pivots d'une matrice en RREF.
    Une colonne est pivot si elle contient un 1 qui est le premier element non nul de sa ligne.
    """
    if not rref_matrix or not rref_matrix[0]:
        return []

    pivot_cols = []
    num_rows = len(rref_matrix)
    num_cols = len(rref_matrix[0])

    # Track the last pivot column found to ensure we only look to the right for the next pivot
    current_pivot_col = -1
    for r in range(num_rows):
        found_pivot_in_row = False
        for c in range(current_pivot_col + 1, num_cols):
            if are_almost_equal(rref_matrix[r][c], 1.0):
                # Check if it's the first non-zero element in its row up to this point
                if all(are_almost_equal(rref_matrix[r][k], 0.0) for k in range(c)):
                    pivot_cols.append(c)
                    current_pivot_col = c
                    found_pivot_in_row = True
                    break
        if not found_pivot_in_row: # Si pas de pivot dans cette ligne, elle est nulle ou dependante
            continue
    return pivot_cols

def get_rank(rref_matrix: list[list[float]]) -> int:
    """Calcule le rang de la matrice (dimension de l'image)."""
    return len(get_pivot_columns(rref_matrix))

def get_basis_image(original_matrix: list[list[float]], rref_matrix: list[list[float]]) -> list[list[float]]:
    """
    Determine une base pour l'image (espace colonne) de la matrice.
    Les vecteurs de base sont les colonnes de la matrice originale correspondant aux colonnes pivots de la RREF.
    """
    if not original_matrix or not original_matrix[0]:
        return []
    pivot_cols = get_pivot_columns(rref_matrix)
    basis = [get_col(original_matrix, col_idx) for col_idx in pivot_cols]
    return basis

def get_basis_kernel(rref_matrix: list[list[float]]) -> list[list[float]]:
    """
    Determine une base pour le noyau (null space) de la matrice.
    Resout A @ x = 0 en exprimant les variables de base en fonction des variables libres.
    """
    if not rref_matrix or not rref_matrix[0]:
        return []

    num_rows = len(rref_matrix)
    num_cols = len(rref_matrix[0])

    pivot_cols = get_pivot_columns(rref_matrix)
    free_vars = [c for c in range(num_cols) if c not in pivot_cols]

    if not free_vars: # Si pas de variables libres, le noyau est reduit au vecteur nul $\{0\}$
        return []

    basis_kernel = []

    # Dictionnaire pour associer l'index de colonne pivot a l'index de ligne ou il est pivot
    pivot_col_to_row = {}
    current_pivot_col_idx = 0
    for r in range(num_rows):
        if current_pivot_col_idx < len(pivot_cols):
            col_idx = pivot_cols[current_pivot_col_idx]
            if are_almost_equal(rref_matrix[r][col_idx], 1.0) and all(are_almost_equal(rref_matrix[r][k], 0.0) for k in range(col_idx)):
                 pivot_col_to_row[col_idx] = r
                 current_pivot_col_idx += 1

    # Pour chaque variable libre, construire un vecteur de base
    for free_var_idx in free_vars:
        basis_vector = [0.0] * num_cols
        basis_vector[free_var_idx] = 1.0 # La variable libre courante est 1

        # Remplir les valeurs des variables de base
        # Pour chaque variable de base x_j (ou j est un pivot_col), trouver son equation
        # et exprimer x_j en fonction des variables libres.
        for basic_var_col_idx in pivot_cols:
            # Trouver la ligne ou basic_var_col_idx est un pivot
            pivot_r = pivot_col_to_row[basic_var_col_idx]

            # L'equation de cette ligne (dans RREF) est:
            # 1 * x_basic_var_col_idx + sum(coeff * x_free) = 0
            # Donc, x_basic_var_col_idx = -sum(coeff * x_free)

            sum_free_coeffs = 0.0
            for fv in free_vars:
                # Seule la variable libre courante (free_var_idx) est 1 dans basis_vector
                sum_free_coeffs += rref_matrix[pivot_r][fv] * basis_vector[fv] # basis_vector[fv] est 1 si fv==free_var_idx, 0 sinon

            basis_vector[basic_var_col_idx] = -sum_free_coeffs

        basis_kernel.append(basis_vector)

    return basis_kernel


def get_nullity(rref_matrix: list[list[float]]) -> int:
    """Calcule la nullite de la matrice (dimension du noyau)."""
    if not rref_matrix:
        return 0
    num_cols = len(rref_matrix[0]) if rref_matrix else 0
    return num_cols - get_rank(rref_matrix) # Theoreme du rang : nullite = dim_domaine - rang

# --- Tests et assertions de validation ---

print("--- Debut des tests ---")

# Matrice A pour les tests (3x4)
A = [
    [1.0, 2.0, 3.0, 4.0],
    [5.0, 6.0, 7.0, 8.0],
    [9.0, 10.0, 11.0, 12.0]
]

# Matrice B (plus simple, 3x3, rang 2)
B = [
    [1.0, 1.0, 2.0],
    [2.0, 2.0, 4.0],
    [0.0, 1.0, 1.0]
]

# Matrice C (3x3, rang 1)
C = [
    [1.0, 2.0, 1.0],
    [2.0, 4.0, 2.0],
    [3.0, 6.0, 3.0]
]

# Matrice D (identite, 3x3, rang 3)
D = [
    [1.0, 0.0, 0.0],
    [0.0, 1.0, 0.0],
    [0.0, 0.0, 1.0]
]

# Matrice E (nulle, 2x3, rang 0)
E = [
    [0.0, 0.0, 0.0],
    [0.0, 0.0, 0.0]
]


# --- Test des operations de base ---
print("\n# Tests des operations de base #")
v1 = [1.0, 2.0, 3.0]
v2 = [4.0, 5.0, 6.0]
assert vec_add(v1, v2) == [5.0, 7.0, 9.0]
assert scalar_vec_mul(2.0, v1) == [2.0, 4.0, 6.0]
assert dot_product(v1, v2) == 32.0 # 1*4 + 2*5 + 3*6 = 4 + 10 + 18 = 32
assert are_almost_equal(vec_norm(v1), math.sqrt(14.0)) # 1^2 + 2^2 + 3^2 = 1 + 4 + 9 = 14

mat_test_mv = [
    [1.0, 2.0],
    [3.0, 4.0]
]
vec_test_mv = [5.0, 6.0]
# [1*5 + 2*6, 3*5 + 4*6] = [5+12, 15+24] = [17, 39]
assert mat_vec_mul(mat_test_mv, vec_test_mv) == [17.0, 39.0]
print("Operations de base : OK")

# --- Test get_rref ---
print("\n# Tests de get_rref #")
rref_A = get_rref(A)
expected_rref_A = [
    [1.0, 0.0, -1.0, -2.0],
    [0.0, 1.0, 2.0, 3.0],
    [0.0, 0.0, 0.0, 0.0]
]
for r_idx in range(len(rref_A)):
    for c_idx in range(len(rref_A[0])):
        assert are_almost_equal(rref_A[r_idx][c_idx], expected_rref_A[r_idx][c_idx]), \
            f"RREF de A - Erreur a ({r_idx},{c_idx}): Attendu {expected_rref_A[r_idx][c_idx]}, Obtenu {rref_A[r_idx][c_idx]}"
print("RREF de A : OK")

rref_B = get_rref(B)
expected_rref_B = [
    [1.0, 0.0, 1.0],
    [0.0, 1.0, 1.0],
    [0.0, 0.0, 0.0]
]
for r_idx in range(len(rref_B)):
    for c_idx in range(len(rref_B[0])):
        assert are_almost_equal(rref_B[r_idx][c_idx], expected_rref_B[r_idx][c_idx]), \
            f"RREF de B - Erreur a ({r_idx},{c_idx}): Attendu {expected_rref_B[r_idx][c_idx]}, Obtenu {rref_B[r_idx][c_idx]}"
print("RREF de B : OK")

rref_C = get_rref(C)
expected_rref_C = [
    [1.0, 2.0, 1.0],
    [0.0, 0.0, 0.0],
    [0.0, 0.0, 0.0]
]
for r_idx in range(len(rref_C)):
    for c_idx in range(len(rref_C[0])):
        assert are_almost_equal(rref_C[r_idx][c_idx], expected_rref_C[r_idx][c_idx]), \
            f"RREF de C - Erreur a ({r_idx},{c_idx}): Attendu {expected_rref_C[r_idx][c_idx]}, Obtenu {rref_C[r_idx][c_idx]}"
print("RREF de C : OK")

rref_D = get_rref(D)
expected_rref_D = [
    [1.0, 0.0, 0.0],
    [0.0, 1.0, 0.0],
    [0.0, 0.0, 1.0]
]
for r_idx in range(len(rref_D)):
    for c_idx in range(len(rref_D[0])):
        assert are_almost_equal(rref_D[r_idx][c_idx], expected_rref_D[r_idx][c_idx]), \
            f"RREF de D - Erreur a ({r_idx},{c_idx}): Attendu {expected_rref_D[r_idx][c_idx]}, Obtenu {rref_D[r_idx][c_idx]}"
print("RREF de D : OK")

rref_E = get_rref(E)
expected_rref_E = [
    [0.0, 0.0, 0.0],
    [0.0, 0.0, 0.0]
]
for r_idx in range(len(rref_E)):
    for c_idx in range(len(rref_E[0])):
        assert are_almost_equal(rref_E[r_idx][c_idx], expected_rref_E[r_idx][c_idx]), \
            f"RREF de E - Erreur a ({r_idx},{c_idx}): Attendu {expected_rref_E[r_idx][c_idx]}, Obtenu {rref_E[r_idx][c_idx]}"
print("RREF de E : OK")

# --- Test get_pivot_columns et get_rank ---
print("\n# Tests de get_pivot_columns et get_rank #")
assert get_pivot_columns(rref_A) == [0, 1]
assert get_rank(rref_A) == 2
print("Pivot columns et rang de A : OK")

assert get_pivot_columns(rref_B) == [0, 1]
assert get_rank(rref_B) == 2
print("Pivot columns et rang de B : OK")

assert get_pivot_columns(rref_C) == [0]
assert get_rank(rref_C) == 1
print("Pivot columns et rang de C : OK")

assert get_pivot_columns(rref_D) == [0, 1, 2]
assert get_rank(rref_D) == 3
print("Pivot columns et rang de D : OK")

assert get_pivot_columns(rref_E) == []
assert get_rank(rref_E) == 0
print("Pivot columns et rang de E : OK")

# --- Test get_basis_image ---
print("\n# Tests de get_basis_image #")
basis_image_A = get_basis_image(A, rref_A)
# Pivot columns for A are 0 and 1. So basis for image should be A[:,0] and A[:,1]
expected_basis_image_A = [[1.0, 5.0, 9.0], [2.0, 6.0, 10.0]]
assert len(basis_image_A) == len(expected_basis_image_A)
for i in range(len(basis_image_A)):
    for j in range(len(basis_image_A[0])):
        assert are_almost_equal(basis_image_A[i][j], expected_basis_image_A[i][j]), \
            f"Base image de A - Erreur a ({i},{j}): Attendu {expected_basis_image_A[i][j]}, Obtenu {basis_image_A[i][j]}"
print("Base de l'image de A : OK")

basis_image_B = get_basis_image(B, rref_B)
# Pivot columns for B are 0 and 1.
expected_basis_image_B = [[1.0, 2.0, 0.0], [1.0, 2.0, 1.0]]
assert len(basis_image_B) == len(expected_basis_image_B)
for i in range(len(basis_image_B)):
    for j in range(len(basis_image_B[0])):
        assert are_almost_equal(basis_image_B[i][j], expected_basis_image_B[i][j]), \
            f"Base image de B - Erreur a ({i},{j}): Attendu {expected_basis_image_B[i][j]}, Obtenu {basis_image_B[i][j]}"
print("Base de l'image de B : OK")

basis_image_C = get_basis_image(C, rref_C)
# Pivot column for C is 0.
expected_basis_image_C = [[1.0, 2.0, 3.0]]
assert len(basis_image_C) == len(expected_basis_image_C)
for i in range(len(basis_image_C)):
    for j in range(len(basis_image_C[0])):
        assert are_almost_equal(basis_image_C[i][j], expected_basis_image_C[i][j]), \
            f"Base image de C - Erreur a ({i},{j}): Attendu {expected_basis_image_C[i][j]}, Obtenu {basis_image_C[i][j]}"
print("Base de l'image de C : OK")

# --- Test get_basis_kernel et get_nullity ---
print("\n# Tests de get_basis_kernel et get_nullity #")
basis_kernel_A = get_basis_kernel(rref_A)
# rref_A:
# [1.0, 0.0, -1.0, -2.0]
# [0.0, 1.0, 2.0, 3.0]
# [0.0, 0.0, 0.0, 0.0]
# Pivot vars: x0, x1. Free vars: x2, x3.
# x0 - x2 - 2x3 = 0  => x0 = x2 + 2x3
# x1 + 2x2 + 3x3 = 0 => x1 = -2x2 - 3x3
# For x2=1, x3=0: x = [1, -2, 1, 0]
# For x2=0, x3=1: x = [2, -3, 0, 1]
expected_basis_kernel_A = [[1.0, -2.0, 1.0, 0.0], [2.0, -3.0, 0.0, 1.0]]
assert len(basis_kernel_A) == len(expected_basis_kernel_A)
# Sort to make comparison robust to order of free vars iteration
basis_kernel_A.sort(key=lambda v: [round(x, 5) for x in v])
expected_basis_kernel_A.sort(key=lambda v: [round(x, 5) for x in v])
for i in range(len(basis_kernel_A)):
    for j in range(len(basis_kernel_A[0])):
        assert are_almost_equal(basis_kernel_A[i][j], expected_basis_kernel_A[i][j]), \
            f"Base noyau de A - Erreur a ({i},{j}): Attendu {expected_basis_kernel_A[i][j]}, Obtenu {basis_kernel_A[i][j]}"
print("Base du noyau de A : OK")
assert get_nullity(rref_A) == 2
print("Nullite de A : OK")

# Verifier que les vecteurs du noyau sont bien annules par la matrice originale A
for vec_k in basis_kernel_A:
    result = mat_vec_mul(A, vec_k)
    assert is_zero_vector(result), f"Vecteur du noyau de A non annule: {vec_k}, Resultat: {result}"
print("Verification des vecteurs du noyau de A : OK")


basis_kernel_B = get_basis_kernel(rref_B)
# rref_B:
# [1.0, 0.0, 1.0]
# [0.0, 1.0, 1.0]
# [0.0, 0.0, 0.0]
# Pivot vars: x0, x1. Free var: x2.
# x0 + x2 = 0 => x0 = -x2
# x1 + x2 = 0 => x1 = -x2
# For x2=1: x = [-1, -1, 1]
expected_basis_kernel_B = [[-1.0, -1.0, 1.0]]
assert len(basis_kernel_B) == len(expected_basis_kernel_B)
for i in range(len(basis_kernel_B)):
    for j in range(len(basis_kernel_B[0])):
        assert are_almost_equal(basis_kernel_B[i][j], expected_basis_kernel_B[i][j]), \
            f"Base noyau de B - Erreur a ({i},{j}): Attendu {expected_basis_kernel_B[i][j]}, Obtenu {basis_kernel_B[i][j]}"
print("Base du noyau de B : OK")
assert get_nullity(rref_B) == 1
print("Nullite de B : OK")
for vec_k in basis_kernel_B:
    result = mat_vec_mul(B, vec_k)
    assert is_zero_vector(result), f"Vecteur du noyau de B non annule: {vec_k}, Resultat: {result}"
print("Verification des vecteurs du noyau de B : OK")


basis_kernel_C = get_basis_kernel(rref_C)
# rref_C:
# [1.0, 2.0, 1.0]
# [0.0, 0.0, 0.0]
# [0.0, 0.0, 0.0]
# Pivot var: x0. Free vars: x1, x2.
# x0 + 2x1 + x2 = 0 => x0 = -2x1 - x2
# For x1=1, x2=0: x = [-2, 1, 0]
# For x1=0, x2=1: x = [-1, 0, 1]
expected_basis_kernel_C = [[-2.0, 1.0, 0.0], [-1.0, 0.0, 1.0]]
assert len(basis_kernel_C) == len(expected_basis_kernel_C)
basis_kernel_C.sort(key=lambda v: [round(x, 5) for x in v])
expected_basis_kernel_C.sort(key=lambda v: [round(x, 5) for x in v])
for i in range(len(basis_kernel_C)):
    for j in range(len(basis_kernel_C[0])):
        assert are_almost_equal(basis_kernel_C[i][j], expected_basis_kernel_C[i][j]), \
            f"Base noyau de C - Erreur a ({i},{j}): Attendu {expected_basis_kernel_C[i][j]}, Obtenu {basis_kernel_C[i][j]}"
print("Base du noyau de C : OK")
assert get_nullity(rref_C) == 2
print("Nullite de C : OK")
for vec_k in basis_kernel_C:
    result = mat_vec_mul(C, vec_k)
    assert is_zero_vector(result), f"Vecteur du noyau de C non annule: {vec_k}, Resultat: {result}"
print("Verification des vecteurs du noyau de C : OK")


basis_kernel_D = get_basis_kernel(rref_D)
assert basis_kernel_D == [] # Matrice identite, noyau reduit au vecteur nul
print("Base du noyau de D : OK")
assert get_nullity(rref_D) == 0
print("Nullite de D : OK")


basis_kernel_E = get_basis_kernel(rref_E)
# rref_E:
# [0.0, 0.0, 0.0]
# [0.0, 0.0, 0.0]
# No pivot vars. Free vars: x0, x1, x2.
# For x0=1, x1=0, x2=0: x = [1, 0, 0]
# For x0=0, x1=1, x2=0: x = [0, 1, 0]
# For x0=0, x1=0, x2=1: x = [0, 0, 1]
expected_basis_kernel_E = [[1.0, 0.0, 0.0], [0.0, 1.0, 0.0], [0.0, 0.0, 1.0]]
assert len(basis_kernel_E) == len(expected_basis_kernel_E)
basis_kernel_E.sort(key=lambda v: [round(x, 5) for x in v])
expected_basis_kernel_E.sort(key=lambda v: [round(x, 5) for x in v])
for i in range(len(basis_kernel_E)):
    for j in range(len(basis_kernel_E[0])):
        assert are_almost_equal(basis_kernel_E[i][j], expected_basis_kernel_E[i][j]), \
            f"Base noyau de E - Erreur a ({i},{j}): Attendu {expected_basis_kernel_E[i][j]}, Obtenu {basis_kernel_E[i][j]}"
print("Base du noyau de E : OK")
assert get_nullity(rref_E) == 3
print("Nullite de E : OK")
for vec_k in basis_kernel_E:
    result = mat_vec_mul(E, vec_k)
    assert is_zero_vector(result), f"Vecteur du noyau de E non annule: {vec_k}, Resultat: {result}"
print("Verification des vecteurs du noyau de E : OK")


# --- Test du theoreme du rang ---
print("\n# Tests du Theoreme du Rang #")
# dim_domain = nombre de colonnes de la matrice originale
# rang + nullite = dim_domain

# Matrice A (3x4)
dim_domain_A = len(A[0])
rank_A = get_rank(rref_A)
nullity_A = get_nullity(rref_A)
assert rank_A + nullity_A == dim_domain_A, f"Theoreme du rang echoue pour A: {rank_A} + {nullity_A} != {dim_domain_A}"
print(f"Theoreme du rang pour A (rang={rank_A}, nullite={nullity_A}, dim_domaine={dim_domain_A}): OK")

# Matrice B (3x3)
dim_domain_B = len(B[0])
rank_B = get_rank(rref_B)
nullity_B = get_nullity(rref_B)
assert rank_B + nullity_B == dim_domain_B, f"Theoreme du rang echoue pour B: {rank_B} + {nullity_B} != {dim_domain_B}"
print(f"Theoreme du rang pour B (rang={rank_B}, nullite={nullity_B}, dim_domaine={dim_domain_B}): OK")

# Matrice C (3x3)
dim_domain_C = len(C[0])
rank_C = get_rank(rref_C)
nullity_C = get_nullity(rref_C)
assert rank_C + nullity_C == dim_domain_C, f"Theoreme du rang echoue pour C: {rank_C} + {nullity_C} != {dim_domain_C}"
print(f"Theoreme du rang pour C (rang={rank_C}, nullite={nullity_C}, dim_domaine={dim_domain_C}): OK")

# Matrice D (3x3)
dim_domain_D = len(D[0])
rank_D = get_rank(rref_D)
nullity_D = get_nullity(rref_D)
assert rank_D + nullity_D == dim_domain_D, f"Theoreme du rang echoue pour D: {rank_D} + {nullity_D} != {dim_domain_D}"
print(f"Theoreme du rang pour D (rang={rank_D}, nullite={nullity_D}, dim_domaine={dim_domain_D}): OK")

# Matrice E (2x3)
dim_domain_E = len(E[0])
rank_E = get_rank(rref_E)
nullity_E = get_nullity(rref_E)
assert rank_E + nullity_E == dim_domain_E, f"Theoreme du rang echoue pour E: {rank_E} + {nullity_E} != {dim_domain_E}"
print(f"Theoreme du rang pour E (rang={rank_E}, nullite={nullity_E}, dim_domaine={dim_domain_E}): OK")


print("\n--- Tous les tests sont passes avec succes ! ---")
\end{verbatim}

\newpage
\chapter*{TP 3 : Transformations Linéaires, Noyau et Image "From Scratch" en Python}

\section*{Objectifs du TP}

Ce troisième TP vise à implémenter des concepts fondamentaux de l'algèbre linéaire, essentiels à la compréhension et à la conception de nombreux algorithmes d'Intelligence Artificielle. L'implémentation sera réalisée \textbf{entièrement "from scratch" en Python pur}, sans l'aide de librairies de calcul scientifique de haut niveau (comme NumPy, SciPy, etc.) pour les opérations d'algèbre linéaire.

\subsection*{Objectifs Mathématiques}
*   \textbf{Comprendre la représentation numérique} des vecteurs et matrices.
*   \textbf{Implémenter les opérations de base} sur les vecteurs et matrices (addition, multiplication scalaire, produit scalaire, multiplication matrice-vecteur, multiplication matrice-matrice).
*   \textbf{Maîtriser l'algorithme d'élimination de Gauss-Jordan} pour réduire une matrice en forme échelonnée réduite (RREF).
*   \textbf{Calculer le rang} d'une matrice.
*   \textbf{Déterminer une base du noyau (espace nul)} d'une application linéaire.
*   \textbf{Déterminer une base de l'image (espace colonne)} d'une application linéaire.
*   \textbf{Vérifier le Théorème du Rang} (\_\_INLINE\_CODE\_0\_\_).

\subsection*{Objectifs IA}
*   \textbf{Fondations des Réseaux de Neurones :} Comprendre comment les multiplications matrice-vecteur et matrice-matrice sont au cœur des couches denses (fully connected layers) sans fonction d'activation.
*   \textbf{Réduction de Dimensionalité et Compression :} Le noyau et l'image d'une matrice linéaire sont des outils clés pour comprendre les transformations de l'espace des caractéristiques, la perte ou la conservation d'information, et les principes derrière des techniques comme l'Analyse en Composantes Principales (ACP/PCA).
*   \textbf{Interprétabilité et Robustesse des Modèles :} L'analyse du noyau et de l'image permet d'identifier quelles combinaisons d'entrées produisent des sorties nulles ou quelles sont les sorties possibles d'un modèle, ce qui est crucial pour l'interprétabilité et la robustesse des systèmes d'IA.
*   \textbf{Algorithmes d'Optimisation :} La résolution de systèmes linéaires (via l'élimination de Gauss) est une brique de base pour de nombreux problèmes d'optimisation en ML.

---

\section*{Implémentation Python "From Scratch"}

Nous allons représenter les vecteurs comme des listes de nombres (\_\_INLINE\_CODE\_1\_\_) et les matrices comme des listes de listes de nombres (\_\_INLINE\_CODE\_2\_\_).

\begin{verbatim}
import math

# Constante pour gerer la precision des calculs flottants
EPSILON = 1e-9

def is_zero(val):
    """Verifie si une valeur est consideree comme nulle."""
    return abs(val) < EPSILON

def is_zero_vector(v):
    """Verifie si tous les elements d'un vecteur sont consideres comme nuls."""
    return all(is_zero(x) for x in v)

def are_vectors_equal(v1, v2):
    """Verifie si deux vecteurs sont egaux avec une tolerance EPSILON."""
    if len(v1) != len(v2):
        return False
    return all(is_zero(v1[i] - v2[i]) for i in range(len(v1)))

def clone_matrix(matrix):
    """Cree une copie profonde d'une matrice."""
    return [row[:] for row in matrix]

def print_matrix(matrix, name="Matrix"):
    """Affiche une matrice de maniere formatee."""
    print(f"{name}:")
    for row in matrix:
        print([f"{x:.4f}" for x in row])
    print("-" * 30)

# --- 1. Operations de base sur les vecteurs et matrices ---

def vector_add(v1, v2):
    """Additionne deux vecteurs."""
    assert len(v1) == len(v2), "Les vecteurs doivent avoir la meme dimension."
    return [v1[i] + v2[i] for i in range(len(v1))]

def scalar_vector_mul(s, v):
    """Multiplie un vecteur par un scalaire."""
    return [s * x for x in v]

def dot_product(v1, v2):
    """Calcule le produit scalaire de deux vecteurs."""
    assert len(v1) == len(v2), "Les vecteurs doivent avoir la meme dimension pour le produit scalaire."
    return sum(v1[i] * v2[i] for i in range(len(v1)))

def matrix_vector_mul(matrix, vector):
    """Multiplie une matrice par un vecteur."""
    num_rows = len(matrix)
    num_cols = len(matrix[0]) if num_rows > 0 else 0
    assert num_cols == len(vector), "Le nombre de colonnes de la matrice doit correspondre a la dimension du vecteur."

    result = [0.0] * num_rows
    for i in range(num_rows):
        result[i] = dot_product(matrix[i], vector)
    return result

def matrix_mul(m1, m2):
    """Multiplie deux matrices."""
    rows1, cols1 = len(m1), len(m1[0]) if len(m1) > 0 else 0
    rows2, cols2 = len(m2), len(m2[0]) if len(m2) > 0 else 0

    assert cols1 == rows2, "Le nombre de colonnes de la premiere matrice doit correspondre au nombre de lignes de la seconde."

    result = [[0.0 for _ in range(cols2)] for _ in range(rows1)]
    for i in range(rows1):
        for j in range(cols2):
            for k in range(cols1):
                result[i][j] += m1[i][k] * m2[k][j]
    return result

def transpose(matrix):
    """Transposer une matrice."""
    rows, cols = len(matrix), len(matrix[0]) if len(matrix) > 0 else 0
    if rows == 0 or cols == 0:
        return []
    transposed = [[0.0 for _ in range(rows)] for _ in range(cols)]
    for i in range(rows):
        for j in range(cols):
            transposed[j][i] = matrix[i][j]
    return transposed

# --- 2. Elimination de Gauss-Jordan pour la Forme Echelonnee Reduite (RREF) ---

def gaussian_elimination(matrix):
    """
    Transforme une matrice en sa Forme Echelonnee Reduite (RREF)
    en utilisant l'elimination de Gauss-Jordan.
    Retourne la matrice en RREF et une liste des indices des colonnes pivots.
    """
    mat = clone_matrix(matrix)
    rows = len(mat)
    cols = len(mat[0]) if rows > 0 else 0

    if rows == 0 or cols == 0:
        return [], []

    pivot_cols_indices = []
    lead = 0  # Indice de la colonne pivot actuelle

    for r in range(rows):
        if lead >= cols:
            break

        i = r
        # Trouver une ligne non nulle pour le pivot dans la colonne 'lead'
        while i < rows and is_zero(mat[i][lead]):
            i += 1

        if i < rows: # Si un pivot est trouve
            # Echanger la ligne courante (r) avec la ligne du pivot (i)
            mat[r], mat[i] = mat[i], mat[r]

            # Normaliser la ligne pivot pour que l'element pivot soit 1
            lv = mat[r][lead]
            for j in range(cols):
                mat[r][j] /= lv

            # Eliminer les autres elements de la colonne 'lead' dans les autres lignes
            for i_other_row in range(rows):
                if i_other_row != r:
                    lv_other_row = mat[i_other_row][lead]
                    for j in range(cols):
                        mat[i_other_row][j] -= lv_other_row * mat[r][j]

            pivot_cols_indices.append(lead)
            lead += 1 # Passer a la colonne pivot suivante
        else: # Si aucun pivot n'est trouve dans cette colonne, passer a la colonne suivante
            lead += 1
            # Redo the current row, as it was skipped.
            # This is important for sparse matrices where 'lead' might advance several times without finding a pivot.
            # Example: [[0,1],[0,0]] -> lead should jump to col 1.
            # The current loop structure `for r in range(rows)` naturally tries to find a pivot for row `r`.
            # If no pivot is found in `lead` for row `r`, we advance `lead` and will try to find a pivot
            # for row `r` in the *next* `lead` column. So `r` shouldn't increment.
            # This is done by decrementing `r` here.
            r -= 1 # Re-evaluate the current row `r` with the new `lead`

    # Arrondir les valeurs tres proches de zero a zero absolu
    for r in range(rows):
        for c in range(cols):
            if is_zero(mat[r][c]):
                mat[r][c] = 0.0

    return mat, pivot_cols_indices

# --- 3. Fonctions d'algebre lineaire avancees ---

def get_rank(matrix):
    """
    Calcule le rang d'une matrice (la dimension de son espace colonne).
    Le rang est le nombre de pivots dans la forme echelonnee reduite.
    """
    if not matrix or not matrix[0]:
        return 0
    _, pivot_cols_indices = gaussian_elimination(matrix)
    return len(pivot_cols_indices)

def get_kernel_basis(matrix):
    """
    Trouve une base pour le noyau (espace nul) d'une matrice.
    """
    m, n = len(matrix), len(matrix[0]) if len(matrix) > 0 else 0
    if n == 0: return []
    if m == 0: return [[1.0 if i == k else 0.0 for i in range(n)] for k in range(n)] # Base case for 0-row matrix, kernel is R^n

    rref_mat, pivot_cols_indices = gaussian_elimination(matrix)

    pivot_vars_set = set(pivot_cols_indices)
    free_vars = sorted(list(set(range(n)) - pivot_vars_set))

    kernel_basis = []

    # Si toutes les colonnes sont des colonnes pivots, le noyau est reduit au vecteur nul {0}
    if not free_vars:
        return []

    # Pour chaque variable libre, construire un vecteur de base
    for free_idx in free_vars:
        sol_vec = [0.0] * n
        sol_vec[free_idx] = 1.0 # Definir cette variable libre a 1

        # Utiliser la RREF pour exprimer les variables pivots en fonction des variables libres
        # On itere sur les lignes qui contiennent un pivot
        for r_pivot_row in range(len(pivot_cols_indices)):
            pivot_col = pivot_cols_indices[r_pivot_row]

            # La relation dans la RREF est: x_pivot_col + sum(c_j * x_j) = 0
            # ou x_j sont d'autres variables, et c_j sont les coefficients de rref_mat.
            # Puisque nous avons fixe les variables libres (sol_vec[k] est 0 ou 1)
            # et que les autres variables pivots sont 0 dans la ligne du pivot,
            # cela simplifie a: x_pivot_col = - (sum(rref_mat[r_pivot_row][k] * sol_vec[k]) pour k != pivot_col)

            val_for_pivot = 0.0
            for k in range(n):
                if k != pivot_col:
                    val_for_pivot += rref_mat[r_pivot_row][k] * sol_vec[k]
            sol_vec[pivot_col] = -val_for_pivot

        # Nettoyer les petites valeurs flottantes
        for i in range(n):
            if is_zero(sol_vec[i]):
                sol_vec[i] = 0.0

        kernel_basis.append(sol_vec)

    return kernel_basis

def get_image_basis(original_matrix):
    """
    Trouve une base pour l'image (espace colonne) d'une matrice.
    La base est constituee des colonnes originales de la matrice
    qui correspondent aux colonnes pivots de sa forme echelonnee reduite.
    """
    m, n = len(original_matrix), len(original_matrix[0]) if len(original_matrix) > 0 else 0
    if m == 0 or n == 0:
        return []

    _, pivot_cols_indices = gaussian_elimination(original_matrix)

    image_basis = []
    # Collecter les colonnes originales correspondant aux indices des colonnes pivots
    for col_idx in pivot_cols_indices:
        column_vector = [original_matrix[row][col_idx] for row in range(m)]
        image_basis.append(column_vector)

    return image_basis

# --- 4. Verification du Theoreme du Rang ---

def verify_rank_nullity_theorem(matrix):
    """
    Verifie le Theoreme du Rang pour une matrice donnee :
    dim(Ker(A)) + dim(Im(A)) = dim(domaine)
    ou dim(domaine) est le nombre de colonnes de A.
    Retourne True si le theoreme est verifie, False sinon.
    """
    m, n = len(matrix), len(matrix[0]) if len(matrix) > 0 else 0
    if n == 0: # If no columns, domain dimension is 0. Rank=0, Nullity=0. 0+0=0.
        return True

    rank_A = get_rank(matrix)
    kernel_basis = get_kernel_basis(matrix)
    nullity_A = len(kernel_basis) # Dimension du noyau

    # Le theoreme est rank + nullity = nombre de colonnes (dimension du domaine)
    return abs((rank_A + nullity_A) - n) < EPSILON

# --- 5. Cas de test et Assertions de validation ---

print("--- Debut des tests ---")

# --- Test 1: Matrice carree de plein rang ---
M1 = [
    [1.0, 2.0],
    [3.0, 4.0]
]
print_matrix(M1, "M1")
assert get_rank(M1) == 2, "Test 1: Le rang de M1 devrait etre 2."
kernel_M1 = get_kernel_basis(M1)
assert len(kernel_M1) == 0, "Test 1: Le noyau de M1 devrait etre reduit au vecteur nul (dimension 0)."
image_M1 = get_image_basis(M1)
assert len(image_M1) == 2, "Test 1: L'image de M1 devrait etre de dimension 2."
assert are_vectors_equal(image_M1[0], [1.0, 3.0])
assert are_vectors_equal(image_M1[1], [2.0, 4.0])
assert verify_rank_nullity_theorem(M1), "Test 1: Le Theoreme du Rang devrait etre verifie pour M1."

# --- Test 2: Matrice rectangulaire avec dependance lineaire ---
M2 = [
    [1.0, 2.0, 3.0],
    [2.0, 4.0, 6.0]
]
print_matrix(M2, "M2")
assert get_rank(M2) == 1, "Test 2: Le rang de M2 devrait etre 1."
kernel_M2 = get_kernel_basis(M2)
assert len(kernel_M2) == 2, "Test 2: Le noyau de M2 devrait etre de dimension 2."
# Verifier que les vecteurs du noyau sont corrects (Ax=0)
assert is_zero_vector(matrix_vector_mul(M2, kernel_M2[0])), "Test 2: Le premier vecteur du noyau ne produit pas le vecteur nul."
assert is_zero_vector(matrix_vector_mul(M2, kernel_M2[1])), "Test 2: Le second vecteur du noyau ne produit pas le vecteur nul."
# Exemple de base attendue pour M2: [[-2, 1, 0], [-3, 0, 1]] ou des multiples
assert are_vectors_equal(kernel_M2[0], [-2.0, 1.0, 0.0]) or are_vectors_equal(kernel_M2[0], [-3.0, 0.0, 1.0])
assert are_vectors_equal(kernel_M2[1], [-3.0, 0.0, 1.0]) or are_vectors_equal(kernel_M2[1], [-2.0, 1.0, 0.0])

image_M2 = get_image_basis(M2)
assert len(image_M2) == 1, "Test 2: L'image de M2 devrait etre de dimension 1."
assert are_vectors_equal(image_M2[0], [1.0, 2.0]), "Test 2: La base de l'image de M2 est incorrecte."
assert verify_rank_nullity_theorem(M2), "Test 2: Le Theoreme du Rang devrait etre verifie pour M2."

# --- Test 3: Matrice identite ---
M3 = [
    [1.0, 0.0, 0.0],
    [0.0, 1.0, 0.0],
    [0.0, 0.0, 1.0]
]
print_matrix(M3, "M3")
assert get_rank(M3) == 3, "Test 3: Le rang de M3 devrait etre 3."
kernel_M3 = get_kernel_basis(M3)
assert len(kernel_M3) == 0, "Test 3: Le noyau de M3 devrait etre reduit au vecteur nul."
image_M3 = get_image_basis(M3)
assert len(image_M3) == 3, "Test 3: L'image de M3 devrait etre de dimension 3."
assert verify_rank_nullity_theorem(M3), "Test 3: Le Theoreme du Rang devrait etre verifie pour M3."

# --- Test 4: Matrice de zeros ---
M4 = [
    [0.0, 0.0, 0.0],
    [0.0, 0.0, 0.0]
]
print_matrix(M4, "M4")
assert get_rank(M4) == 0, "Test 4: Le rang de M4 devrait etre 0."
kernel_M4 = get_kernel_basis(M4)
assert len(kernel_M4) == 3, "Test 4: Le noyau de M4 devrait etre de dimension 3 (R^3)."
assert is_zero_vector(matrix_vector_mul(M4, kernel_M4[0]))
assert is_zero_vector(matrix_vector_mul(M4, kernel_M4[1]))
assert is_zero_vector(matrix_vector_mul(M4, kernel_M4[2]))
image_M4 = get_image_basis(M4)
assert len(image_M4) == 0, "Test 4: L'image de M4 devrait etre de dimension 0 (espace nul)."
assert verify_rank_nullity_theorem(M4), "Test 4: Le Theoreme du Rang devrait etre verifie pour M4."

# --- Test 5: Matrice colonne ---
M5 = [
    [1.0],
    [2.0],
    [3.0]
]
print_matrix(M5, "M5")
assert get_rank(M5) == 1, "Test 5: Le rang de M5 devrait etre 1."
kernel_M5 = get_kernel_basis(M5)
assert len(kernel_M5) == 0, "Test 5: Le noyau de M5 devrait etre reduit au vecteur nul."
image_M5 = get_image_basis(M5)
assert len(image_M5) == 1, "Test 5: L'image de M5 devrait etre de dimension 1."
assert are_vectors_equal(image_M5[0], [1.0, 2.0, 3.0])
assert verify_rank_nullity_theorem(M5), "Test 5: Le Theoreme du Rang devrait etre verifie pour M5."

# --- Test 6: Matrice ligne ---
M6 = [[1.0, 2.0, 3.0]]
print_matrix(M6, "M6")
assert get_rank(M6) == 1, "Test 6: Le rang de M6 devrait etre 1."
kernel_M6 = get_kernel_basis(M6)
assert len(kernel_M6) == 2, "Test 6: Le noyau de M6 devrait etre de dimension 2."
assert is_zero_vector(matrix_vector_mul(M6, kernel_M6[0]))
assert is_zero_vector(matrix_vector_mul(M6, kernel_M6[1]))
image_M6 = get_image_basis(M6)
assert len(image_M6) == 1, "Test 6: L'image de M6 devrait etre de dimension 1."
assert are_vectors_equal(image_M6[0], [1.0]) # La base de l'image est la premiere colonne originale qui est pivot
assert verify_rank_nullity_theorem(M6), "Test 6: Le Theoreme du Rang devrait etre verifie pour M6."

# --- Test 7: Matrice non carree avec plus de lignes que de colonnes ---
M7 = [
    [1.0, 0.0],
    [0.0, 1.0],
    [0.0, 0.0]
]
print_matrix(M7, "M7")
assert get_rank(M7) == 2, "Test 7: Le rang de M7 devrait etre 2."
kernel_M7 = get_kernel_basis(M7)
assert len(kernel_M7) == 0, "Test 7: Le noyau de M7 devrait etre de dimension 0."
image_M7 = get_image_basis(M7)
assert len(image_M7) == 2, "Test 7: L'image de M7 devrait etre de dimension 2."
assert are_vectors_equal(image_M7[0], [1.0, 0.0, 0.0])
assert are_vectors_equal(image_M7[1], [0.0, 1.0, 0.0])
assert verify_rank_nullity_theorem(M7), "Test 7: Le Theoreme du Rang devrait etre verifie pour M7."

# --- Test 8: Matrice avec des fractions et rang plein ---
M8 = [
    [2.0, 1.0, 3.0, 0.0],
    [4.0, 2.0, 6.0, 0.0],
    [1.0, 0.5, 1.5, 0.0]
]
print_matrix(M8, "M8")
assert get_rank(M8) == 1, "Test 8: Le rang de M8 devrait etre 1."
kernel_M8 = get_kernel_basis(M8)
assert len(kernel_M8) == 3, "Test 8: Le noyau de M8 devrait etre de dimension 3."
for vec in kernel_M8:
    assert is_zero_vector(matrix_vector_mul(M8, vec)), f"Test 8: Le vecteur {vec} du noyau ne produit pas le vecteur nul."
assert verify_rank_nullity_theorem(M8), "Test 8: Le Theoreme du Rang devrait etre verifie pour M8."

# --- Test 9: Matrice vide ou 0xN/Nx0 ---
M9_empty_rows = []
assert get_rank(M9_empty_rows) == 0, "Test 9: Le rang d'une matrice vide (0 lignes) doit etre 0."
assert len(get_kernel_basis(M9_empty_rows)) == 0 # depends on interpretation of N
assert len(get_image_basis(M9_empty_rows)) == 0

M9_empty_cols = [[],[]]
assert get_rank(M9_empty_cols) == 0, "Test 9: Le rang d'une matrice vide (0 colonnes) doit etre 0."
assert len(get_kernel_basis(M9_empty_cols)) == 0 # Kernel is always empty if no columns
assert len(get_image_basis(M9_empty_cols)) == 0


print("\n--- Tous les tests ont ete passes avec succes ! ---")
\end{verbatim}

\newpage
\chapter*{TP 4 : Applications Linéaires, Noyau et Image pour l'Analyse de Caractéristiques}

\section*{Objectif Mathématique et IA du TP}

Ce TP a pour objectif de vous familiariser avec les concepts fondamentaux d'algèbre linéaire que sont les applications linéaires, leur noyau (kernel) et leur image (image, ou espace colonne). Ces notions sont au cœur de nombreuses techniques d'apprentissage automatique, notamment pour la transformation de caractéristiques, la réduction de dimensionnalité, l'analyse des espaces de données et la compréhension des modèles.

\textbf{Objectifs Mathématiques :}
1.  \textbf{Implémenter des opérations matricielles de base} : Multiplication matrice-vecteur et manipulation de matrices.
2.  \textbf{Comprendre et implémenter la forme échelonnée réduite (RREF)} : Une étape cruciale pour l'analyse des matrices.
3.  \textbf{Calculer l'image d'un vecteur par une application linéaire} : Représenté par la multiplication d'une matrice par ce vecteur.
4.  \textbf{Déterminer une base du noyau d'une application linéaire} : L'ensemble des vecteurs qui sont projetés sur le vecteur nul.
5.  \textbf{Déterminer une base de l'image d'une application linéaire} : L'ensemble de tous les vecteurs qui peuvent être atteints par la transformation.
6.  \textbf{Vérifier le théorème du rang} : Comprendre la relation entre la dimension du noyau (nullité) et la dimension de l'image (rang) d'une matrice.

\textbf{Objectifs pour l'IA :}
1.  \textbf{Transformation de Caractéristiques (Feature Engineering)} : Les applications linéaires sont utilisées pour projeter des données d'un espace de caractéristiques vers un autre, par exemple pour créer de nouvelles caractéristiques ou pour adapter les données à un modèle.
2.  \textbf{Filtrage et Invariance (Noyau)} : Le noyau peut être interprété comme l'ensemble des "modifications" ou des "variations" dans les données qui sont ignorées ou "filtrées" par la transformation. Par exemple, si une application linéaire représente une compression d'image, son noyau contiendrait les informations d'image qui sont perdues (non essentielles) lors de la compression. Pour l'apprentissage profond, le noyau des couches linéaires peut révéler des invariances apprises.
3.  \textbf{Réduction de Dimensionalité et Représentation (Image)} : L'image d'une application linéaire définit le sous-espace dans lequel les données transformées résident. C'est la base de techniques comme la PCA (Analyse en Composantes Principales), où l'on cherche une projection linéaire des données sur un sous-espace de dimension inférieure qui capture le maximum de variance. Le rang de la matrice donne la dimension de l'espace de caractéristiques transformées atteignables.
4.  \textbf{Comprendre les Limites des Modèles Linéaires} : L'étude du noyau et de l'image aide à comprendre ce qu'un modèle linéaire peut et ne peut pas représenter ou distinguer.

---

\section*{Implémentation en Python Pur}

Toutes les opérations d'algèbre linéaire doivent être implémentées "from scratch", sans l'aide de librairies comme \_\_INLINE\_CODE\_0\_\_. Nous utiliserons des listes de listes pour représenter les matrices et des listes pour les vecteurs.

\begin{verbatim}
import math

# Constante pour gerer les erreurs de virgule flottante
EPSILON = 1e-9

def are_approximately_equal(a, b, epsilon=EPSILON):
    """Verifie si deux nombres flottants sont approximativement egaux."""
    return abs(a - b) < epsilon

def is_zero(val, epsilon=EPSILON):
    """Verifie si un nombre est approximativement zero."""
    return abs(val) < epsilon

def create_matrix(rows, cols, value=0.0):
    """Cree une matrice de taille rows x cols initialisee avec 'value'."""
    return [[value for _ in range(cols)] for _ in range(rows)]

def get_matrix_dimensions(matrix):
    """Retourne (rows, cols) d'une matrice."""
    rows = len(matrix)
    if rows == 0:
        return 0, 0
    cols = len(matrix[0])
    return rows, cols

def matrix_scalar_multiply(matrix, scalar):
    """Multiplie une matrice par un scalaire."""
    rows, cols = get_matrix_dimensions(matrix)
    result = create_matrix(rows, cols)
    for r in range(rows):
        for c in range(cols):
            result[r][c] = matrix[r][c] * scalar
    return result

def matrix_vector_multiply(matrix, vector):
    """
    Multiplie une matrice par un vecteur.
    matrix: M x N
    vector: N x 1
    Retourne: M x 1
    """
    rows, cols = get_matrix_dimensions(matrix)
    if cols != len(vector):
        raise ValueError("Le nombre de colonnes de la matrice doit etre egal a la longueur du vecteur.")

    result = [0.0] * rows
    for r in range(rows):
        current_sum = 0.0
        for c in range(cols):
            current_sum += matrix[r][c] * vector[c]
        result[r] = current_sum
    return result

def transpose_matrix(matrix):
    """Calcule la transposee d'une matrice."""
    rows, cols = get_matrix_dimensions(matrix)
    result = create_matrix(cols, rows)
    for r in range(rows):
        for c in range(cols):
            result[c][r] = matrix[r][c]
    return result

def _swap_rows(matrix, r1, r2):
    """Echange deux lignes d'une matrice."""
    matrix[r1], matrix[r2] = matrix[r2], matrix[r1]

def _scale_row(matrix, r, scalar):
    """Multiplie une ligne d'une matrice par un scalaire."""
    for c in range(len(matrix[r])):
        matrix[r][c] *= scalar

def _add_multiple_of_row(matrix, target_r, source_r, scalar):
    """Ajoute un multiple d'une ligne a une autre ligne."""
    for c in range(len(matrix[target_r])):
        matrix[target_r][c] += matrix[source_r][c] * scalar

def rref(matrix):
    """
    Convertit une matrice en sa forme echelonnee reduite (RREF).
    Implemente l'algorithme d'elimination de Gauss-Jordan.
    """
    mat = [row[:] for row in matrix] # Cree une copie pour ne pas modifier l'originale
    rows, cols = get_matrix_dimensions(mat)
    if rows == 0 or cols == 0:
        return mat

    lead = 0 # Colonne pivot actuelle
    for r in range(rows):
        if lead >= cols:
            break

        # Trouver la ligne avec le plus grand element absolu dans la colonne 'lead'
        # pour la stabilite numerique.
        i = r
        while i < rows and is_zero(mat[i][lead]):
            i += 1

        if i == rows: # Pas de pivot dans cette colonne, passer a la suivante
            lead += 1
            continue

        _swap_rows(mat, r, i)

        # Mettre l'element pivot a 1
        lv = mat[r][lead]
        _scale_row(mat, r, 1.0 / lv)

        # Eliminer les autres elements de la colonne 'lead'
        for i_row in range(rows):
            if i_row != r:
                lv = mat[i_row][lead]
                _add_multiple_of_row(mat, i_row, r, -lv)

        lead += 1

    return mat

def get_pivot_columns(rref_matrix):
    """
    Identifie les indices des colonnes pivots dans une matrice RREF.
    Une colonne est pivot si elle contient un 1 qui est le premier element non nul de sa ligne.
    """
    rows, cols = get_matrix_dimensions(rref_matrix)
    pivot_cols = []
    for r in range(rows):
        for c in range(cols):
            if is_zero(rref_matrix[r][c]):
                continue
            if are_approximately_equal(rref_matrix[r][c], 1.0):
                # Verifier si c'est le premier non-zero de la ligne
                is_pivot = True
                for prev_c in range(c):
                    if not is_zero(rref_matrix[r][prev_c]):
                        is_pivot = False
                        break
                if is_pivot:
                    pivot_cols.append(c)
                    break # Passer a la ligne suivante
    return pivot_cols

def find_image_basis(matrix):
    """
    Trouve une base pour l'image (espace colonne) de la matrice.
    La base est formee par les colonnes de la MATRICE ORIGINALE
    correspondant aux colonnes pivots de sa RREF.
    """
    rows, cols = get_matrix_dimensions(matrix)
    if rows == 0 or cols == 0:
        return []

    rref_mat = rref(matrix)
    pivot_col_indices = get_pivot_columns(rref_mat)

    image_basis = []
    for col_idx in pivot_col_indices:
        column_vector = [matrix[r][col_idx] for r in range(rows)]
        image_basis.append(column_vector)

    return image_basis

def find_kernel_basis(matrix):
    """
    Trouve une base pour le noyau (null space) de la matrice.
    Ceci est fait en resolvant Ax = 0 en utilisant la RREF de A.
    """
    rows, cols = get_matrix_dimensions(matrix)
    if rows == 0 or cols == 0:
        return []

    rref_mat = rref(matrix)
    pivot_col_indices = get_pivot_columns(rref_mat)

    free_var_indices = [c for c in range(cols) if c not in pivot_col_indices]
    kernel_basis = []

    for free_idx in free_var_indices:
        # Creer un vecteur solution ou la variable libre courante est 1
        # et les autres variables libres sont 0.
        solution_vec = [0.0] * cols
        solution_vec[free_idx] = 1.0

        # Resoudre pour les variables pivots en utilisant les equations de la RREF
        for r_idx, p_idx in enumerate(pivot_col_indices):
            # p_idx est l'index de la colonne pivot dans la matrice originale
            # r_idx est l'index de la ligne correspondante dans la RREF (si la matrice n'est pas nulle)
            # Les lignes non nulles de la RREF correspondent aux lignes de la matrice originale
            # C'est-a-dire, si rref_mat[r_idx][p_idx] est un pivot 1, alors cette ligne
            # definit x_p_idx en fonction des variables libres.
            # L'equation est x_p_idx + sum(coeff * x_free) = 0
            # donc x_p_idx = - sum(coeff * x_free)

            # Il faut s'assurer que r_idx reste dans les limites des lignes de rref_mat
            if r_idx < rows:
                val = 0.0
                for f_idx in free_var_indices:
                    val += rref_mat[r_idx][f_idx] * solution_vec[f_idx]
                solution_vec[p_idx] = -val
        kernel_basis.append(solution_vec)

    return kernel_basis

def rank(matrix):
    """Calcule le rang d'une matrice (dimension de l'image)."""
    image_basis = find_image_basis(matrix)
    return len(image_basis)

def nullity(matrix):
    """Calcule la nullite d'une matrice (dimension du noyau)."""
    kernel_basis = find_kernel_basis(matrix)
    return len(kernel_basis)

def verify_rank_nullity_theorem(matrix):
    """
    Verifie le theoreme du rang: rank(A) + nullity(A) = nombre de colonnes de A.
    """
    rows, cols = get_matrix_dimensions(matrix)
    if cols == 0: # Cas de base pour les matrices sans colonnes
        return True

    matrix_rank = rank(matrix)
    matrix_nullity = nullity(matrix)

    return matrix_rank + matrix_nullity == cols

def print_matrix(matrix, name="Matrix"):
    """Affiche une matrice."""
    print(f"{name}:")
    for row in matrix:
        print([f"{x:.4f}" for x in row])
    print()

def print_vector(vector, name="Vector"):
    """Affiche un vecteur."""
    print(f"{name}:")
    print([f"{x:.4f}" for x in vector])
    print()

# --- Tests de validation mathematique ---

print("--- Demarrage des tests de validation ---")

# Matrice simple 2x2
A = [[1.0, 2.0],
     [3.0, 4.0]]
v = [1.0, 0.0]

# Test 1: Multiplication Matrice-Vecteur
Av = matrix_vector_multiply(A, v)
assert are_approximately_equal(Av[0], 1.0) and are_approximately_equal(Av[1], 3.0), \
    f"Test 1 failed: matrix_vector_multiply. Expected [1.0, 3.0], got {Av}"
print("Test 1 (Matrice-Vecteur) passe.")

# Test 2: RREF - Matrice Identite
I = [[1.0, 0.0], [0.0, 1.0]]
rref_I = rref(I)
assert all(are_approximately_equal(rref_I[r][c], I[r][c]) for r in range(2) for c in range(2)), \
    f"Test 2 failed: RREF Identity. Expected {I}, got {rref_I}"
print("Test 2 (RREF Identite) passe.")

# Test 3: RREF - Matrice avec elimination
M1 = [[1.0, 2.0, -1.0],
      [2.0, 4.0, -2.0],
      [-1.0, -2.0, 1.0]]
rref_M1_expected = [[1.0, 2.0, -1.0],
                    [0.0, 0.0, 0.0],
                    [0.0, 0.0, 0.0]]
rref_M1 = rref(M1)
assert all(are_approximately_equal(rref_M1[r][c], rref_M1_expected[r][c]) for r in range(3) for c in range(3)), \
    f"Test 3 failed: RREF M1. Expected {rref_M1_expected}, got {rref_M1}"
print("Test 3 (RREF avec elimination) passe.")

# Test 4: RREF - Matrice plus complexe
M2 = [[1.0, 2.0, 3.0, 4.0],
      [4.0, 5.0, 6.0, 7.0],
      [7.0, 8.0, 9.0, 1.0]]
rref_M2_expected = [[1.0, 0.0, -1.0, 0.0],
                    [0.0, 1.0, 2.0, 0.0],
                    [0.0, 0.0, 0.0, 1.0]] # Calcule manuellement pour verifier
rref_M2 = rref(M2)
# Comparaison avec une tolerance pour les flottants
for r in range(3):
    for c in range(4):
        assert are_approximately_equal(rref_M2[r][c], rref_M2_expected[r][c]), \
            f"Test 4 failed at ({r},{c}): RREF M2. Expected {rref_M2_expected[r][c]}, got {rref_M2[r][c]}"
print("Test 4 (RREF plus complexe) passe.")


# Test 5: Noyau - Matrice nulle
Z = [[0.0, 0.0], [0.0, 0.0]]
kernel_Z = find_kernel_basis(Z)
# Pour une matrice 0, le noyau est l'espace entier (R^2), base canonique [[1,0], [0,1]]
assert len(kernel_Z) == 2 and \
       all(are_approximately_equal(kernel_Z[0][c], [1.0, 0.0][c]) for c in range(2)) and \
       all(are_approximately_equal(kernel_Z[1][c], [0.0, 1.0][c]) for c in range(2)), \
    f"Test 5 failed: Kernel of zero matrix. Expected [[1.0, 0.0], [0.0, 1.0]], got {kernel_Z}"
print("Test 5 (Noyau de matrice nulle) passe.")


# Test 6: Noyau - Matrice avec un noyau non reduit au vecteur nul (M1)
# M1 RREF est [[1, 2, -1], [0, 0, 0], [0, 0, 0]]
# Pivot en colonne 0 (x), libres en colonnes 1 (y), 2 (z)
# x + 2y - z = 0 => x = -2y + z
# Base du noyau:
# y=1, z=0 => x=-2. Vector: [-2, 1, 0]
# y=0, z=1 => x=1. Vector: [1, 0, 1]
kernel_M1 = find_kernel_basis(M1)
assert len(kernel_M1) == 2, f"Test 6 failed: Kernel M1 size. Expected 2, got {len(kernel_M1)}"
# Verifier que les vecteurs du noyau donnent 0 quand multiplies par M1
for kv in kernel_M1:
    res = matrix_vector_multiply(M1, kv)
    assert all(is_zero(x) for x in res), f"Test 6 failed: Kernel M1 vector {kv} not in kernel. Result: {res}"
print("Test 6 (Noyau de matrice M1) passe.")


# Test 7: Image - Matrice M1
# M1 RREF est [[1, 2, -1], [0, 0, 0], [0, 0, 0]]
# Colonne pivot est la 0eme. Donc la 0eme colonne de M1 est la base de l'image.
# Base attendue: [[1.0, 2.0, -1.0]]
image_M1 = find_image_basis(M1)
assert len(image_M1) == 1 and \
       all(are_approximately_equal(image_M1[0][r], M1[r][0]) for r in range(3)), \
    f"Test 7 failed: Image M1. Expected [[1.0, 2.0, -1.0]], got {image_M1}"
print("Test 7 (Image de matrice M1) passe.")

# Test 8: Image - Matrice pleine de rang
# M2 RREF: [[1.0, 0.0, -1.0, 0.0], [0.0, 1.0, 2.0, 0.0], [0.0, 0.0, 0.0, 1.0]]
# Colonnes pivots: 0, 1, 3.
# Base de l'image attendue: Colonnes 0, 1, 3 de M2
image_M2_expected_basis = [[1.0, 4.0, 7.0],
                           [2.0, 5.0, 8.0],
                           [4.0, 7.0, 1.0]]
image_M2 = find_image_basis(M2)
assert len(image_M2) == 3, f"Test 8 failed: Image M2 size. Expected 3, got {len(image_M2)}"
for i in range(3):
    for j in range(3):
        assert are_approximately_equal(image_M2[i][j], image_M2_expected_basis[i][j]), \
            f"Test 8 failed: Image M2. Expected {image_M2_expected_basis[i][j]}, got {image_M2[i][j]}"
print("Test 8 (Image de matrice M2) passe.")

# Test 9: Rang
assert rank(M1) == 1, f"Test 9 failed: Rank M1. Expected 1, got {rank(M1)}"
assert rank(M2) == 3, f"Test 9 failed: Rank M2. Expected 3, got {rank(M2)}"
print("Test 9 (Calcul du rang) passe.")

# Test 10: Nullite
assert nullity(M1) == 2, f"Test 10 failed: Nullity M1. Expected 2, got {nullity(M1)}"
assert nullity(M2) == 1, f"Test 10 failed: Nullity M2. Expected 1, got {nullity(M2)}" # 4 cols - 3 rank = 1
print("Test 10 (Calcul de la nullite) passe.")

# Test 11: Theoreme du Rang
assert verify_rank_nullity_theorem(M1), f"Test 11 failed: Rank-Nullity M1. {rank(M1)} + {nullity(M1)} != {get_matrix_dimensions(M1)[1]}"
assert verify_rank_nullity_theorem(M2), f"Test 11 failed: Rank-Nullity M2. {rank(M2)} + {nullity(M2)} != {get_matrix_dimensions(M2)[1]}"
print("Test 11 (Theoreme du Rang) passe.")

# Test 12: Matrice non carree (plus de lignes que de colonnes)
M3 = [[1.0, 0.0],
      [0.0, 1.0],
      [0.0, 0.0]]
# RREF is M3 itself
# Pivot cols: 0, 1
# Image basis: [[1,0,0], [0,1,0]] (cols 0 and 1 of M3)
# Kernel basis: [] (no free variables)
assert rank(M3) == 2, f"Test 12 failed: Rank M3. Expected 2, got {rank(M3)}"
assert nullity(M3) == 0, f"Test 12 failed: Nullity M3. Expected 0, got {nullity(M3)}"
assert verify_rank_nullity_theorem(M3), f"Test 12 failed: Rank-Nullity M3. {rank(M3)} + {nullity(M3)} != {get_matrix_dimensions(M3)[1]}"
print("Test 12 (Matrice non carree, plus de lignes) passe.")

# Test 13: Matrice non carree (plus de colonnes que de lignes)
M4 = [[1.0, 2.0, 3.0],
      [4.0, 5.0, 6.0]]
# RREF_M4: [[1, 0, -1], [0, 1, 2]] (x + 2y + 3z -> x + 0y -z = 0, 0x + y + 2z = 0)
# Pivot cols: 0, 1
# Free var: 2 (z)
# Image basis: [[1,4], [2,5]] (cols 0 and 1 of M4)
# Kernel basis: z=1 -> x=1, y=-2. Vector: [1, -2, 1]
assert rank(M4) == 2, f"Test 13 failed: Rank M4. Expected 2, got {rank(M4)}"
assert nullity(M4) == 1, f"Test 13 failed: Nullity M4. Expected 1, got {nullity(M4)}"
assert verify_rank_nullity_theorem(M4), f"Test 13 failed: Rank-Nullity M4. {rank(M4)} + {nullity(M4)} != {get_matrix_dimensions(M4)[1]}"
kernel_M4 = find_kernel_basis(M4)
assert len(kernel_M4) == 1 and all(is_zero(x) for x in matrix_vector_multiply(M4, kernel_M4[0])), \
    f"Test 13 failed: Kernel M4 vector {kernel_M4[0]} not in kernel."
print("Test 13 (Matrice non carree, plus de colonnes) passe.")

print("\n--- Tous les tests de validation passes avec succes ! ---")


# --- Exemple d'application IA : Transformation de Caracteristiques ---

print("\n--- Exemple d'application IA : Transformation de Caracteristiques ---")

# Supposons que nous avons des donnees de clients avec 3 caracteristiques:
# [Revenu annuel, Age, Score de fidelite]
client_data = [100000.0, 35.0, 8.0]
print_vector(client_data, "Donnees client originales (Revenu, Age, Fidelite)")

# Application lineaire 1: Normalisation et combinaison de caracteristiques
# Nous voulons creer de nouvelles caracteristiques:
# f1 = Revenu / 10000 (normalisation)
# f2 = Age - (Revenu / 5000) (age ajuste par revenu, une forme de maturite financiere)
# f3 = Score de fidelite * Age (engagement pondere par l'age)

# Matrice de transformation T1 (3x3)
T1 = [[1/10000, 0.0,       0.0],
      [-1/5000, 1.0,       0.0],
      [0.0,     8.0,       1.0]] # Note: le 8.0 sur la 2eme ligne de la derniere colonne est une erreur, devrait etre 0.0.
                                # Pour cet exemple, le calcul serait plutot un `age * score`
                                # On va simplifier pour montrer l'idee:
                                # f3 = 0.1 * age + 0.9 * fidelite

# Correction de T1 pour un exemple plus clair
T1 = [[1/10000, 0.0,     0.0],  # f1 = Revenu / 10000
      [-1/5000, 1.0,     0.0],  # f2 = Age - Revenu/5000
      [0.0,     0.1,     0.9]]  # f3 = 0.1*Age + 0.9*Fidelite

print_matrix(T1, "Matrice de Transformation T1 (Normalisation/Combinaison)")
transformed_data_t1 = matrix_vector_multiply(T1, client_data)
print_vector(transformed_data_t1, "Caracteristiques transformees par T1")
# Attendu:
# f1 = 100000 / 10000 = 10.0
# f2 = 35 - (100000 / 5000) = 35 - 20 = 15.0
# f3 = 0.1 * 35 + 0.9 * 8 = 3.5 + 7.2 = 10.7
assert are_approximately_equal(transformed_data_t1[0], 10.0)
assert are_approximately_equal(transformed_data_t1[1], 15.0)
assert are_approximately_equal(transformed_data_t1[2], 10.7)
print("Transformation T1 verifiee.")


# Application lineaire 2: Reduction de dimensionnalite (projection)
# Nous voulons projeter nos donnees dans un espace de 2 dimensions,
# ou la premiere dimension represente le "potentiel financier"
# et la seconde le "niveau d'engagement".
# potentiel financier = 0.7 * f1 + 0.3 * f2 (ponderation des nouvelles caracteristiques)
# engagement = f3 (la meme caracteristique)

# Matrice de transformation T2 (2x3)
T2 = [[0.7, 0.3, 0.0],
      [0.0, 0.0, 1.0]]

print_matrix(T2, "Matrice de Transformation T2 (Reduction de Dimensionnalite)")
reduced_data = matrix_vector_multiply(T2, transformed_data_t1)
print_vector(reduced_data, "Caracteristiques reduites par T2 (Potentiel Financier, Engagement)")
# Attendu:
# potentiel financier = 0.7 * 10.0 + 0.3 * 15.0 = 7.0 + 4.5 = 11.5
# engagement = 10.7
assert are_approximately_equal(reduced_data[0], 11.5)
assert are_approximately_equal(reduced_data[1], 10.7)
print("Reduction de dimensionnalite T2 verifiee.")

print("\n--- Analyse du Noyau et de l'Image pour T2 (Matrice de Reduction) ---")
print_matrix(T2, "Matrice T2")
print(f"Rang de T2 (dimension de l'image, espace des caracteristiques atteignables) : {rank(T2)}")
print(f"Nullite de T2 (dimension du noyau, ce qui est 'ignore') : {nullity(T2)}")

# Theoreme du Rang pour T2: Rang (2) + Nullite (1) = Nombre de colonnes (3)
assert verify_rank_nullity_theorem(T2)
print(f"Verification du theoreme du rang pour T2 : {verify_rank_nullity_theorem(T2)}")

kernel_T2_basis = find_kernel_basis(T2)
print("Base du Noyau de T2 (vecteurs ignores par la transformation) :")
for i, v in enumerate(kernel_T2_basis):
    print_vector(v, f"  Vecteur {i+1}")

# Pour T2 = [[0.7, 0.3, 0.0], [0.0, 0.0, 1.0]]
# RREF est deja cette matrice.
# Colonnes pivots: 0, 2. Variable libre: 1 (le deuxieme element du vecteur d'entree)
# Les equations:
# 0.7 * x0 + 0.3 * x1 = 0
# x2 = 0
# Si x1 = 1 (variable libre)
# 0.7 * x0 + 0.3 * 1 = 0 => x0 = -0.3 / 0.7 = -3/7
# x2 = 0
# Vecteur du noyau : [-3/7, 1, 0] approx [-0.4286, 1.0, 0.0]
expected_kernel_T2 = [-3/7, 1.0, 0.0]
assert len(kernel_T2_basis) == 1
assert all(are_approximately_equal(kernel_T2_basis[0][i], expected_kernel_T2[i]) for i in range(3))
print("Base du Noyau de T2 verifiee.")

print("\nInterpretation pour T2:")
print(f"- Le rang de {rank(T2)} signifie que les donnees transformees resident dans un espace de {rank(T2)} dimensions, malgre l'entree en {get_matrix_dimensions(T2)[1]} dimensions.")
print(f"- La nullite de {nullity(T2)} signifie qu'il y a {nullity(T2)} dimension(s) dans l'espace des caracteristiques d'entree qui ne contribuent pas a la sortie (sont 'filtrees' ou mappees a zero).")
print(f"- Le vecteur de base du noyau [-0.4286, 1.0, 0.0] indique que si les caracteristiques transformees par T1 varient proportionnellement a ce vecteur (c'est-a-dire si le second terme (f2) augmente de 1 et le premier (f1) diminue de 0.4286, et le troisieme (f3) ne change pas), cela n'affectera pas la projection sur l'axe du 'potentiel financier' dans l'espace image de T2. Plus precisement, cela affecterait *uniquement* la premiere equation. L'effet de la variation de f1 et f2 est annule pour l'axe 'potentiel financier'. La valeur de 'engagement' (f3) est independante.")

image_T2_basis = find_image_basis(T2)
print("Base de l'Image de T2 (espace des caracteristiques atteignables) :")
for i, v in enumerate(image_T2_basis):
    print_vector(v, f"  Vecteur {i+1}")
# Les vecteurs de la base de l'image sont les colonnes pivots de T2:
# Colonne 0: [0.7, 0.0]
# Colonne 2: [0.0, 1.0]
expected_image_T2_basis = [[0.7, 0.0], [0.0, 1.0]]
assert len(image_T2_basis) == 2
for i in range(2):
    for j in range(2):
        assert are_approximately_equal(image_T2_basis[i][j], expected_image_T2_basis[i][j])
print("Base de l'Image de T2 verifiee.")

print("\nInterpretation des bases de l'Image:")
print("- La base de l'image nous donne les 'directions' ou 'composantes principales' qui definissent l'espace de sortie. Dans ce cas, elles correspondent aux deux dimensions que nous avons definies : 'potentiel financier' et 'engagement'.")
\end{verbatim}

\newpage
\chapter*{TP 5 : Géométrie des Transformations Linéaires : Noyau, Image et Rang "from scratch"}

\section*{Objectif Pédagogique (Mathématique et IA)}

Ce TP a pour objectif de solidifier la compréhension et l'implémentation des concepts fondamentaux d'algèbre linéaire que sont les applications linéaires, leur noyau (kernel), leur image, et le théorème du rang. En codant ces concepts "from scratch" en Python, sans l'aide de bibliothèques de haut niveau comme NumPy, nous allons approfondir la mécanique sous-jacente des transformations de données.

\textbf{Objectifs Mathématiques :}
*   Implémenter la représentation matricielle d'une application linéaire.
*   Maîtriser l'algorithme d'élimination de Gauss-Jordan pour transformer une matrice en forme échelonnée réduite (RREF).
*   Déduire le noyau (espace nul) d'une matrice à partir de sa RREF.
*   Déduire l'image (espace colonne) d'une matrice à partir de sa forme échelonnée.
*   Calculer le rang et la nullité d'une matrice.
*   Vérifier le théorème du rang : \_\_INLINE\_CODE\_0\_\_.

\textbf{Objectifs IA et Machine Learning :}
*   \textbf{Compréhension des Transformations de Données :} Les applications linéaires sont au cœur des modèles d'IA, notamment les réseaux de neurones (où les couches sont des transformations linéaires). Comprendre le noyau et l'image permet de saisir comment les informations sont transformées, perdues ou préservées.
*   \textbf{Réduction de Dimensionalité :} Le concept d'image est directement lié à des techniques comme l'Analyse en Composantes Principales (PCA), où l'on projette des données sur un sous-espace d'intérêt (l'image d'une matrice de projection) pour réduire la complexité tout en maximisant la variance.
*   \textbf{Redondance et Dépendance Linéaire :} Le noyau révèle les combinaisons de caractéristiques d'entrée qui n'apportent aucune information utile après une transformation, car elles sont mappées au vecteur nul. Cela aide à identifier la redondance dans les données ou les caractéristiques.
*   \textbf{Capacité des Modèles :} Le rang d'une matrice de poids dans un réseau de neurones peut donner une idée de la capacité expressive du modèle ou de la complexité des transformations qu'il peut effectuer. Une matrice de faible rang peut indiquer un modèle sous-paramétré ou une transformation restrictive.
*   \textbf{Stabilité Numérique et Inversibilité :} Bien que non directement implémentée, la notion de noyau est cruciale pour comprendre l'inversibilité des matrices, un concept fondamental en optimisation et dans la résolution de systèmes linéaires.

Ce TP met l'accent sur l'importance de construire une intuition géométrique et algébrique solide, qui est essentielle pour la conception, l'analyse et le débogage de systèmes d'IA complexes.

\section*{Code Source Python "from scratch"}

\begin{verbatim}
import math

# --- Configuration pour la precision des flottants ---
EPSILON = 1e-9 # Petite valeur pour comparer les flottants a zero

def is_zero(value):
    """Verifie si une valeur flottante est proche de zero."""
    return abs(value) < EPSILON

def format_vector(v):
    """Formate un vecteur pour un affichage lisible."""
    return "[" + ", ".join(f"{x:.4f}" for x in v) + "]"

def format_matrix(m):
    """Formate une matrice pour un affichage lisible."""
    if not m: return "[]"
    return "[\n" + ",\n".join("  " + format_vector(row) for row in m) + "\n]"

# --- Operations de base sur les vecteurs et matrices ---

def add_vectors(v1, v2):
    """Additionne deux vecteurs."""
    if len(v1) != len(v2):
        raise ValueError("Les vecteurs doivent avoir la meme dimension pour etre additionnes.")
    return [v1[i] + v2[i] for i in range(len(v1))]

def scalar_mult_vector(scalar, v):
    """Multiplie un vecteur par un scalaire."""
    return [scalar * x for x in v]

def dot_product(v1, v2):
    """Calcule le produit scalaire de deux vecteurs."""
    if len(v1) != len(v2):
        raise ValueError("Les vecteurs doivent avoir la meme dimension pour le produit scalaire.")
    return sum(v1[i] * v2[i] for i in range(len(v1)))

def matrix_vector_mult(matrix, vector):
    """Multiplie une matrice par un vecteur."""
    num_rows = len(matrix)
    num_cols = len(matrix[0]) if num_rows > 0 else 0

    if num_cols != len(vector):
        raise ValueError("Le nombre de colonnes de la matrice doit etre egal a la dimension du vecteur.")

    result = [0.0] * num_rows
    for i in range(num_rows):
        result[i] = dot_product(matrix[i], vector)
    return result

def transpose_matrix(matrix):
    """Calcule la transposee d'une matrice."""
    if not matrix or not matrix[0]:
        return []
    num_rows = len(matrix)
    num_cols = len(matrix[0])
    transposed = [[0.0 for _ in range(num_rows)] for _ in range(num_cols)]
    for i in range(num_rows):
        for j in range(num_cols):
            transposed[j][i] = matrix[i][j]
    return transposed

def matrix_mult(m1, m2):
    """Multiplie deux matrices."""
    if not m1 or not m1[0] or not m2 or not m2[0]:
        raise ValueError("Les matrices ne peuvent pas etre vides.")

    rows_m1 = len(m1)
    cols_m1 = len(m1[0])
    rows_m2 = len(m2)
    cols_m2 = len(m2[0])

    if cols_m1 != rows_m2:
        raise ValueError("Le nombre de colonnes de la premiere matrice doit etre egal au nombre de lignes de la seconde.")

    m2_t = transpose_matrix(m2) # Transposer m2 pour un acces plus direct aux colonnes
    result = [[0.0 for _ in range(cols_m2)] for _ in range(rows_m1)]

    for i in range(rows_m1):
        for j in range(cols_m2):
            result[i][j] = dot_product(m1[i], m2_t[j])
    return result

# --- Elimination de Gauss-Jordan (RREF) ---

def find_pivot_row(matrix, col, start_row):
    """Trouve la premiere ligne non nulle (pivot) dans une colonne a partir de start_row."""
    num_rows = len(matrix)
    for i in range(start_row, num_rows):
        if not is_zero(matrix[i][col]):
            return i
    return -1

def swap_rows(matrix, r1, r2):
    """Echange deux lignes dans une matrice."""
    matrix[r1], matrix[r2] = matrix[r2], matrix[r1]

def scale_row(matrix, row, scalar):
    """Multiplie une ligne par un scalaire."""
    for j in range(len(matrix[row])):
        matrix[row][j] *= scalar

def add_scaled_row(matrix, target_row, source_row, scalar):
    """Ajoute un multiple scalaire d'une ligne a une autre ligne."""
    for j in range(len(matrix[target_row])):
        matrix[target_row][j] += scalar * matrix[source_row][j]

def row_echelon_form(matrix_original):
    """
    Convertit une matrice en Forme Echelonnee (REF).
    Retourne une nouvelle matrice pour ne pas modifier l'originale.
    """
    matrix = [row[:] for row in matrix_original] # Copie la matrice pour ne pas modifier l'originale
    num_rows = len(matrix)
    if num_rows == 0: return []
    num_cols = len(matrix[0])
    if num_cols == 0: return [[] for _ in range(num_rows)]

    lead = 0 # Indice de colonne du pivot
    for r in range(num_rows):
        if lead >= num_cols:
            break
        i = find_pivot_row(matrix, lead, r)
        if i != -1: # Si un pivot est trouve
            swap_rows(matrix, r, i)
            # Normaliser la ligne pour que le pivot soit 1
            lv = matrix[r][lead]
            scale_row(matrix, r, 1.0 / lv)
            # Eliminer les elements en dessous du pivot
            for i in range(num_rows):
                if i != r:
                    lv_below = matrix[i][lead]
                    add_scaled_row(matrix, i, r, -lv_below)
            lead += 1
        else: # Pas de pivot dans cette colonne, passer a la suivante
            lead += 1

    # Nettoyage des petites valeurs pour eviter des erreurs de flottants
    for r in range(num_rows):
        for c in range(num_cols):
            if is_zero(matrix[r][c]):
                matrix[r][c] = 0.0 # Mettre a zero les valeurs tres proches de zero

    return matrix


def reduced_row_echelon_form(matrix_original):
    """
    Convertit une matrice en Forme Echelonnee Reduite (RREF).
    Retourne une nouvelle matrice pour ne pas modifier l'originale.
    """
    matrix = [row[:] for row in matrix_original] # Copie la matrice pour ne pas modifier l'originale
    num_rows = len(matrix)
    if num_rows == 0: return []
    num_cols = len(matrix[0])
    if num_cols == 0: return [[] for _ in range(num_rows)]

    lead = 0
    for r in range(num_rows):
        if lead >= num_cols:
            break
        i = find_pivot_row(matrix, lead, r)
        if i != -1:
            swap_rows(matrix, r, i)
            # Normaliser la ligne pour que le pivot soit 1
            lv = matrix[r][lead]
            scale_row(matrix, r, 1.0 / lv)
            # Eliminer les elements au-dessus et en dessous du pivot
            for i in range(num_rows):
                if i != r:
                    lv_other = matrix[i][lead]
                    add_scaled_row(matrix, i, r, -lv_other)
            lead += 1
        else: # Pas de pivot dans cette colonne, passer a la suivante
            lead += 1

    # Nettoyage des petites valeurs pour eviter des erreurs de flottants
    for r in range(num_rows):
        for c in range(num_cols):
            if is_zero(matrix[r][c]):
                matrix[r][c] = 0.0

    return matrix

# --- Fonctions pour Noyau, Image, Rang ---

def get_rank(matrix):
    """
    Calcule le rang d'une matrice (dimension de l'image).
    Le rang est le nombre de lignes non nulles dans la forme echelonnee reduite.
    """
    rref = reduced_row_echelon_form(matrix)
    rank = 0
    for row in rref:
        if any(not is_zero(x) for x in row):
            rank += 1
    return rank

def find_image_basis(matrix):
    """
    Trouve une base pour l'image (espace colonne) d'une matrice.
    Les colonnes de la matrice originale correspondant aux colonnes pivot de la RREF
    forment une base pour l'image.
    """
    num_rows = len(matrix)
    if num_rows == 0: return []
    num_cols = len(matrix[0])
    if num_cols == 0: return []

    rref = reduced_row_echelon_form(matrix)

    pivot_columns_indices = []
    current_pivot_row = 0
    for j in range(num_cols):
        if current_pivot_row < num_rows and not is_zero(rref[current_pivot_row][j]):
            pivot_columns_indices.append(j)
            current_pivot_row += 1

    # Extraire les colonnes originales correspondant aux pivots
    image_basis = []
    for col_idx in pivot_columns_indices:
        column_vector = [matrix[i][col_idx] for i in range(num_rows)]
        image_basis.append(column_vector)

    return image_basis

def find_kernel_basis(matrix):
    """
    Trouve une base pour le noyau (espace nul) d'une matrice.
    Utilise la RREF pour identifier les variables libres et pivots.
    """
    num_rows = len(matrix)
    if num_rows == 0: return []
    num_cols = len(matrix[0])
    if num_cols == 0: return [[]] # Vecteur nul de la bonne dimension si l'espace est vide

    rref = reduced_row_echelon_form(matrix)

    pivot_cols = []
    free_cols = []
    pivot_row_for_col = [-1] * num_cols # Associe chaque colonne pivot a sa ligne

    current_pivot_row = 0
    for j in range(num_cols):
        if current_pivot_row < num_rows and not is_zero(rref[current_pivot_row][j]):
            pivot_cols.append(j)
            pivot_row_for_col[j] = current_pivot_row
            current_pivot_row += 1
        else:
            free_cols.append(j)

    kernel_basis = []

    for free_col_idx in free_cols:
        basis_vector = [0.0] * num_cols
        basis_vector[free_col_idx] = 1.0 # La variable libre correspondante est 1

        # Pour chaque variable pivot, exprimer en fonction de cette variable libre
        for pivot_col_idx in pivot_cols:
            r = pivot_row_for_col[pivot_col_idx]
            # Si le terme correspondant a la variable libre est non nul dans l'equation du pivot
            basis_vector[pivot_col_idx] = -rref[r][free_col_idx]

        kernel_basis.append(basis_vector)

    # Si toutes les colonnes sont pivots et que le rang est plein (matrice carree inversible), le noyau est le vecteur nul.
    # Dans ce cas, free_cols sera vide, et kernel_basis sera vide.
    # On ajoute le vecteur nul [0,0,...,0] si l'espace du noyau est juste {0}
    if not free_cols and get_rank(matrix) == num_cols:
         return [] # Noyau reduit au vecteur nul, seul le vecteur nul est dedans. Pas de base non immediate d'apres la definition.

    return kernel_basis

def get_nullity(matrix):
    """
    Calcule la nullite d'une matrice (dimension du noyau).
    """
    return len(find_kernel_basis(matrix))

def check_rank_nullity_theorem(matrix):
    """
    Verifie le theoreme du rang: dim(Ker(A)) + dim(Im(A)) = dim(domaine).
    dim(domaine) est le nombre de colonnes de la matrice.
    """
    if not matrix or not matrix[0]:
        # Cas des matrices vides ou a une seule ligne/colonne vide
        num_cols = len(matrix[0]) if matrix else 0
        return get_nullity(matrix) + get_rank(matrix) == num_cols

    num_cols = len(matrix[0])
    rank = get_rank(matrix)
    nullity = get_nullity(matrix)

    return rank + nullity == num_cols

# --- Assertions de Validation Mathematique ---

print("--- Debut des tests ---")

# Matrice identite 3x3
I3 = [[1.0, 0.0, 0.0],
      [0.0, 1.0, 0.0],
      [0.0, 0.0, 1.0]]

# Matrice nulle 2x3
Z23 = [[0.0, 0.0, 0.0],
       [0.0, 0.0, 0.0]]

# Matrice A (full rank 2x2)
A = [[1.0, 2.0],
     [3.0, 4.0]]

# Matrice B (non-carree)
B = [[1.0, 2.0, 3.0],
     [4.0, 5.0, 6.0]]

# Matrice C (avec dependance lineaire)
C = [[1.0, 2.0, 3.0],
     [2.0, 4.0, 6.0],
     [7.0, 8.0, 9.0]]

# Matrice D (avec dependance lineaire plus complexe)
D = [[1.0, 1.0, 2.0],
     [2.0, 2.0, 4.0],
     [0.0, 1.0, 1.0]]

# Matrice E (pour un noyau non reduit au vecteur nul simple)
E = [[1.0, 1.0, 1.0],
     [0.0, 1.0, 1.0],
     [0.0, 0.0, 0.0]]

# Matrice F (pour un noyau 2D)
F = [[1.0, 2.0, 3.0, 4.0],
     [5.0, 6.0, 7.0, 8.0],
     [9.0, 10.0, 11.0, 12.0]]


print("\n--- Test 1: Operations de base ---")
v1 = [1.0, 2.0]
v2 = [3.0, 4.0]
assert add_vectors(v1, v2) == [4.0, 6.0]
assert scalar_mult_vector(2.0, v1) == [2.0, 4.0]
assert dot_product(v1, v2) == 1.0*3.0 + 2.0*4.0 # 3 + 8 = 11
assert matrix_vector_mult(A, v1) == [1.0*1.0 + 2.0*2.0, 3.0*1.0 + 4.0*2.0] # [5, 11]
assert matrix_vector_mult(A, v1) == [5.0, 11.0]

assert matrix_mult(A, I3[:2]) == A # Test A * I_2x2 (partie de I3)
# Need a 2x2 identity matrix for the above. Let's define one.
I2 = [[1.0, 0.0], [0.0, 1.0]]
assert matrix_mult(A, I2) == A
print("Test 1: Operations de base PASSED")


print("\n--- Test 2: Reduction RREF ---")
# RREF de I3 doit etre I3
assert reduced_row_echelon_form(I3) == I3

# RREF de Z23 doit etre Z23
assert reduced_row_echelon_form(Z23) == Z23

# RREF de A (matrice inversible)
A_rref = reduced_row_echelon_form(A)
expected_A_rref = [[1.0, 0.0], [0.0, 1.0]]
assert all(all(is_zero(A_rref[i][j] - expected_A_rref[i][j]) for j in range(len(A_rref[0]))) for i in range(len(A_rref)))

# RREF de C (dependance lineaire)
C_rref = reduced_row_echelon_form(C)
expected_C_rref = [[1.0, 0.0, -1.0],
                   [0.0, 1.0, 2.0],
                   [0.0, 0.0, 0.0]]
for i in range(len(C_rref)):
    for j in range(len(C_rref[0])):
        assert is_zero(C_rref[i][j] - expected_C_rref[i][j]), f"C_rref[{i}][{j}] = {C_rref[i][j]}, Expected: {expected_C_rref[i][j]}"
print("Test 2: Reduction RREF PASSED")


print("\n--- Test 3: Rang de la matrice ---")
assert get_rank(I3) == 3
assert get_rank(Z23) == 0
assert get_rank(A) == 2 # Matrice 2x2 full rank
assert get_rank(B) == 2 # Matrice 2x3 full row rank
assert get_rank(C) == 2 # Matrice 3x3, rang 2
assert get_rank(D) == 2 # Matrice 3x3, rang 2
assert get_rank(E) == 2 # Matrice 3x3, rang 2
assert get_rank(F) == 2 # Matrice 3x4, rang 2
print("Test 3: Rang de la matrice PASSED")


print("\n--- Test 4: Base de l'Image (espace colonne) ---")
# Image de I3 est R^3, base est les vecteurs canoniques
image_I3_basis = find_image_basis(I3)
assert len(image_I3_basis) == 3
assert all(any(is_zero(x - expected_x) for x, expected_x in zip(v, [1.0, 0.0, 0.0])) for v in image_I3_basis)
assert all(any(is_zero(x - expected_x) for x, expected_x in zip(v, [0.0, 1.0, 0.0])) for v in image_I3_basis)
assert all(any(is_zero(x - expected_x) for x, expected_x in zip(v, [0.0, 0.0, 1.0])) for v in image_I3_basis)

# Image de A (matrice 2x2 full rank) est R^2, base est les colonnes de A
image_A_basis = find_image_basis(A)
assert len(image_A_basis) == 2
assert image_A_basis == [[1.0, 3.0], [2.0, 4.0]] # colonnes originales de A

# Image de C (rang 2)
image_C_basis = find_image_basis(C)
# Les colonnes pivot de C sont la 1ere et la 2eme (indices 0 et 1)
assert len(image_C_basis) == 2
assert image_C_basis[0] == [1.0, 2.0, 7.0]
assert image_C_basis[1] == [2.0, 4.0, 8.0]

# Image de F (rang 2)
image_F_basis = find_image_basis(F)
assert len(image_F_basis) == 2
assert image_F_basis[0] == [1.0, 5.0, 9.0]
assert image_F_basis[1] == [2.0, 6.0, 10.0]
print("Test 4: Base de l'Image (espace colonne) PASSED")


print("\n--- Test 5: Base du Noyau (espace nul) ---")
# Noyau de I3 (matrice inversible) est {0}
kernel_I3_basis = find_kernel_basis(I3)
assert len(kernel_I3_basis) == 0 # Noyau reduit au vecteur nul

# Noyau de A (matrice 2x2 inversible) est {0}
kernel_A_basis = find_kernel_basis(A)
assert len(kernel_A_basis) == 0

# Noyau de C (rang 2, 3 colonnes, nullite 1)
kernel_C_basis = find_kernel_basis(C)
assert len(kernel_C_basis) == 1
# De la RREF de C: x1 - x3 = 0 => x1 = x3, x2 + 2x3 = 0 => x2 = -2x3
# Variable libre x3. Si x3 = 1, alors x1 = 1, x2 = -2. Vecteur: [1, -2, 1]
expected_kernel_C_vec = [1.0, -2.0, 1.0]
# Assurez-vous que le vecteur du noyau trouve est un multiple de l'attendu.
# matrix_vector_mult(C, expected_kernel_C_vec) should be 0 vector
res_C_kernel = matrix_vector_mult(C, kernel_C_basis[0])
assert all(is_zero(x) for x in res_C_kernel)

# Noyau de E (rang 2, 3 colonnes, nullite 1)
kernel_E_basis = find_kernel_basis(E)
assert len(kernel_E_basis) == 1
# RREF de E: [[1, 0, 0], [0, 1, 1], [0, 0, 0]]
# x1 = 0, x2 + x3 = 0 => x2 = -x3. x3 est libre. Si x3=1, x2=-1, x1=0.
# Vecteur: [0, -1, 1]
expected_kernel_E_vec = [0.0, -1.0, 1.0]
res_E_kernel = matrix_vector_mult(E, kernel_E_basis[0])
assert all(is_zero(x) for x in res_E_kernel)

# Noyau de F (rang 2, 4 colonnes, nullite 2)
kernel_F_basis = find_kernel_basis(F)
assert len(kernel_F_basis) == 2
# RREF de F:
# [[1.0, 0.0, -1.0, -2.0],
#  [0.0, 1.0, 2.0, 3.0],
#  [0.0, 0.0, 0.0, 0.0]]
# x1 - x3 - 2x4 = 0  => x1 = x3 + 2x4
# x2 + 2x3 + 3x4 = 0 => x2 = -2x3 - 3x4
# x3, x4 sont libres.
# 1) x3=1, x4=0 => x1=1, x2=-2. Vecteur: [1, -2, 1, 0]
# 2) x3=0, x4=1 => x1=2, x2=-3. Vecteur: [2, -3, 0, 1]
# Verifier que les vecteurs generes sont bien dans le noyau
for v_basis in kernel_F_basis:
    res_F_kernel = matrix_vector_mult(F, v_basis)
    assert all(is_zero(x) for x in res_F_kernel)
print("Test 5: Base du Noyau (espace nul) PASSED")


print("\n--- Test 6: Nullite de la matrice ---")
assert get_nullity(I3) == 0
assert get_nullity(A) == 0
assert get_nullity(C) == 1
assert get_nullity(D) == 1 # rang 2, nullite 1 (pour D = [[1, 1, 2], [2, 2, 4], [0, 1, 1]])
assert get_nullity(E) == 1
assert get_nullity(F) == 2
print("Test 6: Nullite de la matrice PASSED")


print("\n--- Test 7: Theoreme du Rang ---")
# dim(domaine) est le nombre de colonnes
assert check_rank_nullity_theorem(I3) == True # rank=3, nullity=0, cols=3
assert check_rank_nullity_theorem(Z23) == True # rank=0, nullity=3, cols=3 (nullity = len(Z23[0]) - rank )
assert get_nullity(Z23) == 3 # le noyau de Z23 est R^3, de base [[1,0,0],[0,1,0],[0,0,1]]
assert check_rank_nullity_theorem(A) == True # rank=2, nullity=0, cols=2
assert check_rank_nullity_theorem(B) == True # rank=2, nullity=1, cols=3
assert check_rank_nullity_theorem(C) == True # rank=2, nullity=1, cols=3
assert check_rank_nullity_theorem(D) == True # rank=2, nullity=1, cols=3
assert check_rank_nullity_theorem(E) == True # rank=2, nullity=1, cols=3
assert check_rank_nullity_theorem(F) == True # rank=2, nullity=2, cols=4
print("Test 7: Theoreme du Rang PASSED")

print("\n--- Tous les tests sont passes avec succes ! ---")
\end{verbatim}

\newpage
\end{document}